% Lutte contre l'effet de serre — SVT 1ère D — Cours signé OnBuch+
% Programme officiel MINESEC. Compiler avec : tectonic ND11-lutte-effet-serre-d.tex
\newcommand{\DOCMATIERE}{SVT}
\newcommand{\DOCNIVEAU}{1ère D}
\newcommand{\DOCMODULE}{Module 3 : L'Éducation à l'Environnement et au Développement Durable}
\newcommand{\DOCLECON}{ND11}
\newcommand{\DOCTITRE}{Lutte contre l'effet de serre}
% OnBuch+ — préambule « cours signés » ENRICHI (compilé avec tectonic / XeLaTeX)
% Système visuel « Pop » d'OnBuch+ : fond crème (jamais blanc), bordures encre
% épaisses, ombres « dures » décalées, titres Archivo Black en capitales.
% Version MATHÉMATIQUES : graphes (pgfplots/tikz), boîtes pédagogiques
% (définition, propriété, méthode, attention, le savais-tu, exercice résolu,
% exercice, TP, bilan), page de garde et signature OnBuch+ ancrée en bas.
%
% Macros attendues dans main.tex AVANT \input{preamble} :
%   \DOCMATIERE  \DOCNIVEAU  \DOCMODULE  \DOCLECON (numéro)  \DOCTITRE
\documentclass[11pt]{article}

\usepackage{fontspec}
\usepackage[french]{babel}
\usepackage[a4paper,margin=2cm,top=3.4cm,bottom=2.8cm,headheight=1.4cm,footskip=1.3cm]{geometry}
\usepackage{amsmath,amssymb,amsfonts}
\usepackage{mathtools} % \xmapsto, \coloneqq... souvent utilisés par les modèles
\usepackage{cancel} % simplifications barrées (bilans de forces)
% \abs{z} (module, valeur absolue) et \norm{u} (norme), utilisés par les modèles
\providecommand{\abs}{}\let\abs\relax\DeclarePairedDelimiter{\abs}{\lvert}{\rvert}
\providecommand{\norm}{}\let\norm\relax\DeclarePairedDelimiter{\norm}{\lVert}{\rVert}
\usepackage[table]{xcolor} % [table] : fournit aussi \rowcolors (alternance de lignes)
\usepackage{pagecolor}
\usepackage{tikz}
\usetikzlibrary{arrows.meta,positioning,calc,shapes.geometric,shapes.misc,shapes.arrows,decorations.pathmorphing,decorations.pathreplacing,decorations.markings,patterns,shadows,mindmap,trees,fit,backgrounds,matrix,intersections,angles,quotes}
\usepackage{pgfplots}
\usepackage[version=4]{mhchem} % équations nucléaires \ce{^{235}_{92}U}
\usepackage[european,siunitx]{circuitikz} % schémas électriques
\pgfplotsset{compat=1.18}
\usepgfplotslibrary{fillbetween,groupplots} % aires entre courbes (calcul intégral)
\usepackage{enumitem}
\usepackage{fancyhdr}
% [most] charge toutes les bibliothèques tcolorbox (listings, minted, raster,
% documentation...) et gonfle inutilement la mémoire TeX sur un document long
% (~300 boîtes) : on ne charge que ce qui est réellement utilisé.
\usepackage[skins,breakable]{tcolorbox}
\usepackage{graphicx}
\usepackage{colortbl}
\usepackage{booktabs}
\usepackage{tabularx}
\usepackage{diagbox} % case d'en-tête diagonale des tableaux à double entrée
% Colonnes extensibles centrée / gauche / droite (C, L, R), souvent utilisées par les modèles
\newcolumntype{C}{>{\centering\arraybackslash}X}
\newcolumntype{L}{>{\raggedright\arraybackslash}X}
\newcolumntype{R}{>{\raggedleft\arraybackslash}X}
\usepackage{array}
\usepackage{multicol}
\usepackage{siunitx}
% parse-numbers=false : accepte aussi \SI{2,5 \times 10^{-3}}{...}
\sisetup{locale=FR,output-decimal-marker={,},group-minimum-digits=4,parse-numbers=false}
\DeclareSIUnit\ohm{\ensuremath{\symup{\Omega}}} % U+2126 absent de la police
\DeclareSIUnit\division{div} % réglages d'oscilloscope (ms/div, V/div)
\DeclareSIUnit\tr{tr} % tours (vitesse de rotation, ex. \SI{1500}{\tr\per\minute})
\DeclareSIUnit\ch{ch} % cheval-vapeur (puissance moteur, unité usuelle hors SI)
\DeclareSIUnit\franccfa{FCFA} % franc CFA (coûts, contexte camerounais)
\DeclareSIUnit\dioptre{\ensuremath{\delta}} % vergence d'une lentille (dioptrie)
\usepackage{qrcode}
% \degree : commande fréquemment utilisée par les modèles pour un angle ou une
% température en dehors de \SI{}{} (non fournie par les paquets chargés).
\providecommand{\degree}{^{\circ}}
% \qed (fin de démonstration), utilisé par les modèles sans amsthm
\providecommand{\qed}{\hfill$\square$}
% \barre : double barre des valeurs interdites dans un tableau de variations (tkz-tab)
\providecommand{\barre}{\ensuremath{\|}}
% environnement proof (amsthm non chargé) utilisé par les modèles
\ifdefined\proof\else\newenvironment{proof}[1][Démonstration]{\par\noindent\textit{#1.}\ }{\qed\par}\fi
\usepackage{lastpage}
\usepackage{hyperref}
\hypersetup{hidelinks}

% ============================================================
% Palette « Pop » OnBuch+ — mêmes hex que la classe P (plus_ui.dart)
% ============================================================
\definecolor{poporange}{HTML}{F59321}
\definecolor{popdark}{HTML}{E07A0C}
\definecolor{popcream}{HTML}{FFF6E8}
\definecolor{popink}{HTML}{1C1714}
\definecolor{poppurple}{HTML}{7A5AE0}
\definecolor{popgreen}{HTML}{2E9E6A}
\definecolor{poppink}{HTML}{E0457B}
\definecolor{popblue}{HTML}{2F7DE1}
\definecolor{popgold}{HTML}{C98A12}
\definecolor{popmuted}{HTML}{8A7F76}
\definecolor{popline}{HTML}{EADFD2}
\definecolor{popgray}{HTML}{8A7F76}
\colorlet{popgrey}{popgray}
% Alias tolérants (le modèle nomme parfois les couleurs différemment)
\colorlet{popwhite}{white}
\colorlet{popblack}{popink}
\colorlet{popred}{poppink}
\colorlet{popyellow}{popgold}
% Teintes claires pour fonds de tableaux / zones
\colorlet{popblueL}{popblue!10!white}
\colorlet{poporangeL}{poporange!14!white}
\colorlet{popgreenL}{popgreen!12!white}
\colorlet{poppurpleL}{poppurple!12!white}
\colorlet{poppinkL}{poppink!12!white}
\colorlet{popgoldL}{popgold!14!white}

\pagecolor{popcream}

% ============================================================
% Typographie
% ============================================================
\defaultfontfeatures{Ligatures=TeX}
\setmainfont{PlusJakartaSans-Medium.ttf}[Path=../../../fonts/,BoldFont=PlusJakartaSans-Bold.ttf]
% Maths en sans-serif (Fira Math) avec chiffres et lettres droites de Plus
% Jakarta, pour que les formules se fondent dans le texte.
\usepackage[math-style=ISO,bold-style=ISO]{unicode-math}
\setmathfont{FiraMath-Regular.otf}[Path=../../../fonts/]
\setmathfont{PlusJakartaSans-Medium.ttf}[Path=../../../fonts/,range={up/num,up/{Latin,latin},\mathup}]
\usepackage{icomma} % virgule décimale sans espace en mode math
% Caractères mathématiques Unicode absents des polices de texte (Plus Jakarta,
% Archivo Black) : rendus par leur équivalent mathématique, en texte comme en
% formule — sinon ils s'affichent en carré vide (ex. « Arithmétique dans ℤ »).
\usepackage{newunicodechar}
\newunicodechar{ℝ}{\ensuremath{\mathbb{R}}}
\newunicodechar{ℤ}{\ensuremath{\mathbb{Z}}}
\newunicodechar{ℂ}{\ensuremath{\mathbb{C}}}
\newunicodechar{ℕ}{\ensuremath{\mathbb{N}}}
\newunicodechar{ℚ}{\ensuremath{\mathbb{Q}}}
\newunicodechar{𝔻}{\ensuremath{\mathbb{D}}}
\newunicodechar{α}{\ensuremath{\alpha}}
\newunicodechar{β}{\ensuremath{\beta}}
\newunicodechar{γ}{\ensuremath{\gamma}}
\newunicodechar{δ}{\ensuremath{\delta}}
\newunicodechar{ε}{\ensuremath{\varepsilon}}
\newunicodechar{θ}{\ensuremath{\theta}}
\newunicodechar{λ}{\ensuremath{\lambda}}
\newunicodechar{μ}{\ensuremath{\mu}}
\newunicodechar{σ}{\ensuremath{\sigma}}
\newunicodechar{τ}{\ensuremath{\tau}}
\newunicodechar{φ}{\ensuremath{\varphi}}
\newunicodechar{ψ}{\ensuremath{\psi}}
\newunicodechar{ω}{\ensuremath{\omega}}
\newunicodechar{Σ}{\ensuremath{\Sigma}}
% majuscules produites par \MakeUppercase dans les titres (α-aminés -> Α)
\newunicodechar{Α}{\ensuremath{\alpha}}
\newunicodechar{Β}{\ensuremath{\beta}}
\newunicodechar{Γ}{\ensuremath{\Gamma}}
\newunicodechar{Δ}{\ensuremath{\Delta}}
\newunicodechar{Ε}{\ensuremath{\varepsilon}}
\newunicodechar{Θ}{\ensuremath{\Theta}}
\newunicodechar{Λ}{\ensuremath{\Lambda}}
\newunicodechar{Μ}{\ensuremath{\mu}}
\newunicodechar{Τ}{\ensuremath{\tau}}
\newunicodechar{Φ}{\ensuremath{\Phi}}
\newunicodechar{Ψ}{\ensuremath{\Psi}}
\newunicodechar{Ω}{\ensuremath{\Omega}}
\newunicodechar{↦}{\ensuremath{\mapsto}}
\newunicodechar{∈}{\ensuremath{\in}}
\newunicodechar{∉}{\ensuremath{\notin}}
\newunicodechar{∧}{\ensuremath{\wedge}}
\newunicodechar{⇔}{\ensuremath{\Leftrightarrow}}
\newunicodechar{⟺}{\ensuremath{\Longleftrightarrow}}
\newunicodechar{⇒}{\ensuremath{\Rightarrow}}
\newunicodechar{⟹}{\ensuremath{\Longrightarrow}}
\newunicodechar{∘}{\ensuremath{\circ}}
\newunicodechar{∪}{\ensuremath{\cup}}
\newunicodechar{∩}{\ensuremath{\cap}}
\newunicodechar{≡}{\ensuremath{\equiv}}
\newunicodechar{⊂}{\ensuremath{\subset}}
\newunicodechar{⊥}{\ensuremath{\perp}}
\newunicodechar{∃}{\ensuremath{\exists}}
\newunicodechar{∀}{\ensuremath{\forall}}
\newunicodechar{∛}{\ensuremath{\sqrt[3]{\,}}}
\newunicodechar{∜}{\ensuremath{\sqrt[4]{\,}}}
\newunicodechar{⃗}{\ensuremath{{}^{\to}}}
\newunicodechar{ⁿ}{\ifmmode^{n}\else\textsuperscript{\ensuremath{n}}\fi}
\newunicodechar{ˣ}{\ifmmode^{x}\else\textsuperscript{\ensuremath{x}}\fi}
\newunicodechar{ᵇ}{\ifmmode^{b}\else\textsuperscript{\ensuremath{b}}\fi}
\newunicodechar{ʳ}{\ifmmode^{r}\else\textsuperscript{\ensuremath{r}}\fi}
\newunicodechar{ᵘ}{\ifmmode^{u}\else\textsuperscript{\ensuremath{u}}\fi}
\newunicodechar{ʸ}{\ifmmode^{y}\else\textsuperscript{\ensuremath{y}}\fi}
\newunicodechar{ᵃ}{\ifmmode^{a}\else\textsuperscript{\ensuremath{a}}\fi}
\newunicodechar{ᵉ}{\ifmmode^{e}\else\textsuperscript{\ensuremath{e}}\fi}
\newunicodechar{ᵏ}{\ifmmode^{k}\else\textsuperscript{\ensuremath{k}}\fi}
\newunicodechar{ᵖ}{\ifmmode^{p}\else\textsuperscript{\ensuremath{p}}\fi}
\newunicodechar{ᴺ}{\ifmmode^{N}\else\textsuperscript{\ensuremath{N}}\fi}
\newunicodechar{ᶜ}{\ifmmode^{c}\else\textsuperscript{\ensuremath{c}}\fi}
\newunicodechar{ʰ}{\ifmmode^{h}\else\textsuperscript{\ensuremath{h}}\fi}
\newunicodechar{ᵠ}{\ifmmode^{q}\else\textsuperscript{\ensuremath{q}}\fi}
\newunicodechar{ₙ}{\ifmmode_{n}\else\textsubscript{\ensuremath{n}}\fi}
\newunicodechar{ₐ}{\ifmmode_{a}\else\textsubscript{\ensuremath{a}}\fi}
\newunicodechar{ᵢ}{\ifmmode_{i}\else\textsubscript{\ensuremath{i}}\fi}
\newunicodechar{ᵣ}{\ifmmode_{r}\else\textsubscript{\ensuremath{r}}\fi}
\newunicodechar{ₖ}{\ifmmode_{k}\else\textsubscript{\ensuremath{k}}\fi}
\newunicodechar{ᵦ}{\ifmmode_{\beta}\else\textsubscript{\ensuremath{\beta}}\fi}
\newunicodechar{ₓ}{\ifmmode_{x}\else\textsubscript{\ensuremath{x}}\fi}
\newfontfamily\popdisplay{ArchivoBlack-Regular.ttf}[Path=../../../fonts/]
\newfontfamily\poplogo{SpaceGrotesk-Bold.ttf}[Path=../../../fonts/]
\newfontfamily\popsemi{PlusJakartaSans-SemiBold.ttf}[Path=../../../fonts/]
\newfontfamily\popxbold{PlusJakartaSans-ExtraBold.ttf}[Path=../../../fonts/]
\newfontfamily\popxboldls{PlusJakartaSans-ExtraBold.ttf}[Path=../../../fonts/,LetterSpace=3.0]

\setlist[enumerate]{leftmargin=1.7em,itemsep=5pt,topsep=4pt}
\setlist[itemize]{leftmargin=1.7em,itemsep=4pt,topsep=4pt,label={\color{poporange}\textbullet}}
\setlength{\parindent}{0pt}
\setlength{\parskip}{5pt}
\renewcommand{\arraystretch}{1.25}

% pgfplots : style Pop par défaut
\pgfplotsset{
  every axis/.append style={axis line style={popink,line width=1pt},tick style={popink},
    grid style={popline,line width=0.5pt},label style={font=\small},tick label style={font=\footnotesize}},
  popcurve/.style={poporange,line width=1.8pt,no markers,smooth},
}
% Même style disponible dans un \draw TikZ simple (hors pgfplots)
\tikzset{popcurve/.style={poporange,line width=1.8pt}}
\tikzset{popgrid/.style={popline,very thin}}
\colorlet{popcurve}{poporange} % le nom du style est parfois utilisé comme couleur

% ============================================================
% Pill / squiggle
% ============================================================
\newcommand{\poppill}[3]{%
  \tikz[baseline=-0.55ex]\node[draw=popink,line width=1.2pt,rounded corners=2.6pt,%
    fill=#2,inner xsep=6.5pt,inner ysep=3pt]{%
    \popxboldls\fontsize{7}{7}\selectfont\color{#3}\MakeUppercase{#1}};%
}
\newcommand{\popsquiggle}[1]{%
  \tikz[baseline=-0.4ex]{
    \draw[#1,line width=2.6pt,line cap=round]
      plot [smooth] coordinates {(0,0.09) (0.34,0.01) (0.68,0.21) (1.02,0.01) (1.36,0.10)};
  }%
}
% Mot-clé surligné (marqueur orange sous le texte)
\newcommand{\cle}[1]{{\popxbold\color{popdark}#1}}

% ============================================================
% En-tête / pied
% ============================================================
\pagestyle{fancy}
\fancyhf{}
\renewcommand{\headrulewidth}{0pt}
\renewcommand{\footrulewidth}{0pt}
\lhead{\raisebox{-0.15ex}{\poplogo\fontsize{13}{13}\selectfont\color{popink}OnBuch\color{poporange}+}}
\rhead{\poppill{\DOCMATIERE\ \textbullet\ \DOCNIVEAU\ \textbullet\ Leçon \DOCLECON}{white}{popink}}
\lfoot{\popxbold\fontsize{7.6}{7.6}\selectfont\color{popmuted}\MakeUppercase{Cours signé OnBuch+}\ \textbullet\ {\popsemi\DOCTITRE}}
\rfoot{\tikz[baseline=-0.6ex]\node[rounded corners=8.5pt,fill=popink,minimum height=17pt,minimum width=17pt,inner xsep=5pt,inner sep=0]{\popxbold\fontsize{8}{8}\selectfont\color{popcream}\thepage\,/\,\pageref*{LastPage}};}

% ============================================================
% Titres
% ============================================================
\newcommand{\chaptitre}[2]{%
  \begin{tcolorbox}[enhanced,colback=poporange,colframe=popink,boxrule=2.6pt,arc=11pt,
      drop shadow=popink,left=15pt,right=15pt,top=12pt,bottom=11pt,before skip=3pt,after skip=18pt]
    \poppill{#2}{popink}{popcream}\par\vspace{9pt}
    {\raggedright\hyphenpenalty=10000\exhyphenpenalty=10000\popdisplay\fontsize{22}{24}\selectfont\color{popcream}\MakeUppercase{#1}\par}
  \end{tcolorbox}
}
% Section numérotée : pastille orange avec le numéro + titre Archivo
\newcounter{popsec}
\newcommand{\coursec}[1]{%
  \stepcounter{popsec}\setcounter{popsub}{0}%
  \par\vspace{16pt}\noindent
  \begin{minipage}{\linewidth}
  \tikz[baseline=-0.7ex]\node[draw=popink,line width=1.6pt,fill=poporange,rounded corners=4pt,minimum size=22pt,inner sep=0]{\popdisplay\fontsize{12}{12}\selectfont\color{popcream}\thepopsec};%
  \hspace{8pt}{\popdisplay\fontsize{15}{17}\selectfont\color{popink}\MakeUppercase{#1}}\par
  \vspace{4pt}\noindent{\color{poporange}\rule{36pt}{3.6pt}}
  \end{minipage}\par\nopagebreak\vspace{8pt}
}
\newcounter{popsub}
\newcommand{\courssub}[1]{%
  \stepcounter{popsub}%
  \par\vspace{9pt}\noindent{\popxbold\fontsize{11.5}{13}\selectfont\color{popink}{\color{poporange}\thepopsec.\thepopsub}\hspace{6pt}#1}\par\nopagebreak\vspace{4pt}
}

% ============================================================
% Boîtes pédagogiques — même recette : fond blanc, bordure épaisse colorée,
% ombre dure encre, badge plein. Titre optionnel [#1].
% ============================================================
\tcbset{popbase/.style={breakable,enhanced,colback=white,boxrule=2.3pt,arc=9pt,
  shadow={3pt}{-3pt}{0pt}{popink},left=14pt,right=14pt,top=11pt,bottom=11pt,before skip=11pt,after skip=15pt,
  fontupper=\popsemi\fontsize{10.6}{14.5}\selectfont\color{popink},
  fontlower=\popsemi\fontsize{10.6}{14.5}\selectfont\color{popink}}}
% Coche dessinée (le glyphe ✓ n'existe pas dans les polices du document)
\newcommand{\popcheck}[1]{\tikz[baseline=-0.5ex]\draw[#1,line width=1.6pt,line cap=round,line join=round](-0.13,0.0)--(-0.04,-0.09)--(0.13,0.1);}
\providecommand{\coche}{\popcheck{popgreen}}
\providecommand{\resultat}[1]{\par\smallskip\noindent\fcolorbox{popgreen}{popgreenL}{\parbox{\dimexpr\linewidth-2\fboxsep-2\fboxrule\relax}{\textbf{Résultat.} #1}}\par\smallskip}
\newcommand{\popboxhead}[3]{\poppill{#1}{#2}{white}\ifx\relax#3\relax\else\hspace{6pt}{\popxbold\fontsize{10.6}{12}\selectfont\color{popink}#3}\fi\par\vspace{7pt}}

\newtcolorbox{aretenir}[1][]{popbase,colframe=popgreen,before upper={\popboxhead{À retenir}{popgreen}{#1}}}
\newtcolorbox{exemplebox}[1][]{popbase,colframe=poporange,before upper={\popboxhead{Exemple}{poporange}{#1}}}
\newtcolorbox{definition}[1][]{popbase,colframe=popblue,before upper={\popboxhead{Définition}{popblue}{#1}}}
\newtcolorbox{propriete}[1][]{popbase,colframe=poppurple,before upper={\popboxhead{Propriété}{poppurple}{#1}}}
\newtcolorbox{methode}[1][]{popbase,colframe=popgold,before upper={\popboxhead{Méthode}{popgold}{#1}}}
\newtcolorbox{attention}[1][]{popbase,colframe=poppink,before upper={\popboxhead{Attention}{poppink}{#1}}}
\newtcolorbox{experience}[1][]{popbase,colframe=popblue,colback=popblueL,before upper={\popboxhead{Expérience}{popblue}{#1}}}
% « Le savais-tu ? » : carte violette pleine (contexte, culture, Cameroun)
\newtcolorbox{savaistu}[1][]{popbase,colback=poppurple,colframe=popink,
  fontupper=\popsemi\fontsize{10.4}{14.2}\selectfont\color{white},
  before upper={\poppill{Le savais-tu\,?}{white}{poppurple}\ifx\relax#1\relax\else\hspace{6pt}{\popxbold\color{white}#1}\fi\par\vspace{7pt}}}
% Exercice résolu : énoncé en haut, \tcblower puis la solution sur fond crème
\newcounter{popexo}\newcounter{popexr}
\newtcolorbox{exoresolu}[1][]{popbase,colframe=poporange,colbacklower=popcream,
  segmentation style={popink,line width=1.2pt,dashed},
  before upper={\stepcounter{popexr}\popboxhead{Exercice résolu \thepopexr}{poporange}{#1}},
  before lower={\poppill{Solution}{popink}{popcream}\par\vspace{6pt}}}
% Exercice (non résolu en place) — la correction est dans la partie Corrigés
\newtcolorbox{exercice}[1][]{popbase,colframe=popink,
  before upper={\stepcounter{popexo}\popboxhead{Exercice \thepopexo}{popink}{#1}}}
% Corrigé
\newtcolorbox{corrige}[1][]{popbase,colframe=popgreen,colback=popgreenL,
  before upper={\popboxhead{Corrigé}{popgreen}{#1}}}
% Objectifs / prérequis / bilan
\newtcolorbox{objectifs}{popbase,colframe=popink,colback=poporangeL,
  before upper={\popboxhead{Objectifs de la leçon}{popink}{}}}
\newtcolorbox{prerequis}{popbase,colframe=popink,colback=popblueL,
  before upper={\popboxhead{Prérequis}{popblue}{}}}
\newtcolorbox{bilan}{popbase,colframe=popink,colback=white,boxrule=2.8pt,
  before upper={\popboxhead{Fiche bilan}{poporange}{L'essentiel à connaître pour l'examen}}}
% Figure encadrée avec légende
\newtcolorbox{popfigure}[1][]{enhanced,breakable=false,colback=white,colframe=popline,boxrule=1.4pt,arc=8pt,
  left=10pt,right=10pt,top=10pt,bottom=8pt,before skip=10pt,after skip=12pt,halign=center,
  lower separated=false,halign lower=center,
  fontlower=\popsemi\fontsize{9}{11.5}\selectfont\color{popmuted},
  #1}
\newcommand{\legende}[1]{\par\vspace{6pt}{\popsemi\fontsize{9}{11.5}\selectfont\color{popmuted}#1}}

% ============================================================
% Signature OnBuch+
% ============================================================
\newcommand{\poplogoinline}[1]{{\poplogo\fontsize{#1}{#1}\selectfont\color{popink}OnBuch\color{poporange}+}}
\newcommand{\signatureonbuch}{%
  \begin{tcolorbox}[enhanced,colback=popink,colframe=popink,boxrule=2.2pt,arc=10pt,
      drop shadow=poporange,left=14pt,right=14pt,top=10pt,bottom=10pt,before skip=0pt,after skip=0pt]
    \tikz[baseline=-0.7ex]\node[circle,fill=popgreen,draw=popcream,line width=1.3pt,minimum size=22pt,inner sep=0]{\popcheck{white}};%
    \hspace{9pt}\begin{minipage}[c]{0.62\linewidth}
      {\popxbold\fontsize{12}{13}\selectfont\color{popcream}Signé\ \textbullet\ {\poplogo OnBuch\color{poporange}+}}\par\vspace{2pt}
      {\raggedright\popsemi\fontsize{8}{10.5}\selectfont\color{popline}\MakeUppercase{Équipe pédagogique OnBuch+}\\\MakeUppercase{Conforme au programme officiel MINESEC}\par}
    \end{minipage}\hfill
    {\poplogo\fontsize{22}{22}\selectfont\color{popcream}OnBuch\color{poporange}+}
  \end{tcolorbox}%
}

% ============================================================
% Page de garde
% #1 titre · #2 matière · #3 niveau · #4 module · #5 n° leçon · #6 durée ·
% #7 description / objectifs (texte ou itemize)
% ============================================================
\newcommand{\pagegarde}[7]{%
  \thispagestyle{empty}
  \newgeometry{a4paper,margin=1.7cm,top=1.6cm,bottom=1.4cm}%
  \begin{tikzpicture}[remember picture,overlay]
    \fill[poporange!18] ([xshift=-3.2cm,yshift=-3.6cm]current page.north east) circle (4.6cm);
    \fill[poppurple!14] ([xshift=2.4cm,yshift=7.6cm]current page.south west) circle (3.8cm);
  \end{tikzpicture}%
  {\poplogo\fontsize{30}{30}\selectfont\color{popink}OnBuch\color{poporange}+}\hfill
  \raisebox{0.6ex}{\poppill{Cours signé \textbullet\ Premium}{popink}{popcream}}\par
  \vspace{1pt}\popsquiggle{poppurple}\par
  \vspace{8pt}
  \begin{tcolorbox}[enhanced,colback=poporange,colframe=popink,boxrule=2.8pt,arc=13pt,
      drop shadow=popink,left=18pt,right=18pt,top=12pt,bottom=13pt,after skip=10pt]
    \poppill{#2 \textbullet\ #3}{popink}{popcream}\hspace{5pt}\poppill{#4}{white}{popink}\par\vspace{8pt}
    {\popdisplay\fontsize{14}{14}\selectfont\color{popink}LEÇON #5}\par\vspace{4pt}
    {\raggedright\hyphenpenalty=10000\exhyphenpenalty=10000\popdisplay\fontsize{25}{27}\selectfont\color{popcream}\MakeUppercase{#1}\par}
    \vspace{8pt}
    {\popxbold\fontsize{9.5}{9.5}\selectfont\color{popink}\MakeUppercase{Durée conseillée : #6}}
  \end{tcolorbox}
  \begin{tcolorbox}[enhanced,colback=white,colframe=popink,boxrule=2.2pt,arc=10pt,
      drop shadow=popink,left=16pt,right=16pt,top=10pt,bottom=10pt,after skip=8pt]
    \poppill{Dans cette leçon}{poppurple}{white}\par\vspace{5pt}
    \popsemi\fontsize{9.3}{12.6}\selectfont\color{popink}#7
  \end{tcolorbox}
  \noindent\begin{minipage}[t]{0.487\linewidth}
    \begin{tcolorbox}[enhanced,colback=popgreen,colframe=popink,boxrule=2.2pt,arc=10pt,
        drop shadow=popink,left=11pt,right=11pt,top=10pt,bottom=10pt,height=3.1cm,valign=center]
      \tcbox[colback=white,colframe=popink,boxrule=1.3pt,arc=5pt,boxsep=0pt,nobeforeafter,
        left=3pt,right=3pt,top=3pt,bottom=3pt,valign=center]{\qrcode[height=1.9cm]{https://play.google.com/store/apps/details?id=com.luvvix.onbuch&hl=fr}}%
      \hspace{8pt}\begin{minipage}[c]{0.5\linewidth}
        {\popdisplay\fontsize{11}{12.5}\selectfont\color{white}\MakeUppercase{Télécharge OnBuch}}\par\vspace{4pt}
        {\raggedright\popsemi\fontsize{8.4}{11}\selectfont\color{white}Gratuit \textbullet\ Google Play\\Tuteur IA, annales, concours, résultats\par}
      \end{minipage}
    \end{tcolorbox}
  \end{minipage}\hfill
  \begin{minipage}[t]{0.487\linewidth}
    \begin{tcolorbox}[enhanced,colback=poppurple,colframe=popink,boxrule=2.2pt,arc=10pt,
        drop shadow=popink,left=11pt,right=11pt,top=10pt,bottom=10pt,height=3.1cm,valign=center]
      \tikz[baseline=-0.7ex]\node[circle,fill=white,draw=popink,line width=1.3pt,minimum size=1.6cm,inner sep=0]{\popdisplay\fontsize{24}{24}\selectfont\color{poppurple}+};%
      \hspace{8pt}\begin{minipage}[c]{0.6\linewidth}
        {\popdisplay\fontsize{11}{12.5}\selectfont\color{white}\MakeUppercase{Passe à OnBuch+}}\par\vspace{4pt}
        {\raggedright\popsemi\fontsize{8.4}{11}\selectfont\color{white}Tous les cours signés, Tuteur IA illimité, fiches PDF — onglet OnBuch+ de l'app\par}
      \end{minipage}
    \end{tcolorbox}
  \end{minipage}
  \vspace{8pt}
  \popcontenu
  \vfill
  \signatureonbuch
  \restoregeometry
  \newpage
}

% Rangée « Ce cours contient » (page de garde)
\newcommand{\poptile}[3]{% #1 couleur #2 gros texte #3 légende
  \begin{tcolorbox}[enhanced,colback=#1,colframe=popink,boxrule=2pt,arc=9pt,shadow={2.5pt}{-2.5pt}{0pt}{popink},
      left=6pt,right=6pt,top=9pt,bottom=9pt,halign=center,valign=center,height=1.75cm,width=\linewidth,nobeforeafter]
    {\popdisplay\fontsize{12.5}{13}\selectfont\color{white}\MakeUppercase{#2}}\par\vspace{3pt}
    {\centering\popsemi\fontsize{7.8}{9.5}\selectfont\color{white}#3\par}
  \end{tcolorbox}}
\newcommand{\popcontenu}{%
  \par\noindent{\popxboldls\fontsize{7.5}{7.5}\selectfont\color{popmuted}CE COURS SIGNÉ CONTIENT}\par\vspace{7pt}
  \noindent\begin{minipage}[t]{0.235\linewidth}\poptile{popblue}{Cours}{complet, illustré, pas à pas}\end{minipage}\hfill
  \begin{minipage}[t]{0.235\linewidth}\poptile{poporange}{Méthodes}{et exercices résolus}\end{minipage}\hfill
  \begin{minipage}[t]{0.235\linewidth}\poptile{poppink}{Exercices}{type examen + corrigés}\end{minipage}\hfill
  \begin{minipage}[t]{0.235\linewidth}\poptile{popgreen}{Fiche bilan}{l'essentiel à retenir}\end{minipage}\par}

% Sommaire
\newtcolorbox{sommaire}{enhanced,colback=white,colframe=popink,boxrule=2.2pt,arc=10pt,
  drop shadow=popink,left=18pt,right=18pt,top=13pt,bottom=13pt,before skip=6pt,after skip=20pt,
  before upper={\poppill{Sommaire}{poppurple}{white}\par\vspace{9pt}\popsemi\fontsize{10.6}{14.5}\selectfont\color{popink}}}

% Signature de fin de cours — poussée tout en bas de la dernière page
\newcommand{\findecours}{%
  \par\vfill\nopagebreak
  \begin{center}\popsquiggle{poporange}\end{center}\vspace{4pt}
  \signatureonbuch
}
\AtBeginDocument{\def\llbracket{\mathopen{[\mkern-3.5mu[}}\def\rrbracket{\mathclose{]\mkern-3.5mu]}}} % intervalles d'entiers (glyphe ⟦ absent de la police)
% Glyphes absents de Fira Math : redéfinis à partir de glyphes disponibles
\AtBeginDocument{%
  \def\setminus{\mathbin{\backslash}}\let\smallsetminus\setminus
  \def\vdots{\mathord{\vbox{\baselineskip3.2pt\lineskiplimit0pt\kern4pt\hbox{.}\hbox{.}\hbox{.}}}}%
  \def\ddots{\mathinner{\mkern1mu\raise6.5pt\hbox{.}\mkern2mu\raise3.6pt\hbox{.}\mkern2mu\raise0.7pt\hbox{.}\mkern1mu}}%
  \def\triangle{\mathord{\scalebox{0.9}{$\symup{\Delta}$}}}%
  \def\nabla{\mathord{\raisebox{\depth}{\scalebox{1}[-1]{$\symup{\Delta}$}}}}%
  \def\checkmark{\popcheck{popdark}}%
  \def\star{\mathbin{\ast}}\def\bigstar{\mathord{\ast}}%
  \def\diamond{\mathbin{\scalebox{0.6}{$\Diamond$}}}%
  \def\bigcup{\mathop{\mathchoice{\raisebox{-0.3em}{\scalebox{1.7}{$\cup$}}}{\scalebox{1.25}{$\cup$}}{\cup}{\cup}}\displaylimits}%
  \def\bigcap{\mathop{\mathchoice{\raisebox{-0.3em}{\scalebox{1.7}{$\cap$}}}{\scalebox{1.25}{$\cap$}}{\cap}{\cap}}\displaylimits}%
  \def\lesseqqgtr{\mathrel{\lessgtr}}\def\bigoplus{\mathop{\oplus}\displaylimits}%
}
\providecommand{\ssi}{si et seulement si} % abréviation fréquente
\AtBeginDocument{\providecommand{\wideparen}{\overparen}} % arc de cercle

%% FIGFIX-BEGIN v4 — correctifs globaux des figures (docs/onbuch-generation/tools/figfix)
\usepackage{adjustbox}\usepackage{etoolbox}
% toute figure TikZ/pgfplots plus large que la ligne est réduite à la largeur disponible
\BeforeBeginEnvironment{tikzpicture}{\begin{adjustbox}{max width=\linewidth}}
\AfterEndEnvironment{tikzpicture}{\end{adjustbox}}
% légendes pgfplots lisibles par-dessus les courbes
\pgfplotsset{every axis/.append style={legend style={fill=white,fill opacity=0.92,text opacity=1,draw=black!15,inner sep=3pt}}}
% étiquettes d'arêtes (syntaxe edge["..."]) avec fond blanc
\tikzset{every edge quotes/.append style={fill=white,inner sep=1.2pt}}
%% FIGFIX-END
\begin{document}

\pagegarde{Lutte contre l'effet de serre}{SVT}{1ère D}{Module 3 : L'Éducation à l'Environnement et au Développement Durable}{ND11}{12 h}{Cette leçon explique le mécanisme de l'effet de serre naturel, montre comment les activités humaines le renforcent et provoquent le réchauffement climatique, puis présente les conséquences observées et les moyens d'action à l'échelle individuelle, nationale et internationale. Elle s'appuie sur des exemples camerounais concrets (déforestation, transport, énergie) pour former des citoyens capables d'agir.
\vspace{4pt}
\begin{multicols}{2}\small
\begin{itemize}
\item À la fin de cette leçon, je sais définir l'effet de serre et citer les principaux gaz à effet de serre (CO2, méthane, vapeur d'eau)
\item À la fin de cette leçon, je sais expliquer le mécanisme de l'effet de serre à partir d'un schéma (rayonnement solaire, rayonnement infrarouge, piégeage)
\item À la fin de cette leçon, je sais distinguer l'effet de serre naturel de son renforcement d'origine anthropique
\item À la fin de cette leçon, je sais relier une activité humaine donnée (transport, déforestation, élevage, énergie) à une émission de gaz à effet de serre
\item À la fin de cette leçon, je sais décrire au moins quatre conséquences du réchauffement climatique sur l'environnement et les populations, dont des exemples camerounais
\item À la fin de cette leçon, je sais interpréter un graphique d'évolution des températures ou des concentrations en CO2
\end{itemize}\end{multicols}}

\begin{sommaire}
\begin{multicols}{2}
\begin{enumerate}
\item L'effet de serre : un phénomène naturel indispensable
\item Le renforcement de l'effet de serre par les activités humaines
\item Les conséquences du réchauffement climatique
\item Lutter contre l'effet de serre : atténuation et adaptation
\item Les politiques climatiques : du local à l'international
\item Activité d'intégration
\item Méthodes et exercices résolus
\item Exercices
\item Corrigés des exercices
\item Fiche bilan
\end{enumerate}
\end{multicols}
\end{sommaire}

\begin{objectifs}
\begin{itemize}
\item À la fin de cette leçon, je sais définir l'effet de serre et citer les principaux gaz à effet de serre (CO2, méthane, vapeur d'eau)
\item À la fin de cette leçon, je sais expliquer le mécanisme de l'effet de serre à partir d'un schéma (rayonnement solaire, rayonnement infrarouge, piégeage)
\item À la fin de cette leçon, je sais distinguer l'effet de serre naturel de son renforcement d'origine anthropique
\item À la fin de cette leçon, je sais relier une activité humaine donnée (transport, déforestation, élevage, énergie) à une émission de gaz à effet de serre
\item À la fin de cette leçon, je sais décrire au moins quatre conséquences du réchauffement climatique sur l'environnement et les populations, dont des exemples camerounais
\item À la fin de cette leçon, je sais interpréter un graphique d'évolution des températures ou des concentrations en CO2
\item À la fin de cette leçon, je sais proposer des gestes individuels et des politiques d'atténuation adaptés au contexte camerounais
\item À la fin de cette leçon, je sais citer les grandes étapes de la coopération internationale sur le climat (Protocole de Kyoto, Accord de Paris)
\end{itemize}
\end{objectifs}

\begin{prerequis}
\begin{itemize}
\item Notions sur l'atmosphère et sa composition (gaz de l'air) vues au collège
\item Notion d'énergie et de rayonnement solaire (SVT de Seconde)
\item Écosystèmes et équilibres naturels (Module 1, début d'année de Première)
\item Lecture et interprétation de graphiques et de courbes
\item Notion de ressources naturelles et de développement durable (introduction du Module 3)
\end{itemize}
\end{prerequis}

\newpage

\coursec{L'effet de serre : un phénomène naturel indispensable}

Imagine un instant : tu es à Douala, en plein mois d'août. La pluie tombe sans discontinuer depuis trois jours. Les ruisseaux sortent de leur lit, l'eau boueuse envahit les ruelles de Bépanda et de New Bell, emportant murs en banco et véhicules. Les familles s'entassent dans les écoles transformées en centres d'hébergement d'urgence. Maintenant, transporte-toi à 1\,000 km au nord, à Maroua, en mars. Le thermomètre affiche 42 °C à l'ombre depuis des semaines. Le sol est craquelé, les points d'eau sont à sec, les plants de mil peinent à lever, et les éleveurs conduisent leurs troupeaux vers le sud, toujours plus loin, à la recherche d'un brin d'herbe. Sur la côte, à Kribi, les pêcheurs te montrent, amer, la ligne de cocotiers qui s'effondre dans l'océan : la mer a grignoté dix mètres de plage en dix ans. 

Parallèlement, dans la forêt du bassin du Congo, les troncs d'ayous et de sapelli tombent sous les tronçonneuses pour l'exportation, libérant le carbone stocké depuis des siècles. À Yaoundé comme à Bertoua, le bourdonnement incessant des groupes électrogènes diesel et la nuée de motos-taxis \emph{benskineurs} crachent leurs gaz d'échappement dans l'air que nous respirons. 

\textbf{Problématique.} Et si tous ces phénomènes — inondations côtières, sécheresses sahéliennes, érosion du trait de côte, déforestation massive, pollution urbaine — avaient une cause commune, invisible mais puissante : le \textbf{renforcement de l'effet de serre} ? Comment ce mécanisme naturel, qui rend la vie possible sur Terre depuis des milliards d'années, est-il en train de devenir une menace existentielle ? Et surtout, quels leviers d'action le Cameroun, pays de la forêt, de l'eau et du soleil, peut-il actionner pour s'adapter et atténuer ce dérèglement ?

Pour répondre, commençons par comprendre le mécanisme de base, sans lequel notre planète serait une boule de glace stérile.

\courssub{Le bilan radiatif de la Terre : l'énergie venue du Soleil}

Tout commence par le Soleil. Notre étoile envoie vers la Terre un flux colossal d'énergie, le \textbf{rayonnement solaire incident}, d'une puissance moyenne de \SI{340}{W.m^{-2}} au sommet de l'atmosphère (constante solaire divisée par 4 du fait de la géométrie sphérique). Ce rayonnement est constitué pour l'essentiel de \textbf{rayonnement visible} (lumière) et de proche infrarouge (onde courte).

\begin{experience}[Le bilan radiatif simplifié : observation et modélisation]
\textbf{Observation.} Au sommet de l'atmosphère, la Terre reçoit \SI{340}{W.m^{-2}}. Une partie est renvoyée vers l'espace sans chauffer la planète (réflexion par les nuages, les aérosols, la surface claire : glace, sable, nuages). C'est l'\textbf{albédo planétaire} (environ 0,30 en moyenne). Le reste est absorbé par l'atmosphère (ozone, vapeur d'eau, aérosols) et surtout par la surface (océans, continents, forêts).

\textbf{Hypothèse.} La Terre absorbe de l'énergie en continu. Si elle ne la restitue pas, sa température doit monter indéfiniment. Or, la température moyenne est stable (à long terme). La Terre doit donc réémettre vers l'espace une puissance exactement égale à celle qu'elle absorbe.

\textbf{Données / Modélisation.} La surface terrestre, à \SI{15}{\celsius} (\SI{288}{K}), se comporte comme un corps noir approximatif. D'après la loi de Wien ($\lambda_{\max} = b/T$), elle émet un pic de rayonnement vers \SI{10}{\micro\meter} : c'est du \textbf{rayonnement infrarouge thermique} (onde longue), invisible à l'œil nu mais ressenti comme de la chaleur.

\textbf{Interprétation.} L'atmosphère est \emph{transparente} au visible (le Soleil nous chauffe) mais \emph{opaque} à une grande partie de l'infrarouge (la Terre a du mal à se refroidir). C'est l'essence de l'effet de serre.

\textbf{Conclusion.} L'équilibre radiatif s'écrit : 
\[
\text{Absorbé (Solaire)} = \text{Émis (Infrarouge vers l'espace)}
\]
Sans atmosphère, cet équilibre donnerait une température d'équilibre de \SI{-18}{\celsius}. L'atmosphère, par ses gaz, réchauffe la surface de \SI{33}{\celsius} supplémentaires.
\end{experience}

\begin{popfigure}
\begin{tikzpicture}[scale=0.9, every node/.style={font=\small}]
    % Soleil
    \node[draw=poporange, fill=poporange!20, circle, minimum size=1.2cm, align=center] (soleil) at (0,5){\textbf{Soleil}\\\SI{5800}{K}};
    \draw[poporange, thick, ->] (soleil) -- ++(-1.5,0) node[left] {$E_{\text{incident}} = \SI{340}{W.m^{-2}}$};
    
    % Atmosphère - Couche GES
    \draw[popblue, thick, fill=popblue!10, rounded corners] (-4,2.5) rectangle (4,3.5) node[midway, text=popblue] {\textbf{Atmosphère : Gaz à effet de serre (H$_2$O, CO$_2$, CH$_4$...)}};
    
    % Surface Terre
    \draw[popgreen, thick, fill=popgreen!20, rounded corners] (-4,-1) rectangle (4,0) node[midway, text=popgreen] {\textbf{Surface terrestre (Sol, Océans, Végétation)}};
    
    % Flèches Rayonnement Solaire (Visible) - Descendantes
    \draw[poporange, thick, ->] (0,3.5) -- (0,2.5) node[midway, right] {\textcolor{poporange}{\textbf{Rayonnement visible (onde courte)}}};
    \draw[poporange, thick, ->] (-2,3.5) -- (-2,2.5);
    \draw[poporange, thick, ->] (2,3.5) -- (2,2.5);
    
    % Albédo (Réflexion)
    \draw[poporange, thick, dashed, <-] (-3,3.5) -- (-3,4.5) node[above] {\textcolor{poporange}{Albédo $\approx 30\%$ (réfléchi)}};
    \draw[poporange, thick, dashed, <-] (3,3.5) -- (3,4.5);
    
    % Absorption surface
    \draw[poporange, thick, ->] (-1,2.5) -- (-1,0) node[midway, right] {Absorbé par le sol};
    \draw[poporange, thick, ->] (1,2.5) -- (1,0);
    
    % Rayonnement Infrarouge émis par la Terre (Montantes)
    \draw[popred, thick, ->] (0,0) -- (0,2.5) node[midway, left] {\textcolor{popred}{\textbf{Rayonnement infrarouge (onde longue)}}};
    \draw[popred, thick, ->] (-2,0) -- (-2,2.5);
    \draw[popred, thick, ->] (2,0) -- (2,2.5);
    
    % Piégeage par GES (Absorption / Réémission vers le bas)
    \draw[popred, thick, ->] (0,3.5) -- (0,0) node[midway, right] {\textcolor{popred}{Réémission vers le sol (Effet de serre)}};
    \draw[popred, thick, ->] (-2,3.5) -- (-2,0);
    \draw[popred, thick, ->] (2,3.5) -- (2,0);
    
    % Sortie vers l'espace (Fenêtre atmosphérique)
    \draw[popred, thick, ->] (-3.5,3) -- (-5,3) node[left] {\textcolor{popred}{Vers l'espace (\SI{240}{W.m^{-2}})}};
    \draw[popred, thick, ->] (3.5,3) -- (5,3) node[right] {\textcolor{popred}{Vers l'espace}};
    
    % Légendes température
    \node[draw, fill=popcream, rounded corners, anchor=west] at (4.5, -1) {$T_{\text{moy}} \approx \SI{+15}{\celsius}$};
\end{tikzpicture}
\legende{Schéma du bilan radiatif terrestre. Le rayonnement solaire visible (orange) traverse l'atmosphère et chauffe le sol. Le sol réémet de l'infrarouge (rouge). Les gaz à effet de serre (GES) absorbent une grande partie de cet infrarouge et le réémettent dans toutes les directions, y compris vers le sol, réchauffant la surface. Seule une partie s'échappe vers l'espace (flèches latérales). L'albédo (flèches pointillées) renvoie 30\% de l'énergie solaire.}
\end{popfigure}

\begin{definition}[Effet de serre (phénomène naturel)]
Phénomène physique par lequel certains gaz de l'atmosphère terrestre (les \cle{gaz à effet de serre}) absorbent une partie du rayonnement infrarouge émis par la surface terrestre et la réémettent vers le bas, maintenant une température moyenne à la surface compatible avec l'eau liquide et la vie ($\approx \SI{+15}{\celsius}$).
\end{definition}

\begin{definition}[Gaz à effet de serre (GES)]
Gaz atmosphériques dont la structure moléculaire (liaisons covalentes polaires, géométrie non linéaire ou asymétrique) leur confère des modes de vibration/rotation résonnants dans l'infrarouge thermique. Ils absorbent et réémettent le rayonnement infrarouge. Les principaux GES naturels sont : la \textbf{vapeur d'eau} (\ce{H2O}), le \textbf{dioxyde de carbone} (\ce{CO2}), le \textbf{méthane} (\ce{CH4}), le \textbf{protoxyde d'azote} (\ce{N2O}) et l'\textbf{ozone} (\ce{O3}).
\end{definition}

\begin{attention}[Piège lexical : « Gaz à effet de serre » $\neq$ « Polluants locaux »]
Ne confonds pas les GES (effet \emph{global}, durée de vie longue, bien mélangés dans l'atmosphère) et les polluants atmosphériques locaux (particules fines $PM_{2,5}$, \ce{NO2}, \ce{SO2}, CO) qui affectent la qualité de l'air respiré en ville (Douala, Yaoundé) mais ont un impact climatique différent (certains refroidissent par réflexion, d'autres réchauffent). Le \ce{CO2} n'est \emph{pas} toxique à concentration atmosphérique actuelle pour la respiration, mais il est le principal pilote du réchauffement global.
\end{attention}

\courssub{Le piégeage de la chaleur : rôle des gaz à effet de serre}

Pourquoi \ce{CO2}, \ce{CH4} et \ce{H2O} absorbent-ils l'infrarouge alors que \ce{N2} (78\%) et \ce{O2} (21\%), majoritaires, ne le font pas ? C'est une question de \textbf{structure moléculaire} et de \textbf{symétrie}.

\begin{methode}[Pourquoi un gaz est-il un GES ? (Démarche d'investigation)]
\begin{enumerate}
    \item \textbf{Observer les molécules.} \ce{N2} et \ce{O2} sont diatomiques homonucléaires : liaison symétrique, pas de dipôle électrique permanent, pas de changement de dipôle lors de la vibration d'étirement. Elles sont « invisibles » à l'infrarouge.
    \item \textbf{Analyser \ce{CO2}.} Linéaire \ce{O=C=O}. Vibration symétrique : pas de dipôle (inactive IR). Vibrations \emph{asymétrique} (étirement) et \emph{de flexion} (bending) : création d'un dipôle instantané $\rightarrow$ \textbf{absorption forte vers \SI{15}{\micro\meter} (667 cm$^{-1}$)}.
    \item \textbf{Analyser \ce{H2O}.} Géométrie coudée (angle \SI{104.5}{\degree}). Dipôle permanent. Trois modes de vibration actifs IR + rotation $\rightarrow$ \textbf{bandes d'absorption larges} couvrant une grande partie du spectre IR.
    \item \textbf{Analyser \ce{CH4}.} Tétraédrique. Vibrations de déformation créent un dipôle $\rightarrow$ \textbf{absorption forte vers \SI{7.7}{\micro\meter}}.
    \item \textbf{Conclure.} Seules les molécules dont la vibration modifie le moment dipolaire électrique interagissent avec le champ électromagnétique infrarouge. C'est le cas des GES, pas des gaz majoritaires \ce{N2}/\ce{O2}.
\end{enumerate}
\end{methode}

\begin{propriete}[Puissance de réchauffement global (PRG / GWP)]
Tous les GES n'ont pas le même effet par molécule. On rapporte leur effet à celui du \ce{CO2} sur 100 ans (GWP$_{100}$, valeurs GIEC AR6) :
\begin{center}
\begin{tabularx}{\linewidth}{|l|X|c|}\hline
\textbf{Gaz} & \textbf{Principales sources naturelles / anthropiques} & \textbf{PRG$_{100}$ (ref \ce{CO2}=1)} \\ \hline
\ce{CO2} & Respiration, décomposition, combustion fossile, déforestation & 1 \\
\ce{CH4} & Zones humides, rizières, élevage (entérogénèse), fuites gaz/fossiles & \textbf{27--30} \\
\ce{N2O} & Sols, engrais azotés, combustion & \textbf{273} \\
\ce{H2O} & Évaporation océans (rétroaction, pas forçage direct) & Non applicable \\ \hline
\end{tabularx}
\end{center}
\emph{Note :} Le \ce{CH4} a une durée de vie plus courte (\SI{\sim 12}{ans}) que le \ce{CO2} (siècles), d'où l'intérêt d'agir vite sur le méthane pour un bénéfice climatique rapide.
\end{propriete}

\begin{exemplebox}[L'analogie de la serre de jardin : utile mais imparfaite]
Dans une serre horticole (verre), le soleil traverse le verre (transparent au visible). Le sol chauffe et émet de l'infrarouge. Le \textbf{verre est opaque à l'infrarouge} : il bloque la sortie de la chaleur. De plus, le verre \textbf{empêche la convection} (l'air chaud ne peut pas s'échapper par le toit).
Dans l'atmosphère, les GES agissent comme le verre pour le \textbf{rayonnement} (piégeage radiatif), mais l'air \textbf{peut convecter} (mouvements verticaux, orages, alizés). L'analogie « serre » est donc radiativement correcte, mais dynamiquement fausse : l'atmosphère n'a pas de « toit » physique qui bloque la convection.
\end{exemplebox}

\courssub{Un phénomène vital : la Terre sans effet de serre}

Si nous « éteignons » les GES dans notre modèle (atmosphère transparente à l'IR comme au visible), le bilan radiatif s'applique directement à la surface nue.

\begin{popfigure}
\begin{tikzpicture}[scale=1.1, every node/.style={font=\small}]
    % AVEC EFFET DE SERRE
    \begin{scope}[xshift=-4cm]
        \draw[popblue, thick, fill=popblue!10] (-2.5,0) circle (2cm);
        \draw[popgreen, thick, fill=popgreen!30] (-2.5,0) circle (1.5cm) node[midway, text=popgreen!80!black] {\textbf{Terre + Atmosphère}};
        \draw[poporange, thick, ->] (-0.5,1.5) -- (-0.5,0.5) node[midway, right] {Solaire};
        \draw[popred, thick, <->] (-2.5,1.5) -- (-1,2.5) node[above right] {IR piégé / Réémis};
        \node[draw=popred, fill=popred!10, rounded corners, align=center, anchor=south] at (-2.5,-2.2) {$T_{\text{moy}} = \mathbf{+15^\circ C}$\\(\SI{288}{K})\\(Eau liquide, Vie)};
    \end{scope}
    
    % SANS EFFET DE SERRE
    \begin{scope}[xshift=4cm]
        \draw[popgreen, thick, fill=popgreen!30] (2.5,0) circle (1.5cm) node[midway, text=popgreen!80!black] {\textbf{Terre nue}};
        \draw[poporange, thick, ->] (4.5,1.5) -- (4.5,0.5) node[midway, right] {Solaire};
        \draw[popred, thick, ->] (2.5,1.5) -- (2.5,2.5) node[above] {IR $\rightarrow$ Espace};
        \node[draw=popblue, fill=popblue!10, rounded corners, align=center, anchor=south] at (2.5,-2.2) {$T_{\text{moy}} = \mathbf{-18^\circ C}$\\(\SI{255}{K})\\(Glace globale, Pas de vie)};
    \end{scope}
    
    \node[align=center, font=\bfseries\large] at (0,3.5) {Comparaison : Avec / Sans effet de serre naturel};
    \draw[popgold, thick, <->] (-0.5,-1.5) -- (0.5,-1.5) node[midway, below] {\textbf{Gain : +33 °C}};
\end{tikzpicture}
\legende{Comparaison des températures moyennes d'équilibre. À gauche : Terre avec atmosphère et GES (température mesurée \SI{+15}{\celsius}). À droite : Terre hypothétique sans GES (température d'équilibre radiatif \SI{-18}{\celsius}). La différence de \SI{33}{\celsius} est l'amplitude de l'effet de serre naturel.}
\end{popfigure}

\begin{exemplebox}[Calcul de la température d'équilibre (ordre de grandeur)]
Loi de Stefan-Boltzmann : Puissance émise $P = \sigma T^4$ ($\sigma = \SI{5.67e-8}{W.m^{-2}.K^{-4}}$).
Énergie absorbée par la Terre : $(1 - \alpha) \times S_0/4 = 0,70 \times \SI{340}{W.m^{-2}} \approx \SI{238}{W.m^{-2}}$.
Équilibre : $\sigma T^4 = 238 \Rightarrow T^4 = 238 / (5,67 \times 10^{-8}) \approx 4,2 \times 10^9 \Rightarrow T \approx \SI{255}{K} = \SI{-18}{\celsius}$.
Température observée moyenne : \SI{288}{K} = \SI{+15}{\celsius}.
\textbf{Gain de l'effet de serre naturel : $\Delta T = \SI{33}{\celsius}$.}
\end{exemplebox}

\begin{aretenir}[Bilan chiffré clé]
\begin{itemize}
    \item Temp. moyenne réelle (avec effet de serre) : \textbf{$\approx \SI{+15}{\celsius}$}.
    \item Temp. d'équilibre radiatif (sans effet de serre) : \textbf{$\approx \SI{-18}{\celsius}$}.
    \item \textbf{Apport de l'effet de serre naturel : \textbf{+33 °C}.}
    \item Sans lui : océans gelés, pas de cycle de l'eau liquide, pas de vie telle que nous la connaissons.
\end{itemize}
\end{aretenir}

\courssub{Distinction cruciale : effet de serre naturel vs réchauffement climatique}

C'est le point le plus souvent mal compris, y compris dans les copies d'examen.

\begin{attention}[Erreur majeure au Probatoire : Confondre cause et conséquence]
\begin{itemize}
    \item \textbf{Effet de serre} = \textbf{Mécanisme naturel}, permanent, bénéfique, indispensable. Existe depuis 4 milliards d'années. « La couverture thermique de la Terre ».
    \item \textbf{Renforcement de l'effet de serre} = \textbf{Perturbation anthropique}. Augmentation de la concentration des GES (surtout \ce{CO2}, \ce{CH4}, \ce{N2O}) due aux activités humaines depuis l'ère industrielle ($\sim$ 1750).
    \item \textbf{Réchauffement climatique} = \textbf{Conséquence} observable : hausse de la température moyenne globale ($\approx \SI{+1.2}{\celsius}$ depuis 1850--1900), fonte des glaces, montée des mers, événements extrêmes.
\end{itemize}
\textbf{Consigne de rédaction :} N'écris jamais « l'effet de serre augmente ». Écris : « la \textbf{concentration des GES augmente}, ce qui \textbf{renforce l'effet de serre}, entraînant un \textbf{réchauffement climatique} ».
\end{attention}

\begin{savaistu}[Le pilote et le passager : \ce{CO2} vs \ce{H2O}]
La vapeur d'eau (\ce{H2O}) est le \textbf{premier GES} par sa contribution à l'effet de serre naturel (\SI{\sim 50}{\%} de l'effet total, contre \SI{\sim 20}{\%} pour le \ce{CO2}). Mais sa concentration est \textbf{contrôlée par la température} (loi de Clausius-Clapeyron : l'air chaud contient plus de vapeur). C'est une \textbf{rétroaction} (feedback) : le \ce{CO2} (émis par l'homme, durée de vie longue) force le réchauffement initial $\rightarrow$ T° monte $\rightarrow$ évaporation monte $\rightarrow$ \ce{H2O} monte $\rightarrow$ effet de serre amplifié.
\textbf{Image :} Le \ce{CO2} anthropique est le \textbf{pilote} (forçage), la vapeur d'eau est le \textbf{passager} (rétroaction rapide). C'est pourquoi le GIEC cible le \ce{CO2} (et \ce{CH4}, \ce{N2O}) dans les politiques d'atténuation.
\end{savaistu}

\courssub{Synthèse et transition vers la section 2}

\begin{aretenir}[Ce qu'il faut retenir de la Section 1]
\begin{enumerate}
    \item L'effet de serre est un \textbf{phénomène radiatif naturel} : l'atmosphère laisse passer le visible solaire, absorbe l'infrarouge terrestre et le réémet vers le sol.
    \item Les \textbf{GES} (\ce{H2O}, \ce{CO2}, \ce{CH4}, \ce{N2O}, \ce{O3}) sont les « briques » de ce mécanisme ; \ce{N2} et \ce{O2} sont transparents à l'IR.
    \item Sans effet de serre : $T_{\text{moy}} = \SI{-18}{\celsius}$. Avec : $T_{\text{moy}} = \SI{+15}{\celsius}$. Gain = \textbf{+33 °C}. La vie est possible grâce à lui.
    \item \textbf{Problème actuel :} Ce n'est pas l'effet de serre qui est « mauvais », c'est son \textbf{renforcement} par l'ajout massif de GES d'origine humaine (combustion fossile, déforestation, agriculture intensive) qui déséquilibre le bilan radiatif.
\end{enumerate}
\end{aretenir}

\begin{exoresolu}[Interprétation d'un spectre d'absorption atmosphérique]
\textbf{Énoncé.} On donne ci-dessous un spectre simplifié de la transmissivité de l'atmosphère terrestre en fonction de la longueur d'onde $\lambda$.
\begin{itemize}
    \item Fenêtre visible ($0,4$--$0,7\ \mu m$) : transmissivité $\approx 80\%$ (trous d'absorption O$_3$, H$_2$O, O$_2$).
    \item Infrarouge thermique ($5$--$20\ \mu m$) : bandes d'absorption fortes centrées sur $6,3\ \mu m$ (H$_2$O), $9,6\ \mu m$ (O$_3$), $15\ \mu m$ (CO$_2$), $7,7\ \mu m$ (CH$_4$).
    \item « Fenêtre atmosphérique » $8$--$13\ \mu m$ : transmissivité élevée (sauf O$_3$ à $9,6$).
\end{itemize}
\textbf{Questions.}
1. Expliquer pourquoi le rayonnement solaire réchauffe la surface alors que le rayonnement terrestre peine à s'échapper.
2. Identifier sur le spectre la « fenêtre atmosphérique » et expliquer son rôle dans le refroidissement nocturne.
3. Prédire l'effet d'une augmentation de \ce{CO2} sur la bande à $15\ \mu m$ et sur le flux sortant vers l'espace.

\tcblower
\textbf{Solution détaillée.}
1. \textbf{Spectre solaire (pic $\sim 0,5\ \mu m$) :} Se situe dans la fenêtre de haute transmissivité (visible). L'atmosphère est quasi-transparente $\rightarrow$ l'énergie atteint le sol $\rightarrow$ absorption $\rightarrow$ élévation T°.
\textbf{Spectre terrestre (pic $\sim 10\ \mu m$) :} Se situe dans l'infrarouge où les GES ont de fortes bandes d'absorption. L'atmosphère est opaque $\rightarrow$ le rayonnement est absorbé par les GES $\rightarrow$ réémission isotrope (moitié vers l'espace, moitié vers le sol) $\rightarrow$ piégeage radiatif.
2. \textbf{Fenêtre $8$--$13\ \mu m$ :} C'est la seule bande IR large où l'atmosphère est transparente (hors bande O$_3$). Le rayonnement émis par le sol dans cette bande s'échappe \textbf{directement vers l'espace} sans être absorbé. C'est le principal « radiateur » de la Terre. La nuit, par ciel clair, le sol perd beaucoup de chaleur par cette fenêtre $\rightarrow$ fort refroidissement radiatif (gelées). Par ciel nuageux, les nuages (gouttelettes d'eau) bouchent cette fenêtre $\rightarrow$ refroidissement limité.
3. \textbf{Augmentation \ce{CO2} (bande $15\ \mu m$) :} La bande est déjà saturée au centre (optique épaisse). L'augmentation de concentration \textbf{élargit les ailes de la bande d'absorption} (élargissement par pression et concentration) et augmente l'altitude d'émission effective (plus haute, plus froide). Le flux sortant dans cette bande \textbf{diminue} $\rightarrow$ déséquilibre radiatif positif (forçage) $\rightarrow$ la Terre doit se réchauffer pour rétablir l'équilibre (émettre plus dans la fenêtre $8$--$13\ \mu m$ et aux ailes des bandes).
\end{exoresolu}

\begin{exercice}[Application : Bilan radiatif et température]
On considère une Terre hypothétique sans atmosphère (albédo $\alpha = 0,30$ identique).
1. Calculer la température d'équilibre $T_{\text{eq}}$ de cette Terre (constante solaire $S_0 = \SI{1361}{W.m^{-2}}$, $\sigma = \SI{5.67e-8}{W.m^{-2}.K^{-4}}$).
2. Comparer à la température moyenne observée ($T_{\text{obs}} = \SI{288}{K}$). Commenter la différence.
3. Si l'albédo augmente à $0,35$ (plus de nuages/glaces), que devient $T_{\text{eq}}$ ? Quel type de rétroaction est-ce (positive ou négative) par rapport à un réchauffement initial ?
\end{exercice}

\begin{corrige}[Exercice : Bilan radiatif et température]
1. Puissance absorbée par unité de surface : $P_{\text{abs}} = (1-\alpha) S_0/4 = 0,70 \times 1361 / 4 = \SI{238,2}{W.m^{-2}}$.
   Équilibre : $\sigma T_{\text{eq}}^4 = P_{\text{abs}} \Rightarrow T_{\text{eq}} = (238,2 / (5,67 \times 10^{-8}))^{1/4} \approx \SI{254,8}{K} \approx \SI{-18,3}{\celsius}$.
2. $T_{\text{obs}} = \SI{288}{K} (\SI{+15}{\celsius})$. Différence $\Delta T = \SI{33,2}{K}$. C'est l'\textbf{amplitude de l'effet de serre naturel}. L'atmosphère avec ses GES réchauffe la surface de 33 °C par rapport au corps noir nu.
3. Nouveau $P_{\text{abs}} = 0,65 \times 1361 / 4 = \SI{221,2}{W.m^{-2}}$.
   Nouveau $T_{\text{eq}} = (221,2 / (5,67 \times 10^{-8}))^{1/4} \approx \SI{251,1}{K} (\SI{-22,0}{\celsius})$.
   L'augmentation de l'albédo \textbf{baisse} la température d'équilibre. Si un réchauffement initial cause une fonte des glaces (albédo $\downarrow$), c'est une rétroaction \textbf{positive} (amplification). Si un réchauffement cause plus de nuages bas (albédo $\uparrow$), c'est une rétroaction \textbf{négative} (atténuation). Le signe net des rétroactions nuageuses reste la plus grande incertitude des modèles climatiques.
\end{corrige}

\coursec{Le renforcement de l'effet de serre par les activités humaines}

\medskip
\noindent Dans la section précédente, nous avons vu que l'\cle{effet de serre naturel} est un phénomène vital : sans lui, la température moyenne de la Terre serait d'environ $-18\,^\circ\mathrm{C}$, incompatible avec la vie telle que nous la connaissons. Les \cle{gaz à effet de serre} (GES) — principalement la vapeur d'eau ($\ce{H2O}$), le dioxyde de carbone ($\ce{CO2}$) et le méthane ($\ce{CH4}$) — agissent comme une couverture thermique en absorbant une partie du rayonnement infrarouge émis par la surface terrestre et en le réémettant vers le bas.

\medskip
\noindent Cependant, depuis le début de l'ère industrielle (milieu du \textsc{xixe} siècle), les activités humaines ont massivement augmenté les concentrations atmosphériques de ces gaz. Ce \cle{renforcement de l'effet de serre} (ou effet de serre additionnel) rompt l'équilibre radiatif planétaire et entraîne un \cle{réchauffement climatique} global. Comprenons comment nos sociétés produisent ces gaz, en nous appuyant sur des observations concrètes, des données scientifiques robustes et le contexte camerounais.

\courssub{Le principe : des émissions anthropiques qui s'accumulent dans l'atmosphère}

\begin{definition}[Émissions anthropiques]
On appelle \textbf{émissions anthropiques} les rejets de gaz à effet de serre (et d'autres polluants) dans l'atmosphère qui résultent directement des activités humaines (combustion d'énergies fossiles, agriculture, déforestation, industrie, traitement des déchets), par opposition aux sources naturelles (volcans, respiration des êtres vivants, décomposition naturelle dans les zones humides, etc.).
\end{definition}

\medskip
\noindent Le mécanisme est simple mais implacable : 
\begin{enumerate}
    \item Les activités humaines rejettent du $\ce{CO2}$, du $\ce{CH4}$, du $\ce{N2O}$ (protoxyde d'azote) et d'autres GES dans l'atmosphère.
    \item Ces gaz s'y accumulent car les \cle{puits de carbone} naturels (océans, forêts, sols) n'absorbent qu'une fraction de ces émissions excédentaires (environ la moitié pour le $\ce{CO2}$).
    \item La concentration atmosphérique de ces gaz augmente, renforçant l'absorption du rayonnement infrarouge sortant.
    \item Le bilan radiatif de la Terre devient positif (plus d'énergie entrante que sortante) : le système climatique se réchauffe jusqu'à rétablir l'équilibre à une température plus élevée.
\end{enumerate}

\begin{propriete}[Cause principale du réchauffement observé — consensus GIEC]
L'augmentation des concentrations de gaz à effet de serre dans l'atmosphère depuis l'ère préindustrielle, due aux activités humaines, est la cause principale du réchauffement climatique observé depuis le milieu du \textsc{xxe} siècle. Cette conclusion, établie par le \textbf{GIEC} (Groupe d'experts intergouvernemental sur l'évolution du climat) sur la base de milliers d'études, repose sur la concordance entre : 
\begin{itemize}
    \item l'évolution temporelle des concentrations de GES (mesurées dans les carottes de glace et directement à l'observatoire de Mauna Loa depuis 1958) ;
    \item l'évolution de la température moyenne mondiale ;
    \item la signature isotopique du carbone atmosphérique (excès de \ce{^{12}C} typique de la combustion des énergies fossiles) ;
    \item la modélisation climatique qui ne reproduit le réchauffement observé qu'en intégrant les forçages anthropiques.
\end{itemize}
Cette propriété est \textbf{admise} au Probatoire (pas de démonstration exigible, mais il faut savoir la citer et l'expliquer).
\end{propriete}

\courssub{Source majeure n°1 : la combustion des énergies fossiles}

\paragraph{Observation — Situation concrète.}
À Douala, Yaoundé ou Bafoussam, la circulation dense des motos-taxis (\emph{benskin}), des taxis, des cars et des poids lourds rejette en permanence des gaz d'échappement. Dans les zones industrielles (Bonabéri, Bassa), les usines tournent souvent grâce à des groupes électrogènes diesel ou au fioul lourd, surtout lors des délestages d'Eneo. Au niveau mondial, les centrales à charbon, à gaz et au pétrole produisent encore plus de 60\,\% de l'électricité.

\paragraph{Hypothèse.}
La combustion de ces hydrocarbures (pétrole, charbon, gaz naturel) transforme le carbone stocké depuis des millions d'années sous forme fossile en dioxyde de carbone ($\ce{CO2}$) qui s'accumule dans l'atmosphère.

\paragraph{Données — Réaction chimique et bilan.}
La combustion complète d'un hydrocarbure générique $\ce{C_xH_y}$ s'écrit :
\[
\ce{C_xH_y + (x + y/4) O2 -> x CO2 + (y/2) H2O}
\]
Pour l'essence (approximée par l'octane $\ce{C8H18}$) :
\[
\ce{2 C8H18 + 25 O2 -> 16 CO2 + 18 H2O}
\]
\emph{Masse molaire : $\ce{C8H18} \approx 114\,\mathrm{g/mol}$ ; $\ce{CO2} \approx 44\,\mathrm{g/mol}$.}
$2 \times 114 = 228\,\mathrm{g}$ d'essence produisent $16 \times 44 = 704\,\mathrm{g}$ de $\ce{CO2}$.
Soit \textbf{1 kg d'essence $\rightarrow$ environ 3,09 kg de $\ce{CO2}$}.
Comme la densité de l'essence est d'environ $0,75\,\mathrm{kg/L}$ : \textbf{1 L d'essence brûlé $\rightarrow$ environ 2,3 kg de $\ce{CO2}$}.

\begin{exemplebox}[Émissions d'un moto-taxi à Douala]
Un moto-taxi parcourt en moyenne $100\,\mathrm{km}$ par jour. Sa consommation est d'environ $3\,\mathrm{L/100\,km}$.
\begin{itemize}
    \item Consommation journalière : $3\,\mathrm{L}$.
    \item Émission journalière de $\ce{CO2}$ : $3 \times 2,3 = 6,9\,\mathrm{kg}$.
    \item Émission annuelle (300 jours de travail) : $6,9 \times 300 \approx 2\,070\,\mathrm{kg} \approx \textbf{2,1 tonnes de $\ce{CO2}$ par an}$.
\end{itemize}
Avec plus de 200\,000 motos-taxis au Cameroun, le secteur du transport urbain émet plusieurs millions de tonnes de $\ce{CO2}$ chaque année.
\end{exemplebox}

\begin{methode}[Relier une activité humaine à une émission de GES]
Pour analyser n'importe quelle activité, suivez ces trois étapes :
\begin{enumerate}
    \item \textbf{Identifier le procédé} : combustion (transport, industrie, production d'électricité, cuisine au bois/charbon), fermentation (élevage, riziculture), décomposition anaérobie (décharges), réaction chimique industrielle (cimenterie, engrais).
    \item \textbf{Identifier le(s) gaz produit(s)} : $\ce{CO2}$ (combustion, respiration, ciment), $\ce{CH4}$ (fermentation entérique, riziculture, décharges, fuites gazières), $\ce{N2O}$ (engrais azotés, combustion), gaz fluorés (réfrigération, climatisation).
    \item \textbf{Préciser l'effet sur l'effet de serre} : augmentation de la concentration atmosphérique $\rightarrow$ renforcement du piégeage infrarouge $\rightarrow$ forçage radiatif positif.
\end{enumerate}
\end{methode}

\courssub{Source majeure n°2 : la déforestation et les changements d'affectation des sols}

\paragraph{Observation — Contexte camerounais.}
Le Cameroun abrite une partie de la \cle{forêt du bassin du Congo}, deuxième massif forestier tropical au monde après l'Amazonie. Pourtant, le pays perd environ 200\,000 à 300\,000 hectares de forêt par an (données FAO/Global Forest Watch). Les causes : l'agriculture sur brûlis (culture itinérante sur brûlis), l'exploitation forestière (bois d'œuvre et bois énergie), l'expansion des plantations (palmier à huile, hévéa, bananier), l'urbanisation et les infrastructures.

\begin{savaistu}[Le Cameroun, gardien du « deuxième poumon »]
La forêt du bassin du Congo stocke environ 30 à 40 milliards de tonnes de carbone dans sa biomasse et ses sols (tourbières de la Cuvette Centrale). Sa déforestation libère ce carbone sous forme de $\ce{CO2}$ et réduit la capacité de la planète à absorber le $\ce{CO2}$ anthropique. Protéger cette forêt est un enjeu climatique \textbf{mondial}, pas seulement national.
\end{savaistu}

\paragraph{Mécanisme double — Le piège à éviter.}

\begin{attention}[Déforestation : un double effet « kiss-cool » sur le climat]
La déforestation agit \textbf{de deux manières simultanées} sur l'effet de serre :
\begin{enumerate}
    \item \textbf{Perte d'un puits de carbone} : les arbres en croissance absorbent le $\ce{CO2}$ atmosphérique par photosynthèse ($\ce{CO2 + H2O ->[$\text{lumière}$] C6H12O6 + O2}$). En les coupant, on supprime ce pompage naturel.
    \item \textbf{Liberation du stock de carbone} : le carbone accumulé dans le bois (tronc, branches, racines) et dans les sols forestiers (humus, tourbe) est oxydé. Si le bois est brûlé (agriculture sur brûlis, charbon de bois) ou s'il pourrit au sol, le carbone retourne sous forme de $\ce{CO2}$ (et un peu de $\ce{CH4}$) dans l'atmosphère.
\end{enumerate}
\textbf{Erreur fréquente au Probatoire :} dire seulement « la déforestation libère du $\ce{CO2}$ » sans mentionner la perte simultanée de la capacité d'absorption future. Les deux aspects sont indispensables pour une réponse complète.
\end{attention}

\paragraph{Données chiffrées.}
Le secteur « Utilisation des terres, changement d'affectation des terres et foresterie » (UTCATF ou \emph{LULUCF} en anglais) représente environ \textbf{12 à 15\,\% des émissions anthropiques mondiales de $\ce{CO2}$} (GIEC, 6\textsuperscript{e} rapport). Au Cameroun, ce secteur est souvent le \textbf{premier émetteur national} de GES, devant l'énergie.

\courssub{Source majeure n°3 : l'agriculture — élevage et riziculture (méthane $\ce{CH4}$)}

\paragraph{Observation.}
Dans le Nord-Cameroun (régions de l'Adamaoua, du Nord, de l'Extrême-Nord), l'élevage bovin est une activité majeure (plus de 6 millions de têtes). On y pratique aussi la riziculture irriguée (plaines de la Bénoué, du Logone, de Maga). Ces deux activités sont de puissantes sources de \cle{méthane} ($\ce{CH4}$).

\paragraph{Mécanisme biologique — Fermentation entérique.}
Les ruminants (bovins, ovins, caprins) possèdent un rumen où des milliards de micro-organismes (archées méthanogènes) décomposent la cellulose par fermentation anaérobie :
\[
\ce{C6H12O6 -> 2 CH3COOH + 2 CO2 + 4 H2} \quad \text{(fermentation acétique)}
\]
\[
\ce{CO2 + 4 H2 -> CH4 + 2 H2O} \quad \text{(méthanogenèse : réduction du $\ce{CO2}$ par l'hydrogène)}
\]
Le $\ce{CH4}$ est éructé (rots) par l'animal. Une vache émet en moyenne $80$ à $120\,\mathrm{kg}$ de $\ce{CH4}$ par an.

\paragraph{Mécanisme biologique — Riziculture irriguée.}
Dans les rizières inondées, le sol est anaérobie. La matière organique (résidus de riz, engrais organiques) est décomposée par des bactéries méthanogènes :
\[
\ce{2 CH3COOH -> CH4 + CO2} \quad \text{(acétoclastique)} \qquad
\ce{CO2 + 4 H2 -> CH4 + 2 H2O} \quad \text{(hydrogénotrophe)}
\]
Le $\ce{CH4}$ diffuse vers l'atmosphère à travers les tiges de riz (aérenchymes) et par ébullition dans l'eau.

\begin{propriete}[Pouvoir de réchauffement global — PRG (GWP en anglais)]
Le \textbf{PRG} compare l'effet de serre d'un gaz à celui du $\ce{CO2}$ sur un horizon de temps donné (généralement 100 ans).
\begin{itemize}
    \item $\ce{CO2}$ : PRG = 1 (référence).
    \item $\ce{CH4}$ : PRG $\approx$ \textbf{27 à 29,8} (valeur GIEC AR6 : 27,9 sur 100 ans, 79,7 sur 20 ans). \emph{Notion admise : « $\ce{CH4}$ est environ 25 à 28 fois plus puissant que $\ce{CO2}$ sur 100 ans ».}
    \item $\ce{N2O}$ : PRG $\approx$ \textbf{273} (sur 100 ans).
\end{itemize}
Bien que sa concentration soit plus faible ($\sim 1,9\,\mathrm{ppm}$ vs $420\,\mathrm{ppm}$ pour $\ce{CO2}$), le $\ce{CH4}$ contribue à environ \textbf{30\,\% du forçage radiatif anthropique total} à ce jour.
\end{propriete}

\courssub{Source n°4 : déchets et engrais azotés ($\ce{CH4}$ et $\ce{N2O}$)}

\paragraph{Décharges et déchets organiques.}
Dans les décharges non contrôlées (majorité des villes camerounaises : Nkolfoulou à Yaoundé, PK10 à Douala, etc.), les déchets organiques (restes alimentaires, déchets verts, papiers, cartons) se décomposent en l'absence d'oxygène (milieu anaérobie). Les bactéries méthanogènes produisent du \cle{biogaz} : $\approx 50\text{--}60\,\%\,\ce{CH4}$, $40\text{--}50\,\%\,\ce{CO2}$, traces de $\ce{H2S}$. Sans captage, ce $\ce{CH4}$ s'échappe dans l'atmosphère.

\paragraph{Engrais azotés — Protoxyde d'azote ($\ce{N2O}$).}
L'utilisation massive d'engrais de synthèse (urée, nitrate d'ammonium, NPK) dans les plantations (coton, maïs, banane, palmier à huile) et l'élevage intensif (déjections) apportent de l'azote réactif dans les sols. Deux processus microbiens produisent du $\ce{N2O}$ :
\begin{itemize}
    \item \textbf{Nitrification} (aérobie) : $\ce{NH4+ -> NO2- -> NO3-}$ (fuite de $\ce{N2O}$).
    \item \textbf{Dénitrification} (anaérobie) : $\ce{NO3- -> NO2- -> NO -> N2O -> N2}$ (fuite de $\ce{N2O}$).
\end{itemize}
Le $\ce{N2O}$ est un GES très puissant (PRG $\approx 273$) et détruit aussi la couche d'ozone stratosphérique.

\courssub{Preuves observationnelles : l'évolution historique des concentrations et des températures}

\paragraph{Données — La courbe de Keeling et les carottes de glace.}
Depuis 1958, Charles David Keeling mesure en continu la concentration en $\ce{CO2}$ à l'observatoire de Mauna Loa (Hawaï, $3\,400\,\mathrm{m}$ d'altitude, loin des sources locales). Avant 1958, les concentrations sont reconstruites à partir de \textbf{bulles d'air piégées dans les carottes de glace} de l'Antarctique (Vostok, EPICA, Dome C).

\begin{popfigure}
\begin{tikzpicture}
\begin{axis}[
    width=\linewidth,
    height=0.55\linewidth,
    xlabel={Année},
    ylabel={Concentration en $\ce{CO2}$ (ppm)},
    xmin=1750, xmax=2025,
    ymin=260, ymax=430,
    xtick={1750,1800,1850,1900,1950,2000,2020},
    ytick={280,300,320,340,360,380,400,420},
    grid=major,
    grid style={popmuted},
    axis line style={popdark},
    tick style={popdark},
    label style={font=\small,popdark},
    tick label style={font=\footnotesize,popdark},
    legend style={at={(0.02,0.98)},anchor=north west,font=\footnotesize,fill=popcream,draw=popline},
]
% Données simplifiées : carottes de glace (pré-1958) + Mauna Loa (post-1958)
\addplot[popcurve, very thick, domain=1750:1958, samples=100] 
    ({x}, {280 + 0.005*(x-1750) + 0.00002*(x-1750)^2});
\addplot[poporange, very thick, domain=1958:2024, samples=100] 
    ({x}, {315 + 1.5*(x-1958) + 0.01*(x-1958)^2});
\addplot[only marks, mark=*, mark size=1.5, popblue, domain=1958:2024, samples=14] 
    ({x}, {315 + 1.5*(x-1958) + 0.01*(x-1958)^2 + 0.5*rand});
\legend{Reconstitution par carottes de glace, Mesures directes (Mauna Loa), Données mensuelles (variabilité saisonnière)}
\end{axis}
\end{tikzpicture}
\legende{Évolution de la concentration atmosphérique en $\ce{CO2}$ de 1750 à 2024. La courbe montre une stabilité autour de 280 ppm à l'ère préindustrielle, suivie d'une hausse exponentielle depuis 1850, dépassant 420 ppm en 2024. La petite oscillation annuelle (dents de scie) visible après 1958 correspond au cycle saisonnier de la photosynthèse hémisphère Nord.}
\end{popfigure}

\paragraph{Données — Anomalie de température moyenne mondiale.}
L'anomalie de température est l'écart par rapport à une période de référence (souvent 1951–1980 ou 1850–1900).

\begin{popfigure}
\begin{tikzpicture}
\begin{axis}[
    width=\linewidth,
    height=0.55\linewidth,
    xlabel={Année},
    ylabel={Anomalie de température ($^\circ\mathrm{C}$)},
    xmin=1850, xmax=2025,
    ymin=-0.5, ymax=1.5,
    xtick={1850,1900,1950,2000,2020},
    ytick={-0.4,-0.2,0,0.2,0.4,0.6,0.8,1.0,1.2,1.4},
    grid=major,
    grid style={popmuted},
    axis line style={popdark},
    tick style={popdark},
    label style={font=\small,popdark},
    tick label style={font=\footnotesize,popdark},
    legend style={at={(0.02,0.98)},anchor=north west,font=\footnotesize,fill=popcream,draw=popline},
]
% Courbe de température simplifiée (type HadCRUT / NASA GISTEMP)
\addplot[poppurple, very thick, domain=1850:2024, samples=200] 
    ({x}, {-0.3 + 0.0008*(x-1850) + 0.000005*(x-1850)^2 + 0.1*sin(deg(0.5*(x-1850))) + 0.05*sin(deg(2*(x-1850)))});
\addplot[popblue, thick, dashed, domain=1850:2024, samples=100] 
    ({x}, {-0.3 + 0.0008*(x-1850) + 0.000005*(x-1850)^2});
\legend{Anomalie observée (variabilité naturelle + tendance), Tendance longue (réchauffement anthropique)}
\end{axis}
\end{tikzpicture}
\legende{Évolution de l'anomalie de température moyenne mondiale de 1850 à 2024 (référence 1850–1900). La courbe violette montre la variabilité interannuelle (El Niño/La Niña, volcans) superposée à une tendance de fond claire. Le réchauffement observé est d'environ $+1,2\,^\circ\mathrm{C}$ en 2023/2024 par rapport à l'ère préindustrielle.}
\end{popfigure}

\paragraph{Interprétation — Corrélation et causalité.}
\begin{itemize}
    \item Les deux courbes présentent une \textbf{corrélation temporelle frappante} : hausse concomitante du $\ce{CO2}$ et de la température depuis 1850, accélération nette après 1950 (« grande accélération »).
    \item La physique radiative (loi de Stefan-Boltzmann, spectroscopie d'absorption) établit le \textbf{lien causal} : plus de $\ce{CO2}$ $\rightarrow$ plus d'absorption infrarouge $\rightarrow$ forçage radiatif positif $\rightarrow$ réchauffement.
    \item Les modèles climatiques qui intègrent \textbf{seulement} les forçages naturels (solaire, volcanique) ne reproduisent \textbf{pas} le réchauffement observé après 1970. Seuls les modèles intégrant les forçages anthropiques (GES, aérosols, changement d'albédo) y parviennent.
\end{itemize}

\begin{aretenir}[Preuves du renforcement anthropique de l'effet de serre]
\begin{enumerate}
    \item \textbf{Concentration $\ce{CO2}$} : $280\,\mathrm{ppm}$ (1750) $\rightarrow$ $>420\,\mathrm{ppm}$ (2024), hausse de plus de $50\,\%$.
    \item \textbf{Signature isotopique} : baisse du rapport \ce{^{13}C}/\ce{^{12}C} et du \ce{^{14}C} (effet Suess) $\rightarrow$ origine fossile (végétaux anciens appauvris en \ce{^{13}C} et sans \ce{^{14}C}).
    \item \textbf{Bilan de masse} : les émissions anthropiques cumulées ($\approx 2\,500\,\mathrm{Gt\,CO2}$ depuis 1850) sont environ \textbf{deux fois supérieures} à l'augmentation observée dans l'atmosphère ; le reste a été absorbé par océans et biosphère terrestre.
    \item \textbf{Température} : $+1,2\,^\circ\mathrm{C}$ depuis 1850–1900, dernière décade la plus chaude jamais enregistrée.
    \item \textbf{Consensus} : $>99\,\%$ des articles scientifiques climatologiques et tous les grands organismes scientifiques nationaux/internationaux (GIEC, Académie des sciences, NASA, NOAA, WMO, etc.) s'accordent sur l'origine humaine.
\end{enumerate}
\end{aretenir}

\courssub{Répartition sectorielle des émissions anthropiques mondiales}

\paragraph{Diagramme de synthèse.}
Le graphique ci-dessous montre la répartition approximative des émissions totales de GES (en $\ce{CO2}$-équivalent, intégrant le PRG) par grand secteur d'activité (données GIEC AR6, WGIII, 2019, arrondies).

\begin{popfigure}
\begin{tikzpicture}
% Données : Énergie 30%, Transport 16%, Industrie 20%, Agriculture 12%, Forêt/UTCATF 11%, Bâtiments 6%, Déchets 3%, Autres 2%
\definecolor{energie}{HTML}{E74C3C}
\definecolor{transport}{HTML}{E67E22}
\definecolor{industrie}{HTML}{F39C12}
\definecolor{agriculture}{HTML}{27AE60}
\definecolor{foret}{HTML}{1E8A49}
\definecolor{batiments}{HTML}{3498DB}
\definecolor{dechets}{HTML}{9B59B6}
\definecolor{autres}{HTML}{95A5A6}

% Camembert avec calcul d'angles manuel pour robustesse
% Part 1: Énergie 30% -> 108 deg
\fill[energie] (0,0) -- (0:3cm) arc (0:108:3cm) -- cycle;
\node[font=\footnotesize, text=popdark, align=center] at (54:4.2cm) {\textbf{Énergie (prod. élec., chaleur)}\\30\%};

% Part 2: Transport 16% -> 57.6 deg (start 108, end 165.6)
\fill[transport] (0,0) -- (108:3cm) arc (108:165.6:3cm) -- cycle;
\node[font=\footnotesize, text=popdark, align=center] at (136.8:4.2cm) {\textbf{Transport}\\16\%};

% Part 3: Industrie 20% -> 72 deg (start 165.6, end 237.6)
\fill[industrie] (0,0) -- (165.6:3cm) arc (165.6:237.6:3cm) -- cycle;
\node[font=\footnotesize, text=popdark, align=center] at (201.6:4.2cm) {\textbf{Industrie (procédés + énergie)}\\20\%};

% Part 4: Agriculture 12% -> 43.2 deg (start 237.6, end 280.8)
\fill[agriculture] (0,0) -- (237.6:3cm) arc (237.6:280.8:3cm) -- cycle;
\node[font=\footnotesize, text=popdark, align=center] at (259.2:4.2cm) {\textbf{Agriculture (CH$_4$, N$_2$O)}\\12\%};

% Part 5: Forêt/UTCATF 11% -> 39.6 deg (start 280.8, end 320.4)
\fill[foret] (0,0) -- (280.8:3cm) arc (280.8:320.4:3cm) -- cycle;
\node[font=\footnotesize, text=popdark, align=center] at (300.6:4.2cm) {\textbf{Forêt \& Changement affect. sols}\\11\%};

% Part 6: Bâtiments 6% -> 21.6 deg (start 320.4, end 342)
\fill[batiments] (0,0) -- (320.4:3cm) arc (320.4:342:3cm) -- cycle;
\node[font=\footnotesize, text=popdark, align=center] at (331.2:4.2cm) {\textbf{Bâtiments (direct)}\\6\%};

% Part 7: Déchets 3% -> 10.8 deg (start 342, end 352.8)
\fill[dechets] (0,0) -- (342:3cm) arc (342:352.8:3cm) -- cycle;
\node[font=\footnotesize, text=popdark, align=center] at (347.4:4.2cm) {\textbf{Déchets}\\3\%};

% Part 8: Autres 2% -> 7.2 deg (start 352.8, end 360)
\fill[autres] (0,0) -- (352.8:3cm) arc (352.8:360:3cm) -- cycle;
\node[font=\footnotesize, text=popdark, align=center] at (356.4:4.2cm) {\textbf{Autres (fuites)}\\2\%};

% Légende manuelle
\begin{scope}[xshift=5.5cm, yshift=-1.5cm, font=\footnotesize]
    \fill[energie] (0,-1*0.55) rectangle (0.4,-1*0.55+0.35);
    \node[anchor=west] at (0.5,-1*0.55+0.175) {Énergie (production électricité, chaleur)};
    \fill[transport] (0,-2*0.55) rectangle (0.4,-2*0.55+0.35);
    \node[anchor=west] at (0.5,-2*0.55+0.175) {Transport (route, air, mer)};
    \fill[industrie] (0,-3*0.55) rectangle (0.4,-3*0.55+0.35);
    \node[anchor=west] at (0.5,-3*0.55+0.175) {Industrie (ciment, acier, chimie, énergie)};
    \fill[agriculture] (0,-4*0.55) rectangle (0.4,-4*0.55+0.35);
    \node[anchor=west] at (0.5,-4*0.55+0.175) {Agriculture (élevage, riz, engrais)};
    \fill[foret] (0,-5*0.55) rectangle (0.4,-5*0.55+0.35);
    \node[anchor=west] at (0.5,-5*0.55+0.175) {Forêt \& changement d'affectation des sols};
    \fill[batiments] (0,-6*0.55) rectangle (0.4,-6*0.55+0.35);
    \node[anchor=west] at (0.5,-6*0.55+0.175) {Bâtiments (chauffage, cuisson direct)};
    \fill[dechets] (0,-7*0.55) rectangle (0.4,-7*0.55+0.35);
    \node[anchor=west] at (0.5,-7*0.55+0.175) {Déchets (décharges, eaux usées)};
    \fill[autres] (0,-8*0.55) rectangle (0.4,-8*0.55+0.35);
    \node[anchor=west] at (0.5,-8*0.55+0.175) {Autres (fuites charbon, pétrole, gaz)};
\end{scope}
\end{tikzpicture}
\legende{Répartition sectorielle des émissions anthropiques mondiales de gaz à effet de serre (en $\ce{CO2}$-équivalent, année de référence ~2019, source GIEC AR6 WGIII simplifiée). Le secteur de l'énergie (production d'électricité et de chaleur) domine, suivi par l'industrie, le transport, puis l'agriculture et la forêt (UTCATF). Au Cameroun, la part de l'UTCATF (déforestation) et de l'agriculture est proportionnellement plus élevée que la moyenne mondiale.}
\end{popfigure}

\medskip
\noindent \textbf{Analyse pour le Cameroun :} Si la moyenne mondiale place l'énergie en tête, le profil camerounais est particulier :
\begin{itemize}
    \item \textbf{UTCATF (Forêt + Agriculture)} : souvent \textbf{> 50\,\%} des émissions nationales (déforestation massive + élevage extensif + riziculture + agriculture sur brûlis).
    \item \textbf{Énergie} : part croissante (centrales thermiques fioul/gaz, groupes électrogènes, transport routier vieillissant, bois-énergie pour la cuisson $\approx 70\,\%$ des ménages).
    \item \textbf{Déchets} : gestion quasi inexistante du méthane de décharge.
\end{itemize}
Cela oriente les priorités d'atténuation : \textbf{lutter contre la déforestation et moderniser l'agriculture/élevage} sont les leviers majeurs pour le Cameroun, en parallèle de la transition énergétique.

\begin{exoresolu}[Calcul d'émissions évitées par une action concrète]
\textbf{Énoncé.} Un lycée de Yaoundé remplace ses 4 groupes électrogènes diesel (chacun 50 kVA, fonctionnant 4 h/jour en moyenne 200 jours/an) par une installation solaire de 100 kWc avec batteries, couvrant 80\,\% des besoins. Le facteur d'émission du diesel est de $2,68\,\mathrm{kg\,CO2/L}$ ; la consommation spécifique est de $0,25\,\mathrm{L/kWh}$. Calculer la réduction annuelle d'émissions de $\ce{CO2}$ (en tonnes).

\textbf{Solution détaillée.}
\begin{enumerate}
    \item \textbf{Énergie produite annuellement par les groupes :}
    Puissance utile totale = $4 \times 50\,\mathrm{kVA} \times 0,8$ (facteur de puissance) $\approx 160\,\mathrm{kW}$.
    Énergie = $160\,\mathrm{kW} \times 4\,\mathrm{h/j} \times 200\,\mathrm{j} = 128\,000\,\mathrm{kWh/an}$.
    \item \textbf{Consommation de diesel :}
    $128\,000\,\mathrm{kWh} \times 0,25\,\mathrm{L/kWh} = 32\,000\,\mathrm{L/an}$.
    \item \textbf{Émissions $\ce{CO2}$ avant :}
    $32\,000\,\mathrm{L} \times 2,68\,\mathrm{kg/L} = 85\,760\,\mathrm{kg} \approx \textbf{85,8 tonnes $\ce{CO2}$/an}$.
    \item \textbf{Émissions après (20\,\% résiduel sur diesel) :}
    $0,20 \times 85,8 = \textbf{17,2 tonnes $\ce{CO2}$/an}$.
    \item \textbf{Réduction nette :}
    $85,8 - 17,2 = \textbf{68,6 tonnes $\ce{CO2}$ évitées par an}$.
\end{enumerate}
\textbf{Conclusion :} Ce projet solaire évite près de 69 tonnes de $\ce{CO2}$ par an, équivalent à la séquestration annuelle d'environ 3\,000 arbres matures. C'est une action d'atténuation concrète et mesurable.
\end{exoresolu}

\courssub{Synthèse : les leviers d'action identifiés}

\begin{aretenir}[Principales sources anthropiques de GES et leviers d'action au Cameroun]
\begin{center}
\begin{tabularx}{\linewidth}{|l|X|X|X|}
\hline
\rowcolor{popblueL}
\textbf{Secteur / Source} & \textbf{GES émis} & \textbf{Activités clés (Cameroun)} & \textbf{Leviers d'atténuation prioritaires} \\
\hline
Énergie fossile (transport, industrie, groupes, thermique) & $\ce{CO2}$ & Moto-taxis, cars, camions, centrales fioul/gaz, groupes électrogènes & Transport en commun (BRT, tram), véhicules électriques/hybrides, solaire/éolien/hydro, efficacité énergétique \\
\hline
Déforestation / Changement d'affectation des sols & $\ce{CO2}$ & Agriculture sur brûlis, bois-énergie, exploitation illégale, plantations agro-industrielles & Gestion durable des forêts (PEFC/FSC), reboisement, agroforesterie, foyers améliorés, biogaz domestique \\
\hline
Élevage (ruminants) & $\ce{CH4}$ & Élevage bovin extensif (Nord, Adamaoua) & Amélioration alimentation (supplémentation), gestion déjections (biogaz), sélection races productives \\
\hline
Riziculture irriguée & $\ce{CH4}$ & Plaines de la Bénoué, Logone, Maga, Yagoua & Riziculture non inondée (SRI), variétés à faible émission, alternance sec/humide \\
\hline
Engrais azotés / Déjections & $\ce{N2O}$ & Plantations coton, maïs, banane, palmier ; élevage intensif & Agriculture de précision, engrais à libération lente, rotation cultures, légumineuses \\
\hline
Déchets (décharges, eaux usées) & $\ce{CH4}$ & Décharges non contrôlées (Nkolfoulou, PK10, etc.), absence d'assainissement & Tri à la source, compostage, méthanisation (biogaz), captage biogaz décharges existantes \\
\hline
\end{tabularx}
\end{center}
\end{aretenir}

\medskip
\noindent Cette analyse sectorielle nous donne les clés pour comprendre \textbf{où et comment agir}. La section suivante examinera en détail les \textbf{conséquences} de ce réchauffement sur l'environnement et les populations camerounaises, avant d'aborder les stratégies d'\textbf{atténuation} (réduction des émissions) et d'\textbf{adaptation} (réduction de la vulnérabilité), ainsi que les cadres politiques (Accord de Paris, CDN du Cameroun, politiques sectorielles).

\coursec{Les conséquences du réchauffement climatique}

\noindent Dans la section précédente, nous avons vu comment les activités humaines — combustion d'énergies fossiles, déforestation, élevage intensif — augmentent la concentration des gaz à effet de serre (GES) et renforcent le piège thermique naturel. Ce déséquilibre radiatif positif se traduit par un accumul d'énergie dans le système climatique, dont la manifestation la plus directe est l'augmentation de la température moyenne de la planète. Mais le réchauffement ne se résume pas à une simple hausse du thermomètre : il déclenche une cascade de perturbations physiques, biologiques et humaines que nous allons analyser rigoureusement, en prenant appui sur des observations concrètes au Cameroun et dans le monde.

\courssub{Élévation de la température moyenne et impacts physiques directs}

\begin{definition}[Anomalie de température globale]
L'\cle{anomalie de température} est l'écart entre la température moyenne observée sur une période donnée (mois, année) et une valeur de référence calculée sur une période de référence (souvent 1951--1980 ou 1850--1900 pour l'ère préindustrielle). Elle permet de suivre l'évolution du climat en s'affranchissant des variations spatiales absolues.
\end{definition}

\noindent \textbf{Observation.} Depuis le début de l'ère industrielle (vers 1850--1900), la température moyenne à la surface du globe a augmenté d'environ \textbf{+1,1 à +1,2 °C} (GIEC, 6\textsuperscript{e} rapport, 2021). Les dernières décennies sont les plus chaudes jamais enregistrées : les années 2015--2023 occupent les premières places. Au Cameroun, les stations météorologiques de Yaoundé, Douala, Garoua et Maroua confirment cette tendance : on observe une hausse moyenne de \SI{0,5}{\celsius} à \SI{1,0}{\celsius} depuis les années 1970, avec une accélération nette depuis les années 1990.

\textbf{Hypothèse.} Cette augmentation de la température moyenne entraîne deux phénomènes physiques majeurs qui font monter le niveau des mers : (1) la \emph{dilatation thermique} de l'eau de mer (l'eau chaude occupe un volume plus grand) et (2) la \emph{fonte des glaces continentales} (glaciers de montagne, calottes du Groenland et de l'Antarctique) qui transfère de l'eau stockée sur les continents vers les océans.

\textbf{Données et interprétation.} Les mesures satellitaires (altimétrie radar, missions Jason, Sentinel-6) et les marégraphes côtiers montrent une élévation moyenne du niveau des mers de \textbf{\SI{3,3}{\milli\meter} par an} sur la période 1993--2023, avec une accélération récente à \textbf{\SI{4,5}{\milli\meter}/\text{an}}. La dilatation thermique contribue pour environ 40 \% et la fonte des glaces pour 60 \% (proportion croissante).

\begin{popfigure}
\begin{tikzpicture}[
    node distance=1.2cm and 1.5cm,
    box/.style={draw=popdark, thick, rounded corners, fill=popcream, minimum width=2.8cm, minimum height=1.1cm, align=center, font=\small},
    arrow/.style={-Stealth, thick, popdark},
    plus/.style={font=\large, poporange}
]
\node[box, align=center] (emis){Émissions anthropiques\\(\ce{CO2}, \ce{CH4}, \ce{N2O}, etc.)};
\node[box, right=of emis, align=center] (ges){Renforcement de l'effet\\de serre (forçage radiatif +)};
\node[box, right=of ges, align=center] (warm){Réchauffement global\\(air + océan)};
\node[box, below=of warm, xshift=-1.5cm, align=center] (melt){Fonte des glaces\\continentales};
\node[box, below=of warm, xshift=1.5cm, align=center] (therm){Dilatation thermique\\de l'océan};
\node[box, below=of melt, yshift=-0.5cm, align=center] (sealevel){Élévation du niveau\\des mers (+3,3 mm/an)};
\draw[arrow] (emis) -- (ges);
\draw[arrow] (ges) -- (warm);
\draw[arrow] (warm) -- (melt);
\draw[arrow] (warm) -- (therm);
\draw[arrow] (melt) -- (sealevel);
\draw[arrow] (therm) -- (sealevel);
\node[plus, rotate=90] at ($(melt)!0.5!(sealevel) + (-0.8,0)$) {+};
\node[plus, rotate=90] at ($(therm)!0.5!(sealevel) + (0.8,0)$) {+};
\end{tikzpicture}
\legende{Chaîne de causalité : des émissions de GES à l'élévation du niveau marin. Deux mécanismes physiques s'additionnent : la fonte des glaces continentales (apport de masse) et la dilatation thermique (augmentation de volume).}
\end{popfigure}

\begin{exemplebox}[Vitesse actuelle de la montée des eaux]
Le niveau moyen des mers s'élève d'environ \textbf{3 à 4 mm par an} actuellement, un rythme qui \emph{s'accélère} (environ \SI{1,4}{\milli\meter}/\text{an} dans les années 1990, \SI{3,6}{\milli\meter}/\text{an} sur 2006--2018, \SI{4,5}{\milli\meter}/\text{an} récemment). Sur le siècle, cela représente déjà \SI{20}{\centi\meter} de hausse depuis 1900. Pour le littoral camerounais (Kribi, Douala, Limbé, Tiko), cela signifie une érosion côtière accélérée, une salinisation des nappes phréatiques et une submersion plus fréquente lors des marées hautes et tempêtes. À Douala, la zone du port et les quartiers bas (Bessengue, Akwa) sont déjà menacés lors des crues du Wouri couplées aux marées de vives-eaux.
\end{exemplebox}

\courssub{Modification des régimes pluviométriques et hydrologie : sécheresses, inondations et le cas du lac Tchad}

\noindent \textbf{Observation.} Le réchauffement intensifie le cycle de l'eau : une atmosphère plus chaude contient plus de vapeur d'eau (loi de Clausius-Clapeyron : \SI{7}{\%} d'humidité en plus par °C). Cela se traduit par une \textbf{redistribution spatiale et temporelle des précipitations} : les zones humides deviennent plus humides (pluies extrêmes), les zones sèches plus sèches (sécheresses prolongées), et les événements extrêmes se multiplient.

Au Cameroun, ce contraste est saisissant :
\begin{itemize}
    \item \textbf{Extrême-Nord (région soudano-sahélienne) :} pluviométrie moyenne en baisse, saison des pluies plus courte et plus erratique, poches de sécheresse sévères (ex. 2010, 2012, 2021). Les cours d'eau temporaires (mayos) s'assèchent plus tôt.
    \item \textbf{Grandes métropoles (Douala, Yaoundé) :} épisodes de pluies diluviennes de plus en plus fréquents (ex. \SI{150}{\milli\meter} en 24 h à Douala en août 2023), provoquant inondations soudaines, glissements de terrain et pertes humaines, amplifiés par l'urbanisation anarchique et l'imperméabilisation des sols.
\end{itemize}

\textbf{Données : pluviométrie dans l'Extrême-Nord (données simplifiées, période 1950--2020).}

\begin{popfigure}
\begin{tikzpicture}
\begin{axis}[
    width=\linewidth, height=7cm,
    xlabel={Année}, ylabel={Précipitations annuelles (mm)},
    xmin=1950, xmax=2020, ymin=400, ymax=1100,
    xtick={1950,1970,1990,2010,2020}, ytick={500,700,900,1100},
    grid=major, grid style={popline}, legend pos=north west,
    title style={font=\small, align=center}, title={Évolution simplifiée des précipitations moyennes (Extrême-Nord Cameroun)}
]
\addplot[popcurve, mark=*, mark size=1.5, popblue, only marks] coordinates {
    (1950,950) (1955,880) (1960,920) (1965,850) (1970,780) (1975,650) (1980,620) (1985,580) (1990,680) (1995,720) (2000,600) (2005,750) (2010,550) (2015,680) (2020,620)
};
\addlegendentry{Données annuelles}
\addplot[poporange, thick, domain=1950:2020, samples=2] { -1.2*x + 3290 }; % tendance approximative
\addlegendentry{Tendance linéaire (baisse)}
\end{axis}
\end{tikzpicture}
\legende{Évolution simplifiée des précipitations annuelles moyennes dans l'Extrême-Nord (données synthétiques représentatives). On observe une variabilité interannuelle forte et une tendance à la baisse sur le long terme, avec des années très sèches (années 1970--1980, 2010) et des récupérations partielles.}
\end{popfigure}

\noindent \textbf{Interprétation.} La courbe montre une variabilité naturelle forte (oscillations décennales liées à la mousson ouest-africaine et aux températures de surface de l'Atlantique), mais la tendance de fond est à la baisse et aux extrêmes. Cela impacte directement l'agriculture pluviale (mil, sorgho, arachide) et la disponibilité en eau pour les populations et le bétail.

\begin{experience}[Le lac Tchad : un baromètre régional du changement climatique]
\textbf{Matériel :} Cartes satellitaires (Landsat, Sentinel) du lac Tchad sur plusieurs décennies (1963, 1973, 1987, 2000, 2020), données pluviométriques du bassin, prélèvements d'eau pour irrigation.
\textbf{Protocole :} Mesurer la superficie de l'eau libre sur chaque image ; corréler avec l'indice pluviométrique du bassin et les volumes prélevés pour l'agriculture (projet d'irrigation de la South Chad Irrigation Scheme, barrages en amont sur le Chari/Logone).
\textbf{Observations :}
\begin{itemize}
    \item 1963 : \textasciitilde 25 000 km² (lac « Mega-Chad » résiduel).
    \item 1973 : \textasciitilde 17 000 km² (début grande sécheresse sahélienne).
    \item 1984--1985 : minimum historique \textasciitilde 2 000--2 500 km² (lac scindé en deux bassins nord/sud).
    \item Années 1990--2000 : légère récupération (bassin sud \textasciitilde 5 000--10 000 km²), mais instable.
    \item 2020 : \textasciitilde 1 500--2 500 km² selon la saison.
\end{itemize}
\textbf{Interprétation :} La réduction de \textbf{\textasciitilde 90 \% de la superficie depuis les années 1960} résulte de la \textbf{conjonction de deux facteurs} : (1) \emph{climatique} : baisse durable des pluies sur le bassin versant (sécheresses des années 1970--1990) et augmentation de l'évapotranspiration (chaleur) ; (2) \emph{anthropique} : prélèvements massifs pour l'irrigation en amont (barrages de Maga, Yagoua, projets nigérians) qui réduisent l'apport du Chari-Logone (90 \% de l'alimentation du lac).
\textbf{Conclusion :} Le lac Tchad illustre parfaitement l'interaction \emph{changement climatique + pression humaine} : le réchauffement aggrave la raréfaction de l'eau, ce qui pousse à plus de prélèvements, accélérant l'assèchement. Conséquences : effondrement de la pêche, perte de pâturages, migrations de populations (pêcheurs, éleveurs), conflits d'usage de l'eau, insécurité alimentaire pour \textasciitilde 30 millions de personnes dans le bassin (Cameroun, Tchad, Niger, Nigeria).
\end{experience}

\begin{savaistu}[Le lac Tchad, géant en péril]
Le lac Tchad est passé d'environ \textbf{25 000 km² dans les années 1960} à moins de \textbf{2 500 km² certaines années récentes}, affectant des millions de personnes au Cameroun (région de l'Extrême-Nord), au Tchad, au Niger et au Nigeria. C'est l'un des exemples les plus dramatiques de la vulnérabilité des ressources en eau face au couplage « réchauffement + surexploitation ». La Commission du Bassin du Lac Tchad (CBLT) étudie le projet de transfert d'eau du bassin du Congo (Oubangui) pour le recharger — un projet pharaonique aux enjeux géopolitiques, écologiques et financiers immenses.
\end{savaistu}

\courssub{Perturbation des écosystèmes et de la biodiversité}

\noindent \textbf{Mécanisme.} Le réchauffement modifie les conditions abiotiques (température, eau, pH des océans, saisonnalité) plus vite que de nombreuses espèces ne peuvent s'adapter (migration, adaptation génétique, plasticité phénotypique). Trois réponses principales sont observées :
\begin{enumerate}
    \item \textbf{Déplacement des aires de distribution :} vers les pôles ou en altitude (ex. : espèces de montagne du Mont Cameroun contraintes de remonter, perte d'habitat au sommet).
    \item \textbf{Modification des phénologies :} floraison plus précoce, migration des oiseaux décalée, décalage proies-prédateurs (ex. : pic de chenilles par rapport à la nidification des mésanges).
    \item \textbf{Extinction locale ou globale :} pour les espèces endémiques à aire restreinte, incapables de migrer (ex. : certains amphibiens des montagnes camerounaises, coraux).
\end{enumerate}

\noindent \textbf{Cas emblématique : le blanchissement des coraux.} L'océan absorbe \textasciitilde 90 \% de l'excès de chaleur et \textasciitilde 25--30 \% du \ce{CO2} anthropique (acidification). Une anomalie de température de \SI{1}{\celsius} à \SI{2}{\celsius} au-dessus du maximum saisonnier pendant quelques semaines suffit à provoquer l'expulsion des zooxanthelles (algues symbiotiques) par le corail : le corail blanchit, meurt si le stress persiste. Au Cameroun, les récifs coralliens du golfe de Guinée (zone de Kribi, Limbé) sont menacés, avec des épisodes de blanchissement observés en 1998, 2010, 2016, 2020. Cela déstructure l'habitat de nombreuses espèces de poissons d'intérêt halieutique.

\begin{propriete}[Déséquilibre des chaînes alimentaires]
Le décalage phénologique (ex. : floraison avant pollinisateurs, éclosion insectes avant arrivée oiseaux insectivores) rompt les interactions trophiques synchronisées par la sélection naturelle. Cela peut provoquer des effondrements de populations en cascade (ex. : baisse des populations d'oiseaux migrateurs au Sahel liée à la désynchronisation avec les pics d'insectes à la reproduction).
\end{propriete}

\courssub{Conséquences sanitaires, agricoles, alimentaires et migrations climatiques}

\noindent \textbf{1. Santé humaine.}
\begin{itemize}
    \item \textbf{Extension des zones de paludisme :} l'anophèle (vecteur) et le plasmodium (parasite) voient leur développement accéléré par la chaleur. Les zones d'altitude (Ouest, Adamaoua, Mont Cameroun), autrefois trop fraîches, deviennent propices. On observe une remontée du paludisme vers \SI{1500}{\meter}--\SI{2000}{\meter} d'altitude.
    \item \textbf{Autres maladies tropicales :} dengue, chikungunya (vecteur \emph{Aedes aegypti}/albopictus), choléra (lié aux inondations et chaleur), méningite (saison sèche prolongée, poussières).
    \item \textbf{Canicules :} stress thermique direct (coups de chaleur, déshydratation, aggravation cardiovasculaire/répiratoire), surtout en ville (îlot de chaleur urbain : Douala, Yaoundé peuvent être \SI{3}{\celsius}--\SI{5}{\celsius} plus chaudes que la campagne la nuit).
\end{itemize}

\noindent \textbf{2. Agriculture et sécurité alimentaire.}
\begin{itemize}
    \item \textbf{Baisse des rendements :} stress hydrique (sécheresse), stress thermique (stérilité du pollen au-dessus de \SI{35}{\celsius} pour le maïs, le mil), raccourcissement du cycle cultural. Modélisations : -10 à -20 \% pour le mil/sorgho au Sahel d'ici 2050 ; -10 à -30 \% pour le cacao au Sud-Cameroun (sensible à l'excès d'eau et aux parasites favorisés par l'humidité).
    \item \textbf{Insécurité alimentaire :} hausse des prix, dépendance aux importations (blé, riz), malnutrition chronique (Extrême-Nord, Nord).
    \item \textbf{Conflits agriculteurs-éleveurs :} raréfaction des pâturages et de l'eau (Extrême-Nord, Adamaoua, Nord) pousse les éleveurs transhumants vers le sud plus tôt et plus loin, empiétant sur les champs (mil, maïs, coton). Tensions violentes récurrentes.
\end{itemize}

\noindent \textbf{3. Migrations climatiques et déplacements de populations.}
Le GIEC et l'OIM estiment que le changement climatique pourrait déplacer \textbf{200 millions de personnes d'ici 2050} (migrations internes surtout). Au Cameroun :
\begin{itemize}
    \item \textbf{Extrême-Nord :} populations riveraines du lac Tchad et des plaines d'inondation (Logone, Chari) contraintes de fuir l'assèchement, l'insécurité (Boko Haram exacerbé par la pauvreté) → déplacés internes vers Maroua, Mora, ou vers le Sud (Yaoundé, Douala).
    \item \textbf{Littoral :} érosion côtière (Kribi, Limbé, Douala) + inondations → relogement forcé de quartiers précaires.
    \item \textbf{Zones de montagne :} perte de productivité agricole (sols lessivés, parasites) → exode rural des jeunes.
\end{itemize}

\begin{attention}[Erreur fréquente : météo vs climat]
Ne pas attribuer \textbf{chaque événement météo isolé} (une tempête, une vague de chaleur, une inondation) \emph{uniquement} au réchauffement climatique. La météo est l'état de l'atmosphère à un instant et un lieu donnés (échelle jour/semaine) ; le climat est la \textbf{moyenne statistique sur au moins 30 ans} (normales 1991--2020). Le réchauffement \emph{modifie les probabilités} : il rend les canicules plus fréquentes, plus intenses, plus longues ; il rend les pluies extrêmes plus probables. On parle d'\emph{attribution probabiliste} (ex. : « cette canicule 2023 était 10 fois plus probable à cause du réchauffement »), pas de causalité unique pour un événement unique. Au Probatoire, distinguez toujours \textbf{météo} (court terme, chaotique) et \textbf{climat} (long terme, statistique).
\end{attention}

\begin{methode}[Analyser une conséquence du réchauffement climatique]
Pour répondre à une question du type « Expliquez comment le renforcement de l'effet de serre entraîne telle conséquence », suivez cette chaîne logique en 4 étapes :
\begin{enumerate}
    \item \textbf{Cause racine :} Identifiez l'augmentation des GES (\ce{CO2}, \ce{CH4}...) et le forçage radiatif positif.
    \item \textbf{Mécanisme physique global :} Réchauffement de l'atmosphère et de l'océan (accumulation d'énergie).
    \item \textbf{Mécanisme physique spécifique :} Reliez le réchauffement au processus direct (ex. : dilatation thermique + fonte glaces → montée eaux ; évaporation accrue + circulation atmosphérique modifiée → sécheresses/inondations ; température > seuil thermique → blanchissement coraux).
    \item \textbf{Impact humain/environnemental :} Déclinez la conséquence concrète (submersion côtière, échec agricole, extension paludisme, migration) en citant un exemple local (Cameroun) si possible.
\end{enumerate}
\textbf{Exemple :} « Montée des eaux à Douala » → 1. +GES → 2. +Température globale → 3. Dilatation océan + Fonte glaces → Niveau mer +3,3 mm/an → 4. Submersion quartiers bas + érosion + salinisation nappes.
\end{methode}

\begin{aretenir}[Météo vs Climat : distinction exigible au Probatoire]
\begin{itemize}
    \item \textbf{Météo} : état instantané de l'atmosphère (température, vent, pluie, pression) à un lieu et un moment donnés. Prévisible sur quelques jours (modèles numériques). Ex. : « Il pleut à Douala ce mardi 14 mai, 28 °C. »
    \item \textbf{Climat} : synthèse statistique des conditions météorologiques sur une longue période (\textbf{minimum 30 ans} selon l'OMM). Caractérisé par des moyennes, des extrêmes, des fréquences. Ex. : « À Douala, le mois de mai reçoit en moyenne 250 mm de pluie sur la période 1991--2020. »
    \item \textbf{Variabilité climatique} : fluctuations naturelles autour de la moyenne (El Niño/La Niña, oscillations décennales).
    \item \textbf{Changement climatique} : modification statistiquement significative de la moyenne et/ou de la variabilité du climat, persistant sur plusieurs décennies, d'origine naturelle ou anthropique (actuelle : anthropique).
\end{itemize}
\textbf{Piège Probatoire :} On ne dit pas « le climat d'aujourd'hui est pluvieux » (c'est la météo), on dit « le climat de cette région est de type équatorial de transition ». On ne conclut pas au réchauffement sur un seul été chaud, mais sur la tendance sur 30 ans.
\end{aretenir}

\begin{exoresolu}[Analyse d'un article de presse sur les inondations à Douala]
\textbf{Énoncé.} Un article titre : « Inondations catastrophiques à Douala : le changement climatique en cause ». L'article rapporte \SI{180}{\milli\meter} de pluie en 6 heures, des quartiers entiers sous l'eau, 12 morts. La mairie évoque « des pluies sans précédent dues au réchauffement ». En vous basant sur vos connaissances, critiquez l'argumentaire et complétez-le.
\tcblower
\textbf{Solution détaillée.}
\begin{enumerate}
    \item \textbf{Critique :} L'article commet l'erreur classique d'attribution unique (voir \textbf{Attention} ci-dessus). Un épisode pluvieux isolé, aussi extrême soit-il, relève de la \textbf{météo}. On ne peut pas dire « c'est \emph{la} cause » mais « le réchauffement a \emph{augmenté la probabilité et l'intensité} d'un tel événement ».
    \item \textbf{Complément scientifique (chaîne causale) :}
    \begin{itemize}
        \item +1,2 °C global → atmosphère plus chaude → capacité de rétention d'eau +7 \% par °C (Clausius-Clapeyron) → potentiel de pluies plus intenses.
        \item Réchauffement de l'Atlantique tropical → plus d'évaporation → plus d'humidité advectée vers le golfe de Guinée par la mousson.
        \item \textbf{Facteurs locaux aggravants (non climatiques) :} urbanisation galopante (imperméabilisation \textasciitilde 70 \% des sols), occupation des zones inondables (marécages comblés), absence de drainage adapté, déforestation des collines environnant Douala (ruissellement accéléré).
    \end{itemize}
    \item \textbf{Conclusion nuancée :} Le réchauffement climatique \emph{charge les dés} pour des pluies plus extrêmes, mais la \emph{vulnérabilité} du territoire (aménagement, pauvreté, gouvernance) transforme l'aléa en catastrophe. L'adaptation (drainage, planification urbaine, systèmes d'alerte) est aussi urgente que l'atténuation.
\end{enumerate}
\end{exoresolu}

\begin{popfigure}
\begin{tikzpicture}[
    scale=0.85,
    region/.style={draw=popdark, thick, fill=popcream, rounded corners, minimum width=2.2cm, minimum height=1.2cm, align=center, font=\small},
    impact/.style={draw=poporange, thick, fill=poporange!20, rounded corners, minimum width=2.5cm, minimum height=1cm, align=center, font=\small\bfseries},
    arrow/.style={-Stealth, thick, popdark},
    coast/.style={draw=popblue, very thick},
    river/.style={draw=popblue, thick},
]
% Carte schématique Cameroun
\draw[coast] (0,0) .. controls (1,0.5) and (2,1) .. (3,1.5) .. controls (4,2) and (5,2.5) .. (6,3) .. controls (7,3.5) and (8,3) .. (9,2) .. controls (10,1) and (11,0) .. (11,-0.5);
\draw[river] (4,1.5) -- (4,4) -- (3.5,5) -- (3,6); % Wouri
\draw[river] (7,2) -- (7,4.5) -- (6.5,5.5) -- (6,6.5); % Sanaga approx
\draw[river] (10.5,0) -- (10.5,3) -- (10,4) -- (9.5,5); % Logone/Chari vers Nord
% Villes
\node[region] (douala) at (3.5,1.8) {Douala};
\node[region] (kribi) at (1.5,1.2) {Kribi};
\node[region] (limbe) at (5.5,2.2) {Limbé};
\node[region] (yaounde) at (4.5,3.5) {Yaoundé};
\node[region] (maroua) at (9.5,5.5) {Maroua};
\node[region] (garoua) at (8.5,4.5) {Garoua};
% Impacts
\node[impact, anchor=north, align=center] (flood) at (3.5,0.5){Inondations\\récurrentes};
\node[impact, anchor=north, align=center] (coastal) at (1.5,0){Montée des eaux\\+ érosion};
\node[impact, anchor=north, align=center] (coastal2) at (5.5,1){Montée des eaux\\+ érosion};
\node[impact, anchor=south, align=center] (drought) at (9.5,6.8){Sécheresse\\prolongée};
\node[impact, anchor=south, align=center] (lake) at (10.5,6.5){Lac Tchad\\-90\%};
\node[impact, anchor=west, align=center] (conflict) at (7,5.5){Conflits\\agri/élev.};
\node[impact, anchor=east, align=center] (health) at (4.5,4.8){Paludisme\\altitude};
\draw[arrow, poporange] (douala.south) -- (flood.north);
\draw[arrow, poporange] (kribi.south) -- (coastal.north);
\draw[arrow, poporange] (limbe.south) -- (coastal2.north);
\draw[arrow, poporange] (maroua.north) -- (drought.south);
\draw[arrow, poporange] (garoua.north) -- (lake.south);
\draw[arrow, poporange] (yaounde.east) -- (health.west);
\draw[arrow, poporange] (yaounde.north east) -- (conflict.south west);
% Légende intégrée
\node[draw=popdark, fill=white, rounded corners, font=\footnotesize, align=left, anchor=south west] at (0.2,-0.5) {
    \textbf{Impacts climatiques observés au Cameroun}\\
    \textcolor{poporange}{■} Côtes : montée eaux, érosion, inondations\\
    \textcolor{poporange}{■} Métropoles : inondations pluviales extrêmes\\
    \textcolor{poporange}{■} Montagnes : paludisme en altitude, érosion biodiversité\\
    \textcolor{poporange}{■} Nord/Extrême-Nord : sécheresse, lac Tchad, conflits
};
\end{tikzpicture}
\legende{Carte schématique du Cameroun localisant les principaux impacts du réchauffement climatique observés ou projetés. Les flèches oranges relient les zones urbaines/régionales aux impacts spécifiques documentés.}
\end{popfigure}

\noindent \textbf{Synthèse.} Le réchauffement climatique n'est pas une menace future : il est \textbf{actuel, mesurable, et déjà coûteux} pour le Cameroun. Ses conséquences s'enchaînent logiquement depuis l'accumulation d'énergie (effet de serre renforcé) jusqu'aux impacts sanitaires, agricoles, économiques et migratoires, en passant par les perturbations physiques (eau, glace, cycle hydrologique) et biologiques. La compréhension de cette chaîne causale — du global au local — est la clé pour proposer des actions d'\textbf{atténuation} (réduire les émissions) et d'\textbf{adaptation} (réduire la vulnérabilité), qui feront l'objet de la prochaine section.

\coursec{Lutter contre l'effet de serre : atténuation et adaptation}

Après avoir analysé les conséquences déjà visibles du réchauffement climatique — élévation du niveau de la mer menaçant nos côtes du Golfe de Guinée, perturbation des saisons agricoles au Sahel, multiplication des épisodes caniculaires à Yaoundé et Douala — une question s'impose avec urgence : \textbf{que pouvons-nous faire, concrètement, ici et maintenant ?} La réponse scientifique et politique structure l'action climatique en deux piliers complémentaires : l'atténuation, qui s'attaque aux \emph{causes} (les émissions de gaz à effet de serre), et l'adaptation, qui gère les \emph{effets} déjà inévitables. Comprendre cette distinction, c'est déjà agir en citoyen éclairé.

\courssub{Deux stratégies complémentaires : atténuation et adaptation}

Imaginons une baignoire qui déborde : l'eau qui coule du robinet représente nos émissions de gaz à effet de serre (GES), le niveau d'eau représente la concentration atmosphérique, et l'inondation du sol représente le réchauffement et ses dégâts. Deux gestes s'imposent : \textbf{fermer le robinet} (atténuation) et \textbf{protéger le sol avec des sacs de sable ou une pompe} (adaptation). Fermer le robinet sans protéger le sol laisse l'eau déjà présente faire des dégâts ; protéger le sol sans fermer le robinet condamne à une lutte sans fin contre un débit croissant. Les deux sont indispensables.

\begin{definition}[Atténuation (Mitigation)]
L'\cle{atténuation} désigne l'ensemble des actions humaines visant à \textbf{réduire les émissions de gaz à effet de serre} (en agissant sur les sources : énergie, transport, industrie, agriculture, déforestation) \textbf{ou à augmenter les puits de carbone} (forêts, océans, sols, technologies de capture) pour limiter l'ampleur du réchauffement climatique futur.
\end{definition}

\begin{definition}[Adaptation]
L'\cle{adaptation} désigne l'ensemble des mesures (techniques, institutionnelles, comportementales) visant à \textbf{réduire la vulnérabilité} des sociétés, des économies et des écosystèmes aux effets actuels et futurs du réchauffement climatique, et à \textbf{renforcer leur résilience} (capacité à absorber un choc et à se rétablir).
\end{definition}

\begin{aretenir}[Distinction fondamentale]
\begin{itemize}
    \item \textbf{Atténuation} = agir sur la \emph{cause} (émissions/puits) $\rightarrow$ bénéfice \emph{global} et \emph{long terme} (bénéficie à toute la planète, effets différés).
    \item \textbf{Adaptation} = agir sur la \emph{conséquence} (vulnérabilité) $\rightarrow$ bénéfice \emph{local} et \emph{court/moyen terme} (protège une communauté, un écosystème précis, effets immédiats).
    \item \textbf{Synergie} : une même action peut servir les deux (ex. : reboisement = puits de carbone \textbf{ET} protection contre l'érosion/les crues).
\end{itemize}
\end{aretenir}

\begin{popfigure}
\begin{tikzpicture}[
    node distance=1.2cm and 2cm,
    box/.style={draw=popdark, thick, rounded corners, fill=popcream, minimum width=3.5cm, minimum height=1.2cm, align=center, font=\small},
    titlebox/.style={draw=popblue, thick, rounded corners, fill=popblueL, minimum width=4cm, minimum height=1cm, align=center, font=\bfseries\small},
    arrow/.style={-Stealth, thick, popdark},
    brace/.style={decorate, decoration={brace, amplitude=5pt}, thick, popblue}
]
% Titre central
\node[titlebox] (titre) {Stratégies face au changement climatique};

% Colonne Atténuation
\node[box, below left=of titre, xshift=-1cm, align=center] (att1){Énergies renouvelables\\ (hydro, solaire, biomasse)};
\node[box, below=of att1, align=center] (att2){Efficacité énergétique\\ (bâtiments, industrie)};
\node[box, below=of att2, align=center] (att3){Foresterie : reboisement,\\ agroforesterie, zéro déforestation};
\node[box, below=of att3, align=center] (att4){Transports durables\\ (TC, vélo, véhicules propres)};
\node[box, below=of att4, align=center] (att5){Gestion déchets/agriculture\\ (compost, méthanisation, engrais)};

% Colonne Adaptation
\node[box, below right=of titre, xshift=1cm, align=center] (adap1){Protection côtière\\ (digues, mangroves, relocalisation)};
\node[box, below=of adap1, align=center] (adap2){Agriculture résiliente\\ (variétés sècheresse, agroécologie)};
\node[box, below=of adap2, align=center] (adap3){Gestion de l'eau\\ (retenues, irrigation goutte-à-goutte)};
\node[box, below=of adap3, align=center] (adap4){Systèmes d'alerte précoce\\ (inondations, canicules, santé)};
\node[box, below=of adap4, align=center] (adap5){Urbanisme adapté\\ (îlots de fraîcheur, drainage)};

% Flèches et accolades
\draw[brace] ($(att1.north west) + (-0.3,0.2)$) -- ($(att5.south west) + (-0.3,-0.2)$) node[midway, left=8pt, rotate=90, font=\bfseries\small, popblue] {ATTÉNUATION (Causes)};
\draw[brace] ($(adap1.north east) + (0.3,0.2)$) -- ($(adap5.south east) + (0.3,-0.2)$) node[midway, right=8pt, rotate=-90, font=\bfseries\small, poporange] {ADAPTATION (Effets)};

% Flèches de synergie
\draw[arrow, popgreen, dashed] (att3.east) -- node[midway, above, sloped, font=\tiny, popgreen] {Synergie} (adap2.west);
\draw[arrow, popgreen, dashed] (att5.east) -- node[midway, above, sloped, font=\tiny, popgreen] {Synergie} (adap3.west);
\end{tikzpicture}
\legende{Tableau comparatif : actions d'atténuation (réduire les causes) vs actions d'adaptation (gérer les effets). Les flèches vertes pointillées illustrent les synergies (ex. : reboisement et agriculture résiliente).}
\end{popfigure}

\courssub{L'atténuation : réduire les émissions à la source (fermer le robinet)}

Au Cameroun, notre profil d'émissions est particulier : le secteur \textbf{Foresterie et Utilisation des Terres (FUT)} domine (déforestation, dégradation des forêts du bassin du Congo), suivi par l'\textbf{Agriculture} (élevage, riziculture, brûlis) et l'\textbf{Énergie} (transport, production électrique thermique, biomasse traditionnelle). L'atténuation doit cibler ces secteurs prioritaires.

\paragraph{1. Transition énergétique : l'atout hydroélectrique et le potentiel solaire.}
Le Cameroun dispose d'un potentiel hydroélectrique énorme (2\textsuperscript{e} d'Afrique après la RDC, $\approx$ 23 GW exploitables). Les barrages de \textbf{Song Loulou, Edea, Memve'ele, Nachtigal} fournissent une électricité \textbf{bas-carbone} (faibles émissions par kWh). Pourtant, le mix électrique intègre encore du thermique (fioul, gaz) pour combler les pointes et les sécheresses.
\begin{itemize}
    \item \textbf{Action :} Accélérer les grands barrages (Lom Pangar pour la régulation, Grand Eweng) et le raccordement des réseaux Nord/Sud.
    \item \textbf{Action Nord Cameroun :} Déployer massivement le \textbf{solaire photovoltaïque} (irradiation \textgreater 2000 kWh/m²/an) pour l'électrification rurale, le pompage agricole et le remplacement des groupes électrogènes diesel.
    \item \textbf{Efficacité énergétique :} Isoler les bâtiments (peintures réfléchissantes, ventilation naturelle), promouvoir les ampoules LED, lutter contre les pertes sur le réseau (vétusté, fraude).
\end{itemize}

\paragraph{2. Forêts : le poumon du bassin du Congo, notre levier majeur.}
La forêt dense humide du Sud-Cameroun (partie du bassin du Congo, 2\textsuperscript{e} massif forestier tropical mondial) stocke des milliards de tonnes de carbone. Sa déforestation (agriculture sur brûlis, exploitation forestière illégale, extension cacaoyère, infrastructure) libère ce stock.
\begin{itemize}
    \item \textbf{Reboisement / Restauration :} Planter des essences locales à croissance rapide (\emph{Terminalia superba} - Fraké, \emph{Triplochiton scleroxylon} - Ayous) sur les terres dégradées. Objectif national : 12 millions d'hectares restaurés à l'horizon 2030 (Initiative AFR100).
    \item \textbf{Agroforesterie :} Associer arbres et cultures (cacaoyers sous ombrage, caféiers, plantains + arbres d'engrais comme \emph{Acacia}, \emph{Leucaena}). Cela séquestre du carbone \textbf{ET} améliore les rendements (ombre, azote, matière organique).
    \item \textbf{Lutte contre la déforestation :} Certification FSC/PEFC des concessions forestières, surveillance par satellite (Global Forest Watch), sécurisation des droits fonciers des communautés locales (gestion communautaire durable).
\end{itemize}

\begin{exemplebox}[Le cycle vertueux du reboisement au Cameroun]
Un paysan de la région du Centre plante 1 hectare d'arbres (espacement 3m x 3m $\approx$ 1100 plants). Après 10 ans, la jeune forêt séquestre environ \textbf{5 à 10 tonnes de CO\textsubscript{2} par an et par hectare}. En 20 ans, c'est 100--200 t CO\textsubscript{2} stockés. Si ce paysan vend des \textbf{crédits carbone} (marché volontaire ou mécanisme REDD+), il tire un revenu supplémentaire qui finance l'entretien de la plantation et l'école de ses enfants. La forêt protège aussi le bassin versant qui alimente son village en eau potable. \emph{Atténuation + Adaptation + Développement local.}
\end{exemplebox}

\begin{popfigure}
\begin{tikzpicture}[
    node distance=1.5cm and 1cm,
    proc/.style={draw=popgreen, thick, rounded corners, fill=popgreenL, minimum width=3cm, minimum height=1cm, align=center, font=\small},
    arrow/.style={-Stealth, thick, popgreen},
    startstop/.style={draw=popblue, thick, ellipse, fill=popblueL, minimum width=3cm, minimum height=1cm, align=center, font=\bfseries\small},
]
\node[startstop] (start) {Terre dégradée / Déforestée};
\node[proc, below=of start, align=center] (plant){Plantation d'essences locales\\(Reboisement / Agroforesterie)};
\node[proc, below=of plant, align=center] (croiss){Croissance : Photosynthèse intense\\\ce{6CO2 + 6H2O -> C6H12O6 + 6O2}};
\node[proc, below=of croiss, align=center] (stock){Stockage Carbone\\Biomasse (tronc, racines) + Sol (humus)};
\node[proc, below=of stock, align=center] (effet){Diminution [CO\textsubscript{2}] atmosphérique\\\textrightarrow Limitation réchauffement global};
\node[proc, right=of croiss, xshift=1.5cm, fill=poporangeL, draw=poporange, align=center] (co){Co-bénéfices locaux :\textbullet\ Protection sols/eau\\ \textbullet\ Biodiversité\\ \textbullet\ Revenus (bois, fruits, carbone)\\ \textbullet\ Microclimat frais};
\draw[arrow] (start) -- (plant);
\draw[arrow] (plant) -- (croiss);
\draw[arrow] (croiss) -- (stock);
\draw[arrow] (stock) -- (effet);
\draw[arrow, dashed, poporange] (croiss) -- (co);
\end{tikzpicture}
\legende{Cycle vertueux : le reboisement (atténuation) déclenche une cascade d'effets bénéfiques pour le climat global (stockage CO\textsubscript{2}) et l'environnement local (adaptation, revenus).}
\end{popfigure}

\paragraph{3. Transports et urbanisme : désengorger Douala et Yaoundé.}
Le parc automobile camerounais vieillit (âge moyen \textgreater 15 ans), roulant souvent au supercarburant frelaté ou au diesel soufré, sans contrôle technique effectif. Les embouteillages quotidiens (pertes de temps, carburant brûlé au ralenti, pollution locale NOx, PM2,5) sont une source massive d'émissions évitables.
\begin{itemize}
    \item \textbf{Transports en commun structurants :} Bus BRT (Bus Rapid Transit) sur voies dédiées à Douala/Yaoundé, tramway léger, train de banlieue (ex. ligne Douala-Nkongsamba réhabilitée).
    \item \textbf{Rénovation du parc :} Interdiction progressive d'importation des véhicules \textgreater 10 ans, incitation aux véhicules hybrides/électriques (taxis, motos électriques), contrôle technique obligatoire et sanctionné.
    \item \textbf{Mobilité active :} Pistes cyclables sécurisées, trottoirs praticables, zones piétonnes en centre-ville.
    \item \textbf{Arbre de décision :} Pour choisir l'action adaptée à une source d'émission transport, suivez cette logique :
\end{itemize}

\begin{popfigure}
\begin{tikzpicture}[
    node distance=1.3cm and 0.8cm,
    decision/.style={diamond, draw=popblue, thick, fill=popblueL, minimum width=2.5cm, minimum height=1.2cm, align=center, font=\small, aspect=2},
    action/.style={rectangle, draw=popgreen, thick, fill=popgreenL, rounded corners, minimum width=3cm, minimum height=1cm, align=center, font=\small},
    start/.style={ellipse, draw=popdark, thick, fill=popcream, minimum width=2.5cm, minimum height=1cm, align=center, font=\bfseries\small},
    arrow/.style={-Stealth, thick, popdark},
    label/.style={font=\tiny, pos=0.5, sloped, above}
]
\node[start] (src) {Source d'émission identifiée};
\node[decision, below=of src] (d1) {Type de source ?};
\node[action, below left=of d1, xshift=-1.5cm, align=center] (a1){Véhicules anciens\\polluants};
\node[action, below=of a1, align=center] (a1b){Interdiction import >10 ans\\Contrôle technique strict\\Prime à la casse / Conversion};
\node[action, left=of a1b, xshift=-1cm, align=center] (a1c){Flottes publiques / Taxis :\\Électrification progressive};
\node[action, below right=of d1, xshift=1.5cm, align=center] (a2){Embouteillages\\chroniques};
\node[action, below=of a2, align=center] (a2b){Voies dédiées Bus (BRT)\\Pistes cyclables / Piétonnisations\\Urbanisme polycentrique};
\node[action, right=of a2b, xshift=1cm, align=center] (a2c){Covoiturage / Apps mobilité\\Stationnement relais (P+R)};
\node[decision, below=of d1, yshift=-3.5cm] (d2) {Transport de marchandises ?};
\node[action, below left=of d2, xshift=-1cm] (a3) {Camions vieux / Surchargés};
\node[action, below=of a3, align=center] (a3b){Renouvellement parc poids lourds\\Report modal : Fer / Fleuve (Wouri, Sanaga)\\Logistique mutualisée};
\node[action, below right=of d2, xshift=1cm, align=center] (a4){Motos taxis (benskin)\\2 temps polluants};
\node[action, below=of a4, align=center] (a4b){Conversion motos électriques\\Formation mécaniciens\\Bornes de recharge solaire};

\draw[arrow] (src) -- (d1);
\draw[arrow] (d1) -- node[label, left] {Particuliers / Taxis} (a1);
\draw[arrow] (d1) -- node[label, right] {Urbanisme / Flux} (a2);
\draw[arrow] (a1) -- (a1b);
\draw[arrow] (a1b) -- (a1c);
\draw[arrow] (a2) -- (a2b);
\draw[arrow] (a2b) -- (a2c);
\draw[arrow] (d1) -- ++(0,-2.5) -| (d2);
\draw[arrow] (d2) -- node[label, left] {Fret routier} (a3);
\draw[arrow] (d2) -- node[label, right] {Dernier km / Moto} (a4);
\draw[arrow] (a3) -- (a3b);
\draw[arrow] (a4) -- (a4b);
\end{tikzpicture}
\legende{Arbre de décision : « Quelle action d'atténuation pour quelle source d'émission dans le transport ? » La réponse dépend de la nature du problème (vétusté, congestion, fret, moto-taxi).}
\end{popfigure}

\paragraph{4. Déchets et Agriculture : méthane (CH\textsubscript{4}) et protoxyde d'azote (N\textsubscript{2}O).}
Ces gaz ont un pouvoir de réchauffement global (PRG) bien supérieur au CO\textsubscript{2} (PRG\textsubscript{100ans} : CH\textsubscript{4} $\approx$ 28--30, N\textsubscript{2}O $\approx$ 265--298).
\begin{itemize}
    \item \textbf{Déchets solides (Douala, Yaoundé, Bafoussam) :} Les décharges à ciel ouvert (ex. \textbf{Cosmos} à Douala, \textbf{Nkolfoulou} à Yaoundé) émettent du CH\textsubscript{4} par méthanisation anaérobie des déchets organiques ($\approx$ 60\% des ordures).
        \begin{itemize}
            \item \textbf{Solution :} Tri à la source (bac vert/bac jaune), \textbf{compostage} industriel des biodéchets (engrais pour agriculture périurbaine), \textbf{méthanisation contrôlée} (biogaz pour cuisson/électricité + digestat fertilisant). Fermeture/réhabilitation des décharges sauvages avec captage de biogaz.
        \end{itemize}
    \item \textbf{Agriculture / Élevage :} Riziculture inondée (plaines de Ndop, Logone) $\rightarrow$ CH\textsubscript{4}. Élevage bovin (Nord, Adamaoua) $\rightarrow$ CH\textsubscript{4} entérique + N\textsubscript{2}O (déjections). Engrais azotés mal dosés $\rightarrow$ N\textsubscript{2}O.
        \begin{itemize}
            \item \textbf{Solution :} Riziculture \textbf{SRI} (Système de Riziculture Intensive) : alternance sec/humide, moins d'eau, moins CH\textsubscript{4}, plus de rendement. \textbf{Alimentation du bétail} améliorée (fourrages de qualité, additifs) pour réduire méthane entérique. \textbf{Gestion des déjections} : biogaz domestique (digesteurs familiaux au Nord), compostage. \textbf{Agriculture de précision} : apport d'engrais au bon moment, dose, place (économie + moins N\textsubscript{2}O).
        \end{itemize}
\end{itemize}

\begin{methode}[Proposer une solution d'atténuation adaptée au contexte camerounais]
Pour toute proposition d'action (devoir, oral, projet), vérifiez ces 3 critères :
\begin{enumerate}
    \item \textbf{Efficacité ciblée :} L'action réduit-elle \emph{directement} la source d'émission identifiée ? (Ex. : planter des arbres en ville ne réduit pas les émissions des pots d'échappement, mais les capte partiellement. Remplacer une moto thermique par une électrique \emph{supprime} l'émission locale).
    \item \textbf{Faisabilité locale :} Les ressources, compétences, infrastructures existent-elles ? (Ex. : voiture électrique = bornes de recharge + électricité stable + mécaniciens formés. Au Nord, moto électrique + panneau solaire = faisable. À Douala centre, réseau instable = difficile aujourd'hui).
    \item \textbf{Acceptabilité socio-économique :} Coût pour les ménages ? Impact sur l'emploi ? Justice sociale ? (Ex. : interdire les vieilles voitures sans alternative de transport public pénalise les plus pauvres. Subventionner le compostage crée des emplois verts).
\end{enumerate}
\end{methode}

\courssub{L'adaptation : vivre avec le changement déjà engagé (protéger le sol)}

Même si nous stoppions toutes émissions aujourd'hui, l'inertie climatique (océans, calottes glacières) garantit un réchauffement supplémentaire de $\approx$ 0,5--1°C. Le Cameroun, pays à « vulnérabilité élevée, capacité d'adaptation faible » (indice ND-GAIN), doit s'adapter \textbf{maintenant}.

\paragraph{1. Côtes et ressources en eau : le front de mer et le bassin du Lac Tchad.}
\begin{itemize}
    \item \textbf{Montée des eaux (Golfe de Guinée) :} Érosion côtière à Kribi, Limbe, Douala (quartiers Bonabéri, Youpwe). \textbf{Actions :} Restauration des \textbf{mangroves} (barrière naturelle, nurserie poissons, puits de carbone « bleu »), construction de \textbf{digues} et enrochements (protection dure), \textbf{relocalisation planifiée} des populations les plus exposées (ex. île de Manoka).
    \item \textbf{Ressources en eau / Sécheresse (Nord, Extrême-Nord) :} Rétrécissement du Lac Tchad (-90\% surface depuis 1960), tarissement des mayo (rivières saisonnières). \textbf{Actions :} Barrages de retenue (Lom Pangar, Memve'ele, projets sur la Benoué), \textbf{irrigation goutte-à-goutte} solaire pour maraîchage, forage solaires, gestion intégrée des bassins versants (reboisement têtes de bassin).
\end{itemize}

\paragraph{2. Agriculture et sécurité alimentaire : nourrir 30 millions de Camerounais en 2035.}
\begin{itemize}
    \item \textbf{Variétés résilientes :} Diffusion par l'IRAD (Institut de Recherche Agricole pour le Développement) de variétés de maïs, mil, sorgho, niébé \textbf{tolérantes à la sécheresse} et aux fortes températures (cycle court 90--100 jours).
    \item \textbf{Agroécologie / Agriculture de conservation :} Couverture permanente du sol (paillis, cultures de couverture), non-labour, rotations/diversification. Améliore la rétention d'eau (+ résilience sécheresse) et stocke du carbone dans le sol (+ atténuation).
    \item \textbf{Assurance indicielle :} Produits financiers (ex. ACRE Africa, partenaires locaux) indemnisant l'agriculteur si l'indice pluviométrique (satellite) passe sous un seuil, sans expertise de terrain coûteuse.
\end{itemize}

\paragraph{3. Santé et alertes précoces : prévenir les épidémies et coups de chaleur.}
Le réchauffement étend l'aire de répartition des vecteurs (moustiques \emph{Anopheles} paludisme, \emph{Aedes} dengue/chikungunya) vers les hautes terres (Ouest, Nord-Ouest) et allonge la saison de transmission.
\begin{itemize}
    \item \textbf{Systèmes d'alerte précoce (SAP) :} Couplage météo (METEO CAMEROUN) + santé (MINSANTE) + surveillance épidémiologique. Alerte SMS / radio communautaire : « Risque canicule 40°C demain à Maroua : hydratation, heures de travail décalées, surveillance personnes âgées ». « Risque crue Logone : évacuation zones basses ».
    \item \textbf{Urbanisme santé :} Végétalisation des cours d'écoles et places publiques (îlots de fraîcheur), toits verts, corridors verts urbains (Douala, Yaoundé) pour réduire l'effet d'îlot de chaleur urbain (jusqu'à +5--7°C la nuit en centre-ville vs périphérie).
\end{itemize}

\begin{attention}[Piège majeur : l'adaptation sans atténuation est une impasse]
Ne croyez pas que « nous nous adapterons ». \textbf{Sans atténuation forte, le réchauffement dépassera +3°C, puis +4°C.} À ce niveau :
\begin{itemize}
    \item Les digues ne tiendront pas face à une montée des mers de \textgreater 1 m (effondrement calottes glaciaires).
    \item Aucune variété de maïs ne poussera sous 45°C récurrent avec 3 mois de sécheresse.
    \item Les écosystèmes (forêts, récifs coralliens, mangroves) s'effondreront, perdant leurs services (eau, pêche, bois, régulation climat).
\end{itemize}
L'adaptation a des \textbf{limites dures} (physiologiques, écologiques, financières). L'atténuation est la condition \emph{sine qua non} pour que l'adaptation reste possible et abordable.
\end{attention}

\courssub{Gestes individuels, éducation et pouvoir d'agir : « Pensez global, agissez local »}

Chaque Camerounais, chaque lycéen, a un pouvoir d'action réel. L'addition des gestes individuels change la demande, pousse les entreprises et l'État à agir, et forge la culture climatique.

\begin{exemplebox}[Bilan carbone personnel simplifié : où agir en priorité ?]
Prenons l'exemple d'un lycéen à Yaoundé :
\begin{center}
\begin{tabularx}{\linewidth}{|l|X|c|}\hline
\textbf{Poste} & \textbf{Action concrète} & \textbf{Impact CO\textsubscript{2}e/an (estimation)} \\\hline
Transport & Prendre le bus scolaire / car rapide au lieu de taxi moto (benskin) seul & $\approx$ -300 à -500 kg \\\hline
Énergie & Éteindre ventilateur/lampe en quittant la pièce ; débrancher chargeurs & $\approx$ -50 à -100 kg \\\hline
Alimentation & Manger local (plantain, macabo, ndolé) au lieu de riz/blé importés ; réduire gaspillage & $\approx$ -200 à -400 kg \\\hline
Déchets & Trier ; composter épluchures pour jardin/association ; refuser sac plastique & $\approx$ -50 à -100 kg (évite CH\textsubscript{4} décharge) \\\hline
Forêt & Participer à la \textbf{Journée Nationale de l'Arbre (1er août)} ; planter/parrainer 1 arbre/an & \textbf{+ 10 à 30 kg séquestrés/an/arbre} \\\hline
\end{tabularx}
\end{center}
\textbf{Bilan :} En changeant transport + alimentation + électricité, ce lycéen économise \textbf{$\approx$ 1 tonne CO\textsubscript{2}e/an}, soit l'équivalent des émissions d'un aller-retour Yaoundé-Paris en avion. \emph{C'est massif, mesurable, et ça commence aujourd'hui.}
\end{exemplebox}

\begin{exoresolu}[Calculer l'impact d'une action de reboisement scolaire]
\textbf{Énoncé :} Un club environnement d'un lycée de Bertoua (Est) plante 200 plants d'\emph{Acacia mangium} (espacement 2,5m x 2,5m) sur une friche de 1,25 ha. On admet qu'une plantation d'Acacia de 10 ans séquestre en moyenne 15 t CO\textsubscript{2}/ha/an. Le lycée émet environ 50 t CO\textsubscript{2}/an (électricité, générateur, transport personnel). Combien d'années pour compenser 1 an d'émissions du lycée ? Quelle surface aurait fallu planter pour être « neutre en carbone » dès la 10\textsuperscript{e} année ?

\textbf{Solution détaillée :}
\begin{enumerate}
    \item \textbf{Séquestration annuelle de la plantation :} Surface = 1,25 ha. Taux = 15 t CO\textsubscript{2}/ha/an.
    \[ \text{Séquestration/an} = 1,25 \times 15 = \mathbf{18,75\ t\ CO_2/an} \]
    \item \textbf{Années pour compenser 1 an d'émissions (50 t) :}
    \[ \text{Durée} = \frac{50}{18,75} \approx \mathbf{2,67\ ans} \ (\text{soit 2 ans et 8 mois}). \]
    \item \textbf{Surface pour neutralité à 10 ans (stock cumulé = émissions annuelles $\times$ 10 = 500 t) :}
    Stock cumulé sur 10 ans par ha = $15 \times 10 = 150$ t CO\textsubscript{2}/ha.
    Surface requise = $\frac{500}{150} \approx \mathbf{3,33\ ha}$.
    \item \textbf{Interprétation :} 1,25 ha compensent 1 an d'émissions en $\approx$ 2,7 ans. Pour absorber 10 ans d'émissions sur la durée de vie de la plantation (10 ans), il faut $\approx$ 3,3 ha. Le club doit viser l'extension progressive (ex. 0,5 ha/an pendant 6 ans) et protéger la plantation (feux de brousse, pâturage) pour garantir la permanence du stockage.
\end{enumerate}
\end{exoresolu}

\begin{savaistu}[Le Cameroun à la COP : acteur de la « Coalition pour les Forêts du Bassin du Congo »]
Le Cameroun, par la voix du Président de la République à la COP27 (Charm el-Cheikh, 2022) et la COP28 (Dubaï, 2023), porte la revendication des pays du bassin du Congo : \textbf{« Pollueur-payeur » inversé.} Nos forêts rendent un service planétaire (stockage carbone, régulation pluies Sahel, biodiversité) sans compensation financière à la hauteur. Le mécanisme \textbf{REDD+ (Réduction des Emissions liées à la Déforestation et à la Dégradation des forêts)} et les \textbf{marchés carbone (Article 6 Accord de Paris)} doivent financer la conservation \emph{et} le développement local. En 2023, le Cameroun a signé un accord-cadre avec la Banque Mondiale (FCPF) pour des paiements pour résultats (réduction vérifiée de déforestation). Votre génération négociera ces contrats : comprenez la comptabilité carbone (MRV : Mesure, Rapportage, Vérification) !
\end{savaistu}

\courssub{Le rôle clé de l'éducation et de la sensibilisation des jeunes}

L'article 12 de l'Accord de Paris (Action pour l'Accès au Climat - ACE) mandate l'éducation, la formation, la sensibilisation, l'accès à l'information et la participation du public. Au Cameroun, le MINESEC intègre l'EEDD (Éducation à l'Environnement et au Développement Durable) dans les programmes (SVT, Géographie, Éducation à la Citoyenneté).
\begin{itemize}
    \item \textbf{Compétence visée :} « Analyser une situation environnementale locale pour proposer des actions d'atténuation/adaptation. » C'est \emph{exactement} ce que vous faites en résolvant l'exercice ci-dessus ou en menant un diagnostic énergétique de votre lycée.
    \item \textbf{Projets concrets :} « Lycées verts » (tri, compost, pépinière, solaire), « Jeunes Reporters pour l'Environnement » (enquêtes terrain, diffusion radio/réseaux sociaux), participation aux consultations pour les \textbf{Contributions Déterminées au niveau National (CDN/NDC)} du Cameroun (révisées tous les 5 ans, prochaine échéance 2025).
    \item \textbf{Carrières vertes :} Ingénieur carbone / MRV, technicien énergie solaire, agroforesterie, gestionnaire d'aires protégées, urbaniste climat, juriste droit de l'environnement. Le marché de l'emploi vert explose.
\end{itemize}

\begin{aretenir}[Synthèse : Votre boîte à outils « Action Climat »]
\begin{itemize}
    \item \textbf{Atténuation (Cause) :} Énergies renouvelables (Hydro/Solaire) $\rightarrow$ Efficacité $\rightarrow$ Forêts (Reboisement/Agroforesterie/Zero Déforestation) $\rightarrow$ Transports durables $\rightarrow$ Déchets/Agri (Méthanisation/Compost/SRI).
    \item \textbf{Adaptation (Effet) :} Côtes (Mangroves/Digues) $\rightarrow$ Eau (Retenues/Goutte-à-goutte) $\rightarrow$ Agri (Variétés courtes/Agroécologie/Assurance) $\rightarrow$ Santé (Alertes/Îlots fraîcheur) $\rightarrow$ Urbanisme résilient.
    \item \textbf{Individuel :} Transport $\rightarrow$ Alimentation locale $\rightarrow$ Énergie/Gaspi $\rightarrow$ Déchets $\rightarrow$ Planter/Parrainer arbre.
    \item \textbf{Collectif :} Vote / Plaidoyer $\rightarrow$ Projets lycée/quartier $\rightarrow$ Métiers verts $\rightarrow$ Suivi CDN Cameroun.
\end{itemize}
\end{aretenir}

\begin{exercice}[Mise en situation : Diagnostic et plan d'action pour votre établissement]
Votre lycée (1500 élèves, zone périurbaine Yaoundé) dispose : d'un générateur diesel (coupures fréquentes), d'une cantine (gaspi alimentaire, déchets plastiques), d'une cour bétonnée (chaleur extrême midi), de 2 ha de terrain vague en friche, d'un forage pompe thermique. Les élèves viennent majoritairement en taxi moto (benskin) ou taxi partagé.
\begin{enumerate}
    \item Réalisez un \textbf{diagnostic simplifié des sources d'émission} (scope 1 : directes ; scope 2 : électricité ; scope 3 : transport, déchets, achats).
    \item Proposez \textbf{3 actions d'atténuation} (une par secteur prioritaire) et \textbf{2 actions d'adaptation}, en justifiant chacune par les critères de la \textbf{Méthode} (efficacité, faisabilité, acceptabilité).
    \item Chiffrez approximativement le potentiel de réduction CO\textsubscript{2}e/an pour l'action « installation 20 kWc solaire + batteries pour remplacer générateur » (émissions générateur $\approx$ 2,5 kg CO\textsubscript{2}/L ; conso 10 L/jour ; 200 jours/an).
\end{enumerate}
\end{exercice}

\begin{corrige}[Exercice n°1 - Éléments de correction]
\begin{enumerate}
    \item \textbf{Diagnostic :}
        \begin{itemize}
            \item Scope 1 : Générateur diesel (CO\textsubscript{2}, NOx, bruit), Pompe forage thermique, Déchets cantine (CH\textsubscript{4} si enfouis).
            \item Scope 2 : Électricité réseau (mix hydro/thermique $\rightarrow$ émissions indirectes).
            \item Scope 3 : Transport élèves/enseignants (benskin/taxi $\rightarrow$ majoritaire), Achats cantine (riz importé, emballages), Déplacements admin.
        \end{itemize}
    \item \textbf{Propositions (exemples validés) :}
        \begin{itemize}
            \item \textbf{Atténuation 1 (Énergie) :} Installation solaire 20 kWc + stockage (remplace générateur). \textit{Justification :} Efficacité directe sur 1\textsuperscript{re} source ; faisable (entreprises locales, soleil) ; acceptabilité forte (économie carburant, silence).
            \item \textbf{Atténuation 2 (Déchets/Agri) :} Compostage biodéchets cantine + jardin pédagogique sur friche (amendement compost). \textit{Justification :} Évite CH\textsubscript{4} décharge ; produit engrais ; support pédagogique SVT ; faible coût.
            \item \textbf{Atténuation 3 (Transport) :} Convention avec syndicat taxis/bus pour lignes dédiées « Lycée » aux heures de pointe + abri sécurisé + sensibilisation covoiturage enseignants. \textit{Justification :} Cible scope 3 major ; faisable (dialogue social) ; acceptabilité si gain temps/argent.
            \item \textbf{Adaptation 1 (Chaleur) :} Végétalisation cour (arbres d'ombrage \emph{Terminalia mantaly}, \emph{Delonix regia}) + toits peints blanc (albédo) + brumisateurs solaires. \textit{Justification :} Réduit T° ressentie 3--5°C ; améliore concentration ; faisable (pépinière MINFOF/Communes).
            \item \textbf{Adaptation 2 (Eau) :} Récupération eau de pluie toits (cuves 10 m\textsuperscript{3}) pour arrosage jardin / lavage sols / chasse d'eau. \textit{Justification :} Sécurise ressource en cas coupure CAMWATER / sécheresse ; économise eau traitée ; simple.
        \end{itemize}
    \item \textbf{Chiffrage Solaire :}
        \begin{align*}
            \text{Conso annuelle diesel} &= 10\ \text{L/j} \times 200\ \text{j} = 2000\ \text{L/an} \\
            \text{Émissions évitées} &= 2000\ \text{L} \times 2,5\ \text{kg CO}_2/\text{L} = \mathbf{5000\ kg\ CO_2/an} = \mathbf{5\ t\ CO_2/an}.
        \end{align*}
        \textit{Note :} L'empreinte carbone de la fabrication des panneaux (« dette carbone ») se rembourse en $\approx$ 1--2 ans au Cameroun (fort ensoleillement). Sur 25 ans, gain net $\approx$ 120 t CO\textsubscript{2}.
\end{enumerate}
\end{corrige}

\coursec{Les politiques climatiques : du local à l'international}

Après avoir vu comment l'atténuation (réduire les émissions) et l'adaptation (vivre avec le changement) s'articulent concrètement sur le terrain, une question s'impose : \emph{qui décide, qui paye, et comment s'assurer que tous les pays avancent dans la même direction ?} Le climat ne connaît pas de frontières : le $\ce{CO2}$ émis à Douala ou à Yaoundé se mélange à l'atmosphère planétaire en quelques semaines et influence la température à New York ou à Pékin. Cette section explore comment l'humanité tente de gouverner ce \cle{bien commun mondial} qu'est le climat, du sommet de Rio à l'engagement concret du Cameroun.

\courssub{Pourquoi une coopération internationale ? Le climat, bien commun mondial}

\emph{Observation :} Imaginez un grand lac partagé par plusieurs villages. Si un village rejette ses déchets dans l'eau, tous les autres en subissent la pollution. Pourtant, chaque village a intérêt à continuer de polluer pour économiser le coût d'une station d'épuration, en espérant que les autres feront l'effort à sa place. C'est le \textbf{dilemme du passager clandestin} (ou \emph{free rider}).

\emph{Hypothèse :} Sans règle commune contraignante, la somme des intérêts individuels mène à la ruine collective (la « tragédie des communs » décrite par Garrett Hardin).

\emph{Données :} L'atmosphère est un \cle{bien commun mondial} : elle est non-exclusive (personne ne peut en être exclu) et rivale (l'absorption de $\ce{CO2}$ par l'atmosphère diminue sa capacité d'absorption pour les autres). Les émissions de gaz à effet de serre (GES) sont une \cle{externalité négative} : leur coût (réchauffement, catastrophes) est supporté par l'humanité entière, pas seulement par l'émetteur.

\begin{aretenir}[Le problème fondamental de la gouvernance climatique]
L'atmosphère est un bien commun planétaire. Les émissions d'un pays affectent tous les autres (externalité négative transfrontière). Sans coordination, chaque pays a un intérêt individuel à ne pas réduire ses émissions (passager clandestin), ce qui conduit collectivement à une catastrophe. La coopération internationale vise à internaliser cette externalité par des accords contraignants ou incitatifs.
\end{aretenir}

\courssub{Les grandes étapes de la gouvernance climatique internationale}

La construction d'un régime climatique mondial suit une démarche d'investigation politique : \emph{observation scientifique (GIEC) $\rightarrow$ alerte $\rightarrow$ négociation $\rightarrow$ accord $\rightarrow$ mise en œuvre $\rightarrow$ révision}.

\begin{popfigure}
\begin{tikzpicture}[
    timeline/.style={very thick, popblue},
    event/.style={circle, fill=popblue, inner sep=2pt},
    eventlabel/.style={anchor=west, font=\small, text width=4.5cm, align=left},
    yearlabel/.style={anchor=north, font=\small\bfseries, text=popdark},
    milestoneline/.style={dashed, popmuted, thin}
]
% Ligne de temps
\draw[timeline] (0,0) -- (16,0);
% Années
\foreach \x/\year in {0/1992, 3.5/1997, 8.5/2009, 12.5/2015, 16/2023} {
    \draw[milestoneline] (\x,0) -- (\x,-0.3);
    \node[yearlabel] at (\x,-0.5) {\year};
}
% Événements
% 1992 Rio
\node[event] (rio) at (0,0) {};
\node[eventlabel] at (0.3,0.8) {\textbf{Sommet de la Terre (Rio)}\newline Convention-cadre (CCNUCC)\newline Principe RCPD};
% 1997 Kyoto
\node[event] (kyoto) at (3.5,0) {};
\node[eventlabel] at (3.8,0.8) {\textbf{Protocole de Kyoto}\newline Engagements chiffrés (Annexe I)\newline Mécanismes de flexibilité};
% 2009 Copenhague (échec relatif)
\node[event] (cph) at (8.5,0) {};
\node[eventlabel, anchor=east] at (8.2,0.8) {\textbf{Copenhague (COP15)}\newline Accord non contraignant\newline nouvelle dynamique};
% 2015 Paris
\node[event] (paris) at (12.5,0) {};
\node[eventlabel] at (12.8,0.8) {\textbf{Accord de Paris}\newline Objectif $<2^\circ$C / $1,5^\circ$C\newline CDN tous les 5 ans};
% Aujourd'hui
\node[event] (now) at (16,0) {};
\node[eventlabel, anchor=east] at (15.7,0.8) {\textbf{Global Stocktake (Bilan mondial)}\newline Réhaussement de l'ambition\newline adaptation/pertes};
\end{tikzpicture}
\legende{Frise chronologique des grandes étapes de la gouvernance climatique internationale (1992--2023).}
\end{popfigure}

\textbf{1992 -- Rio de Janeiro : La Convention-Cadre (CCNUCC).}\par
Le « Sommet de la Terre » pose les fondations juridiques. La \cle{CCNUCC} (Convention-Cadre des Nations Unies sur les Changements Climatiques) entre en vigueur en 1994. Elle établit le principe cardinal des \textbf{responsabilités communes mais différenciées (RCPD)} : tous les pays ont la responsabilité de protéger le climat, mais les pays développés (historiquement grands émetteurs) doivent agir les premiers et aider les pays en développement (financement, transfert technologique). Aucune cible chiffrée de réduction n'est encore fixée.

\textbf{1997 -- Kyoto : Premiers engagements chiffrés.}\par
Le \cle{Protocole de Kyoto} (entré en vigueur en 2005) engage \emph{uniquement} les pays industrialisés (pays de l'« Annexe I ») à réduire leurs émissions globales d'environ \SI{5}{\%} par rapport à 1990 sur la période 2008--2012. Il introduit des \emph{mécanismes de flexibilité} (marché du carbone, Mise en œuvre conjointe, Mécanisme pour un développement propre - MDP) pour réduire le coût de l'action. \textbf{Limite :} les États-Unis ne le ratifient pas ; les grands émergents (Chine, Inde) n'ont pas d'objectifs contraignants.

\textbf{2015 -- Paris : Un accord universel et dynamique.}\par
La COP21 aboutit à l'\cle{Accord de Paris}, juridiquement contraignant, ratifié par 195 Parties. Trois piliers :
\begin{itemize}
    \item \textbf{Objectif température :} Maintenir le réchauffement « bien en dessous de \SI{2}{\celsius} » et poursuivre les efforts pour le limiter à \SI{1,5}{\celsius} (par rapport à l'ère préindustrielle).
    \item \textbf{CDN (Contributions Déterminées au niveau National) :} \emph{Chaque} pays (développé ou non) dépose sa propre trajectoire de réduction (atténuation) et d'adaptation, révisable \emph{tous les 5 ans} (mécanisme de « cliquet » : on ne peut que renforcer l'ambition).
    \item \textbf{Transparence et bilan mondial (Global Stocktake) :} Tous les 5 ans, un bilan collectif évalue l'écart entre les engagements et l'objectif \SI{1,5}{\celsius}.
\end{itemize}

\begin{exemplebox}[Objectif de l'Accord de Paris -- Chiffres clés Probatoire]
L'Accord de Paris (2015) vise à limiter l'élévation de la température moyenne mondiale à \textbf{bien en dessous de $+2\,^\circ\mathrm{C}$} par rapport aux niveaux préindustriels, et à \textbf{poursuivre les efforts pour la limiter à $+1,5\,^\circ\mathrm{C}$}. 
\begin{itemize}
    \item \textbf{Préindustriel :} référence 1850--1900.
    \item \textbf{État actuel :} nous sommes déjà à environ $+1,2\,^\circ\mathrm{C}$ (année 2023 record à $+1,45\,^\circ\mathrm{C}$).
    \item \textbf{CDN :} Contributions Déterminées au niveau National. Déposées tous les 5 ans (2020, 2025, 2030...). Elles contiennent des objectifs d'atténuation (réduction GES) et d'adaptation.
\end{itemize}
\end{exemplebox}

\begin{attention}[Piège classique : Couche d'ozone vs Effet de serre]
Ne confondez \textbf{JAMAIS} les deux problèmes à l'examen :
\begin{center}
\begin{tabularx}{\linewidth}{|l|X|X|}\hline
\rowcolor{popblueL} \textbf{Problème} & \textbf{Couche d'ozone (Stratosphère)} & \textbf{Effet de serre (Troposphère)} \\ \hline
\textbf{Gaz responsables} & CFC, halons (Protocole de \textbf{Montréal}, 1987) & $\ce{CO2}$, $\ce{CH4}$, $\ce{N2O}$, vapeur d'eau (Accord de \textbf{Paris}, 2015) \\ \hline
\textbf{Conséquence} & « Trou » d'ozone $\rightarrow$ UV dangereux (cancers, cataractes) & Réchauffement global $\rightarrow$ météo extrême, montée des mers \\ \hline
\textbf{Statut} & Succès : couche en voie de guérison (protocole respecté) & Urgence : émissions encore en hausse, trajectoire $>2,5\,^\circ\mathrm{C}$ \\ \hline
\end{tabularx}
\end{center}
\textbf{Astuce :} Montréal = Ozone (M comme « Mauvais UV ») ; Paris = Climat (P comme « Planète qui chauffe »).
\end{attention}

\begin{aretenir}[Dates et objectifs à connaître par cœur]
\begin{itemize}
    \item \textbf{1992 (Rio) :} CCNUCC -- Cadre général, principe RCPD.
    \item \textbf{1997 (Kyoto) :} 1er protocole contraignant \emph{pour pays développés seulement} ($-5\%$ vs 1990).
    \item \textbf{2015 (Paris) :} Accord universel -- Objectif $<2\,^\circ\mathrm{C}$ / $1,5\,^\circ\mathrm{C}$ -- CDN tous les 5 ans -- Bilan mondial (Global Stocktake).
\end{itemize}
\textbf{Question type :} « Citez l'accord qui a instauré le principe de responsabilités communes mais différenciées. » $\rightarrow$ CCNUCC (Rio 1992). « Quel est l'objectif de température de l'Accord de Paris ? » $\rightarrow$ Bien en dessous de $2\,^\circ\mathrm{C}$, efforts vers $1,5\,^\circ\mathrm{C}$.
\end{aretenir}

\courssub{Le rôle central du GIEC : la science au service de la décision}

Comment les négociateurs savent-ils qu'il faut viser $1,5\,^\circ\mathrm{C}$ et non $3\,^\circ\mathrm{C}$ ? Ils s'appuient sur le \cle{GIEC} (Groupe d'experts Intergouvernemental sur l'Évolution du Climat), créé en 1988 par l'OMM et le PNUE.

\begin{definition}[GIEC -- Groupe d'experts Intergouvernemental sur l'Évolution du Climat]
Organisme scientifique intergouvernemental chargé d'évaluer de manière exhaustive, objective, ouverte et transparente les informations scientifiques, techniques et socio-économiques pertinentes pour comprendre le changement climatique d'origine humaine, ses impacts potentiels et les options d'atténuation et d'adaptation. Il ne fait \emph{pas} de recherche originale : il \textbf{synthétise} la littérature peer-reviewed (des milliers d'articles) en \textbf{Rapports d'évaluation} (tous les 6--7 ans) et \textbf{Rapports spéciaux}.
\end{definition}

\emph{Fonctionnement :} Trois groupes de travail :
\begin{enumerate}
    \item \textbf{GI} : Bases physiques (réchauffement, cycle du carbone, scénarios SSP).
    \item \textbf{GII} : Impacts, adaptation, vulnérabilité.
    \item \textbf{GIII} : Atténuation (énergie, transport, industrie, AFOLU\footnote{Agriculture, Foresterie et Autres Utilisations des Terres.}).
\end{enumerate}
Les \textbf{Résumés à l'intention des décideurs (SPM -- Summary for Policymakers)} sont approuvés ligne par ligne par les gouvernements : c'est la \emph{traduction politique consensuelle} de la science.

\begin{savaistu}[Le GIEC et le Cameroun]
Des scientifiques camerounais (climatologues, agronomes, géographes) participent aux rapports du GIEC comme auteurs principaux ou relecteurs. Le \emph{Centre Régional AGRHYMET} (Niamey) et l'\emph{Observatoire du Climat pour l'Afrique Centrale} relaient les données du GIEC pour adapter les modèles climatiques aux réalités du bassin du Congo et de la zone soudano-sahélienne.
\end{savaistu}

\courssub{Les échelles d'action emboîtées : de l'individu à la communauté internationale}

L'action climatique n'est pas l'affaire des seuls chefs d'État. Elle opère par \textbf{subsidiarité} : chaque niveau agit là où il est le plus efficace, en cohérence avec les niveaux supérieurs.

\begin{popfigure}
\begin{tikzpicture}[
    level/.style={rounded corners=1.5cm, draw=popblue, thick, minimum height=2.2cm, minimum width=4.5cm, align=center, font=\small},
    arrow/.style={-Stealth, thick, poporange},
    label/.style={font=\bfseries\small, text=popdark, anchor=north}
]
% Niveau 4 : International
\node[level, fill=popblue!10, minimum width=14cm, align=center] (int) at (0,0){
    \textbf{COMMUNAUTÉ INTERNATIONALE} (ONU, CCNUCC, GIEC, Fonds Vert)\\
    Fixe le cadre global (Objectif $1,5\,^\circ\mathrm{C}$), finance, transfert techno, bilan mondial
};
\node[label] at (int.north) (lint) {Échelle planétaire};

% Niveau 3 : État (Cameroun)
\node[level, fill=popgreen!10, minimum width=10cm, align=center] (etat) at (0,-3.2){
    \textbf{ÉTAT (Cameroun)} -- CDN, Loi-cadre environnement, Plan National d'Adaptation (PNA)\\
    Politique forestière, énergie (barrages, solaire), fiscalité carbone, éducation
};
\node[label] at (etat.north) (letat) {Échelle nationale};

% Niveau 2 : Collectivités / Communes
\node[level, fill=poporange!10, minimum width=6.5cm, align=center] (comm) at (0,-6.4){
    \textbf{COMMUNES \& RÉGIONS}\\
    Plan Climat Air Énergie Territorial (PCAET), transport urbain, déchets, PLU, reboisement local
};
\node[label] at (comm.north) (lcomm) {Échelle locale / régionale};

% Niveau 1 : Individu / Ménage / Entreprise
\node[level, fill=poppink!10, minimum width=4.5cm, align=center] (ind) at (0,-9.6){
    \textbf{INDIVIDU / MÉNAGE / ENTREPRISE}\\
    Mobilité douce, efficacité énergétique, tri, alimentation locale, consommation responsable, RSE
};
\node[label] at (ind.north) (lind) {Échelle micro / individuelle};

% Flèches de cohérence
\draw[arrow] (int.south) -- (etat.north) node[midway, right, font=\tiny, text=popdark] {Cadre, financements, objectifs};
\draw[arrow] (etat.south) -- (comm.north) node[midway, right, font=\tiny, text=popdark] {Loi, décentralisation, moyens};
\draw[arrow] (comm.south) -- (ind.north) node[midway, right, font=\tiny, text=popdark] {Services publics, incitations, normes};
\draw[arrow, bend left=20] (ind.east) to node[midway, right, font=\tiny, text=popdark] {Pression citoyenne, vote, innovation} (int.east);
\end{tikzpicture}
\legende{Les échelles d'action emboîtées pour le climat : l'international fixe le cap, l'État (Cameroun) légifère et planifie, les collectivités territorialisent, l'individu et l'entreprise agissent au quotidien. La cohérence verticale (haut $\leftrightarrow$ bas) et horizontale (secteurs) est la clé de l'efficacité.}
\end{popfigure}

\emph{Exemple concret au Cameroun :}
\begin{itemize}
    \item \textbf{International :} L'Accord de Paris exige une CDN mise à jour tous les 5 ans.
    \item \textbf{National :} Le Cameroun dépose sa CDN révisée (2021), vote une \emph{Loi-cadre sur l'environnement} (1996, révisée), crée le \emph{Fonds National pour l'Environnement et le Développement Durable (FNE)}.
    \item \textbf{Local (Commune de Yaoundé VI ou Douala V) :} Élaboration d'un \emph{Plan Climat Air Énergie Territorial (PCAET)} : pistes cyclables, tri sélectif des ordures, plantation d'arbres d'alignement, marché bio hebdomadaire.
    \item \textbf{Individuel :} Un élève de Première à Ngaoundéré choisit le vélo ou la marche pour aller au lycée, refuse les sachets plastiques (interdits mais persistants), sensibilise sa famille à la cuisson améliorée (foyers « foyers améliorés » type \emph{Mayon Turbo} ou \emph{Green Charcoal}) pour réduire le bois de chauffe et la déforestation.
\end{itemize}

\courssub{L'engagement du Cameroun : la CDN nationale et les actions concrètes}

Le Cameroun, bien que faible émetteur (\textasciitilde \SI{0,1}{\%} des émissions mondiales), est \textbf{très vulnérable} (agriculture pluviale, côtes basses, Mont Cameroun, bassin du Lac Tchad). Sa \cle{CDN révisée (2021)} s'inscrit dans la \emph{Stratégie Nationale de Développement (SND30)} et la \emph{Vision 2035}.

\begin{savaistu}[La CDN du Cameroun -- Chiffres 2021--2035]
Le Cameroun s'engage à réduire ses émissions de GES d'environ \textbf{32 \% d'ici 2035} par rapport à un scénario de référence (Business As Usual), \textbf{sous condition de soutien international} (financement, technologie, renforcement des capacités).
\begin{itemize}
    \item \textbf{Inconditionnel (effort propre) :} \textasciitilde 12 \% de réduction.
    \item \textbf{Conditionnel (avec aide) :} \textasciitilde 20 \% supplémentaires.
    \item \textbf{Secteurs clés :} Forêt (REDD+), Énergie (hydroélectricité, solaire), Agriculture (agroforesterie), Déchets (méthanisation).
\end{itemize}
\end{savaistu}

\textbf{1. La forêt : poumon du bassin du Congo et levier REDD+.} \par
Le Cameroun abrite \textasciitilde \SI{22}{millions d'hectares} de forêts (2\textsuperscript{e} massif africain après la RDC). La politique forestière (Loi 1994, révisée 2023) vise :
\begin{itemize}
    \item \textbf{Zéro déforestation nette} à l'horizon 2030 (engagement de Glasgow COP26).
    \item \textbf{REDD+ :} Réduction des émissions dues à la déforestation et à la dégradation des forêts, conservation, gestion durable, augmentation des stocks de carbone. Projets pilotes dans le \emph{Tri-National de la Sangha (TNS)} et le \emph{Paysage TRIDOM}.
    \item \textbf{Reboisement / Restauration :} Initiative \emph{« Green Sahel »} (Grande Muraille Verte), objectif \SI{12}{millions d'hectares} restaurés à l'horizon 2030 (engagement AFR100).
\end{itemize}

\textbf{2. L'énergie : hydroélectricité et solaire.} \par
\begin{itemize}
    \item \textbf{Barrages :} Nachtigal (\SI{420}{MW}, mis en service 2023--2024), Memve'ele (\SI{211}{MW}), projet Grand Eweng (\SI{1000}{MW} potentiel). Objectif : \SI{3000}{MW} à l'horizon 2030 (vs \textasciitilde \SI{1500}{MW} en 2020).
    \item \textbf{Solaire :} Centrale de Maroua (\SI{30}{MW}), Guider, projets solaires flottants sur barrages (hybridation). Programme \emph{« Énergie pour Tous »} (électrification rurale solaire).
    \item \textbf{Transition :} Remplacer le thermique (fioul/lourd) à la base (Kribi, Dibamba) par l'hydro/solaire.
\end{itemize}

\textbf{3. Agriculture et résilience.} \par
\begin{itemize}
    \item \textbf{Agroforesterie :} Systèmes cacaoyers/caféiers sous couvert forestier (zones Sud/Centre), parcs à karité/niébé (Nord).
    \item \textbf{Variétés adaptées :} Diffusion de semences à cycle court, tolérantes à la sécheresse (maïs \emph{CMS 8704}, sorgho \emph{S 35}, riz \emph{NERICA}).
    \item \textbf{Élevage :} Amélioration des pâturages, biogaz à partir de bouse (réduction $\ce{CH4}$, énergie propre).
\end{itemize}

\begin{exoresolu}[Analyse de la cohérence d'un projet camerounais avec l'Accord de Paris]
\textbf{Énoncé :} Le projet « \emph{Nachtigal Hydro Power Company (NHPC)} » (\SI{420}{MW} sur la Sanaga) est présenté comme un projet phare de la CDN camerounaise. En vous appuyant sur les mécanismes de l'Accord de Paris, expliquez en quoi ce projet contribue à l'atteinte de l'objectif \SI{1,5}{\celsius} et identifiez un risque d'incohérence (maladaptation) potentiel.

\textbf{Solution détaillée :}
\begin{enumerate}
    \item \textbf{Contribution à l'atténuation (Objectif \SI{1,5}{\celsius}) :}
    \begin{itemize}
        \item Remplace la production thermique (centrales fioul/gaz : Kribi, Dibamba, Limbe) qui émet \textasciitilde \SI{600}{g \ce{CO2}/kWh}.
        \item L'hydroélectricité \emph{au fil de l'eau} (Nachtigal) émet \textasciitilde \SI{10-20}{g \ce{CO2}/kWh} (cycle de vie).
        \item Évitement estimé : \textasciitilde \SI{1,5}{Mt \ce{CO2}/an}. Cela compte dans le volet « Énergie » de la CDN conditionnelle du Cameroun.
        \item S'inscrit dans le \textbf{Mécanisme de Développement Propre (MDP) / Article 6.4 Accord de Paris} : possibilité de générer des \emph{crédits carbone} (ITMO -- Internationally Transferred Mitigation Outcomes) vendus à des pays/pollueurs pour financer le projet $\rightarrow$ financement climatique international.
    \end{itemize}
    \item \textbf{Risque de maladaptation / Incohérence (Justice climatique locale) :}
    \begin{itemize}
        \item \textbf{Déplacement de populations :} \textasciitilde 3000 personnes relogées. Perte de terres agricoles, de sites sacrés, perturbation des pêches (migration des poissons bloquée par le barrage malgré passes à poissons).
        \item \textbf{Effet de serre du réservoir :} Dans les zones tropicales, la mise en eau noie la végétation $\rightarrow$ décomposition anaérobie $\rightarrow$ émission de \textbf{méthane ($\ce{CH4}$)}, puissant GES. Bien que faible pour du fil de l'eau, non nul.
        \item \textbf{Vulnérabilité climatique :} Si les sécheresses s'intensifient (scénario GIEC), le débit de la Sanaga baisse $\rightarrow$ production électrique chute $\rightarrow$ retour au thermique (charbon/fioul) $\rightarrow$ cercle vicieux. \textbf{Diversification (solaire/éolien) indispensable.}
    \end{itemize}
    \item \textbf{Conclusion :} Projet \textbf{aligné sur l'atténuation globale} (CDN, Art. 6), mais nécessite des \textbf{mesures d'adaptation sociale/environnementale fortes} (Plan de Gestion Environnementale et Sociale - PGES rigoureux) pour ne pas violer le principe « ne pas nuire » (Do No Significant Harm) de la taxonomie verte.
\end{enumerate}
\end{exoresolu}

\courssub{Justice climatique et responsabilités communes mais différenciées (RCPD)}

\emph{Problème éthique :} Les pays les moins responsables (Afrique subsaharienne : \textasciitilde \SI{3-4}{\%} des émissions historiques) subissent les impacts les plus durs (sécheresses, inondations, perte de rendements, maladies vectorielles). Les pays riches (OCDE, Chine) ont bâti leur développement sur les énergies fossiles.

\begin{propriete}[Principe des Responsabilités Communes mais Différenciées (RCPD) -- Art. 3.1 CCNUCC]
\begin{enumerate}
    \item \textbf{Responsabilité commune :} Tous les États doivent protéger le climat (bien commun).
    \item \textbf{Responsabilité différenciée :} Les pays développés portent une responsabilité \emph{plus grande} du fait de leurs émissions historiques (stock de $\ce{CO2}$ dans l'atmosphère) et de leurs capacités financières/technologiques supérieures.
    \item \textbf{Conséquences opérationnelles :}
    \begin{itemize}
        \item \textbf{Atténuation :} Pays développés $\rightarrow$ objectifs absolus de réduction (ex: UE $-55\%$ 2030 vs 1990). Pays en développement $\rightarrow$ objectifs relatifs (intensité carbone, déviation par rapport au BAU) ou conditionnels.
        \item \textbf{Financement :} Objectif \textbf{\$100 milliards/an} (dès 2020, prolongé à 2025) des pays développés vers pays en développement (Fonds Vert pour le Climat -- FVC). \textbf{Nouvel objectif post-2025 (NCQG) en négociation.}
        \item \textbf{Transfert technologique :} Accès facilité aux brevets verts (solaire, batteries, capture carbone, variétés résilientes).
        \item \textbf{Pertes et Préjudices (Loss \& Damage) :} Mécanisme financier spécifique adopté à la COP27 (Charm el-Cheikh) et opérationnalisé COP28 (Dubaï) pour indemniser les impacts \emph{irréversibles} (îles submergées, glaciers disparus, espèces éteintes) au-delà de l'adaptation possible.
    \end{itemize}
\end{enumerate}
\end{propriete}

\emph{Illustration camerounaise :} Le Cameroun réclame l'accès au \textbf{Fonds Vert pour le Climat (FVC)} pour financer la partie \emph{conditionnelle} de sa CDN (-\SI{20}{\%} supplémentaires). Projets financés : \emph{Projet d'Adaptation aux Changements Climatiques dans les Zones Côtières (PACC-ZC)} (lutte contre l'érosion côtière à Kribi, Limbé), \emph{Projet d'Agroforesterie Résiliente dans le Grand Nord}.

\begin{methode}[Réussir une question de Probatoire sur les politiques climatiques]
\begin{enumerate}
    \item \textbf{Identifiez l'accord cité} (Rio/CCNUCC, Kyoto, Paris) $\rightarrow$ citez la \textbf{date} et l'\textbf{objectif principal}.
    \item \textbf{Distinuez Atténuation / Adaptation / Pertes \& Préjudices.} Ne les mélangez pas.
    \item \textbf{Appliquez le principe RCPD :} Qui paie ? Qui réduit le plus ? Qui aide ?
    \item \textbf{Contextualisez au Cameroun :} Citez la CDN (\SI{32}{\%} conditionnel), un projet concret (Nachtigal, REDD+, Grande Muraille Verte), un risque (déforestation, érosion côtière).
    \item \textbf{Évitez le piège Ozone/Serre} (Montréal $\neq$ Paris).
\end{enumerate}
\end{methode}

\begin{exercice}[Entraînement Probatoire -- Politiques climatiques]
\begin{enumerate}
    \item Citez les trois piliers de l'Accord de Paris (2015).
    \item Expliquez la différence entre la logique de Kyoto (1997) et celle de Paris (2015) concernant la répartition des efforts.
    \item Le Cameroun vise \SI{32}{\%} de réduction d'ici 2035 dans sa CDN. Précisez la part inconditionnelle et la part conditionnelle. Donnez \textbf{deux} exemples de mesures du secteur « Forêt » et \textbf{deux} du secteur « Énergie » pour atteindre cet objectif.
    \item Définissez le principe RCPD. Donnez une conséquence concrète sur le financement climatique (Fonds Vert).
    \item \textbf{Piège :} Un élève affirme : « Le Protocole de Montréal a réglé le problème du réchauffement climatique en interdisant les CFC. » Corrigez cette affirmation en distinguant les deux problèmes.
\end{enumerate}
\end{exercice}

\begin{corrige}[Exercice -- Politiques climatiques]
\begin{enumerate}
    \item \textbf{Trois piliers de Paris :} (1) Objectif température ($<2\,^\circ\mathrm{C}$, viser $1,5\,^\circ\mathrm{C}$). (2) CDN universelles, révisées tous les 5 ans (mécanisme de cliquet). (3) Transparence / Bilan mondial (Global Stocktake) + Financement (\$100 Md/an) + Transfert techno + Pertes \& Préjudices.
    \item \textbf{Kyoto :} Approche « haut-bas » (top-down). Objectifs contraignants \emph{uniquement} pour pays développés (Annexe I). Pas d'engagement chiffré pour Chine, Inde, Brésil, etc. \textbf{Paris :} Approche « bas-haut » (bottom-up). \emph{Tous} les pays déposent une CDN (différenciée mais universelle). Mécanisme de réhaussement régulier.
    \item \textbf{CDN Cameroun :} Inconditionnel \textasciitilde \SI{12}{\%} / Conditionnel \textasciitilde \SI{20}{\%} (total \SI{32}{\%} vs BAU 2035).
    \begin{itemize}
        \item \textbf{Forêt :} (1) Lutte contre la déforestation illégale / certification FSC. (2) Reboisement / Restauration paysages dégradés (AFR100, Grande Muraille Verte). (3) REDD+ (paiement pour services écosystémiques).
        \item \textbf{Énergie :} (1) Développement hydroélectricité (Nachtigal, Memve'ele, Grand Eweng). (2) Solaire (centrales Maroua, Guider, kits solaires ruraux). (3) Efficacité énergétique (ampoules LED, foyers améliorés).
    \end{itemize}
    \item \textbf{RCPD :} Responsabilité commune (tous concernés) mais différenciée (pays développés = émetteurs historiques + capacités financières). \textbf{Conséquence financement :} Obligation pour pays développés de mobiliser \$100 Md/an (Fonds Vert pour le Climat) pour aider pays en développement (atténuation + adaptation). Nouveau but collectif quantifié (NCQG) post-2025 en discussion.
    \item \textbf{Correction :} \textbf{Faux.} Protocole de \textbf{Montréal (1987)} = protection de la \textbf{couche d'ozone} (interdiction CFC). Problème : trou d'ozone $\rightarrow$ UV cancérigènes. \textbf{Accord de Paris (2015)} = lutte contre \textbf{l'effet de serre / réchauffement} (réduction $\ce{CO2}$, $\ce{CH4}$, $\ce{N2O}$). Les CFC sont aussi des GES puissants, donc Montréal a \emph{co-bénéfice} climat, mais ce n'est pas son objectif premier ni suffisant pour limiter le réchauffement à $1,5\,^\circ\mathrm{C}$.
\end{enumerate}
\end{corrige}

\newpage

\coursec{Activité d'intégration}

\coursec{Lutte contre l'effet de serre}

\courssub{Objectifs de l'activité}
Cette activité d'intégration vise à mobiliser l'ensemble des savoir-faire de la leçon « Lutte contre l'effet de serre » à travers une démarche d'investigation complète. À l'issue de ce travail, vous devez être capable de :
\begin{itemize}
    \item \textbf{Expliquer} le mécanisme de l'effet de serre naturel et son renforcement anthropique à partir d'un schéma annoté.
    \item \textbf{Identifier} et \textbf{quantifier} les sources d'émissions de gaz à effet de serre (GES) liées à des activités humaines concrètes dans un contexte camerounais.
    \item \textbf{Relier} chaque source d'émission à l'activité socio-économique qui la génère (transport, énergie, déchets, agriculture, déforestation).
    \item \textbf{Proposer} et \textbf{hiérarchiser} des mesures d'atténuation (réduction des émissions) et d'adaptation (réduction de la vulnérabilité) pertinentes, chiffrées et faisables localement.
    \item \textbf{Communiquer} vos résultats sous forme d'une affiche synthétique (diagnostic + plan d'action) et \textbf{débattre} de la faisabilité et de la justice climatique de ces mesures.
\end{itemize}

\courssub{Situation : Diagnostic carbone du quartier « Mokolo-Nord » à Yaoundé}
\label{sit:mokolo}
Le quartier \emph{Mokolo-Nord} est un quartier périphérique dense de Yaoundé, comptant environ \textbf{5 000 habitants} répartis en \textbf{1 000 ménages}. Comme de nombreux quartiers en forte croissance au Cameroun, il fait face à des défis énergétiques, de transport et de gestion des déchets qui contribuent au renforcement de l'effet de serre.

Une association de jeunes du quartier, « \emph{Vert Avenir} », souhaite réaliser un \textbf{diagnostic carbone simplifié} pour l'année 2023 et élaborer un \textbf{Plan Climat Quartier (PCQ)} à horizon 2030. Ils collectent les données suivantes auprès des services municipaux, de la SONARA (pétrole), de l'ANAFOR (forêts) et par enquête de terrain :

\begin{center}
\begin{tabularx}{\linewidth}{|l|X|}\hline
\textbf{Source d'émission / Puits} & \textbf{Données collectées (Année 2023)} \\ \hline
\textbf{Transport : Moto-taxis (\emph{benskins})} & 300 moto-taxis en activité quotidienne. Consommation moyenne : 5 L d'essence / jour / moto. Fonctionnement : 300 jours/an. \\ \hline
\textbf{Énergie : Groupes électrogènes (ménages \& commerces)} & 200 groupes électrogènes (essence). Moyenne : 10 L/jour/groupe pendant 150 jours/an (coupures fréquentes). \\ \hline
\textbf{Cuisson : Bois de chauffe \& charbon de bois} & 800 ménages utilisent le bois/charbon. Consommation moyenne : 1,5 tonne de bois équivalent / an / ménage. \textit{Note : 1 tonne de bois sec brûlé $\approx$ 1,8 t \ce{CO2} émis (partie non renouvelable). On considère que 50\% du bois prélevé n'est pas renouvelé (déforestation nette).} \\ \hline
\textbf{Déchets : Décharge sauvage à ciel ouvert} & Volume estimé de déchets organiques mis en décharge : 1 200 tonnes/an. Décomposition anaérobie $\rightarrow$ émission de \ce{CH4}. Facteur d'émission simplifié : 0,5 t éq. \ce{CO2} / tonne de déchets organiques. \\ \hline
\textbf{Déforestation : Extension urbaine \& agriculture de brûlis} & 200 arbres adultes (essences locales, 30 ans) abattus dans l'année pour extension de parcelles et bois de chauffe non géré. \\ \hline
\textbf{Puits de carbone existant : Arbres urbains \& reboisement} & 500 arbres adultes présents dans le quartier (absorbent ~20 kg \ce{CO2}/an/chacun). Projet de plantation de 1 000 jeunes plants (taux de survie estimé 70\%, absorption effective dans 5 ans). \\ \hline
\end{tabularx}
\end{center}

\textbf{Coefficients d'émission simplifiés fournis (valeurs moyennes pour le calcul) :}
\begin{itemize}
    \item Essence (transport \& groupes) : \cle{2,3 kg \ce{CO2} / L} (combustion complète).
    \item Arbre adulte abattu (perte de stock + décomposition) : \cle{25 kg \ce{CO2} / arbre} (émission nette instantanée).
    \item Absorption par arbre adulte (photosynthèse nette) : \cle{20 kg \ce{CO2} / an / arbre}.
    \item Potentiel de réchauffement global (PRG) du \ce{CH4} sur 100 ans : 28 (déjà intégré dans le facteur déchets).
\end{itemize}

\courssub{Consignes}
Travail en groupes de 4 à 5 élèves. Restitution sous forme d'une \textbf{affiche A1} (ou 2 A3 assemblées) et présentation orale de 5 min + 5 min de débat.

\begin{enumerate}
    \item \textbf{Modélisation de l'effet de serre (Rappel) :} Sur l'affiche, dessinez un schéma annoté du mécanisme de l'effet de serre (rayonnement solaire entrant, rayonnement infrarouge sortant, absorption/réémission par les GES : \ce{CO2}, \ce{CH4}, \ce{H2O}). Indiquez par une flèche rouge le « renforcement » dû aux émissions du quartier.
    \item \textbf{Inventaire des émissions (Calculs) :} Calculez les émissions annuelles de \ce{CO2} équivalent (en tonnes éq. \ce{CO2}/an) pour chacune des 5 sources identifiées (Transport, Groupes, Cuisson, Déchets, Déforestation). Calculez l'absorption annuelle actuelle par les 500 arbres existants. Déduisez le \textbf{bilan carbone net} du quartier (Émissions $-$ Absorption).
    \item \textbf{Analyse graphique :} Construisez un \textbf{diagramme en barres} (échelle en t éq. \ce{CO2}/an) classant les 5 sources par ordre décroissant. Ajoutez une barre « Absorption » (valeur négative) pour visualiser le déficit.
    \item \textbf{Élaboration du Plan Climat Quartier (PCQ) :} Proposez \textbf{5 actions d'atténuation} (réduction émissions / augmentation puits) et \textbf{2 actions d'adaptation}. Pour chaque action, remplissez le tableau ci-dessous sur l'affiche :
    \begin{center}
    \begin{tabularx}{\linewidth}{|l|X|X|X|X|}\hline
    \textbf{Action} & \textbf{Type (Atténuation / Adaptation)} & \textbf{Gain \ce{CO2} estimé (t/an) ou Résilience} & \textbf{Faisabilité locale (Coût, Acceptabilité, Délai)} & \textbf{Justification chiffrée (lien avec données)} \\ \hline
    \end{tabularx}
    \end{center}
    \item \textbf{Débat contradictoire :} En classe, défendez votre hiérarchisation. Argument : « Faut-il prioriser l'interdiction des \emph{benskins} (emploi jeunes) ou le reboisement (temps long) ? » « Qui finance les groupes solaires ? »
\end{enumerate}

\courssub{Ressources à mobiliser}
\begin{itemize}
    \item \textbf{Mécanisme de l'effet de serre :} Rayonnement solaire (ondes courtes) traverse l'atmosphère $\rightarrow$ Sol chauffe $\rightarrow$ Émet rayonnement infrarouge (ondes longues) $\rightarrow$ GES (\ce{CO2}, \ce{CH4}, \ce{H2O}, \ce{N2O}) absorbent et réémettent vers le sol $\rightarrow$ Réchauffement surface. Renforcement = augmentation [GES] $\rightarrow$ plus d'IR piégés.
    \item \textbf{Formule de calcul d'émission :} $E = A \times FE$ \\ $E$ : Émission (kg ou t \ce{CO2}), $A$ : Activité (L, km, t, nombre), $FE$ : Facteur d'émission (kg \ce{CO2}/unité d'activité).
    \item \textbf{Bilan carbone :} $Bilan = \sum \text{Émissions sources} - \sum \text{Absorption puits}$.
    \item \textbf{Principales sources anthropiques (GIEC) :} Énergie (combustion fossile), Transport, Agriculture/Foresterie (déforestation, élevage), Déchets, Procédés industriels.
    \item \textbf{Atténuation vs Adaptation :} Atténuation = agir sur les causes (réduire GES). Adaptation = agir sur les conséquences (réduire vulnérabilité).
\end{itemize}

\courssub{Démarche et corrigé commenté}
\begin{exoresolu}[Corrigé type de l'activité d'intégration — Diagnostic carbone Mokolo-Nord]
\textbf{1. Schéma de l'effet de serre (Consigne 1)}
\emph{Attendu :} Schéma montrant : Soleil $\xrightarrow{\text{UV/Visible}}$ Atmosphère (traversée) $\xrightarrow{\text{IR}}$ Sol $\xrightarrow{\text{IR}}$ Atmosphère (absorption par \ce{CO2}, \ce{CH4}, \ce{H2O}) $\xrightarrow{\text{Réémission vers le bas}}$ Réchauffement. Flèche rouge « Activités humaines (Mokolo-Nord) $\rightarrow$ + \ce{CO2}, + \ce{CH4} $\rightarrow$ Effet de serre renforcé $\rightarrow$ + Températures ».

\begin{popfigure}
\begin{tikzpicture}[>=Stealth, scale=0.9, every node/.style={font=\small}]
% Sun
\node[draw=poporange, fill=poporange!20, circle, minimum size=1.5cm] (sun) at (0,4) {\textbf{Soleil}};
% Incoming radiation
\draw[poporange, very thick, ->] (sun) -- ++(0,-1.5) node[midway, right] {Rayonnement solaire (ondes courtes)};
% Atmosphere top
\draw[popblue!50, fill=popblue!10, rounded corners] (-4,2) rectangle (4,1.5);
\node at (0,1.75) {Atmosphère (gaz traces : \ce{CO2}, \ce{CH4}, \ce{H2O}, \ce{N2O})};
% Surface
\draw[popdark, fill=popgreen!20, thick] (-4,0) rectangle (4,-0.5);
\node at (0,-0.25) {Surface terrestre (sol, océans, végétation)};
% Outgoing IR from surface
\draw[popred, very thick, ->] (0,-0.5) -- (0,-1.2) node[midway, right] {Rayonnement infrarouge (IR)};
% Absorption/Re-emission
\draw[popred, thick, ->] (0,-1.2) -- (-2.5,0.5) node[midway, left, align=center] {Absorption par les GES};
\draw[popred, thick, ->] (-2.5,0.5) -- (0,-0.5) node[midway, right, align=center] {Réémission vers le bas};
\draw[popred, thick, ->] (0,-1.2) -- (2.5,0.5) node[midway, right, align=center] {Absorption par les GES};
\draw[popred, thick, ->] (2.5,0.5) -- (0,-0.5) node[midway, left, align=center] {Réémission vers le bas};
% Escape to space
\draw[popred, dashed, ->] (0,1.5) -- (0,2.5) node[above] {Partie s'échappe vers l'espace};
% Human reinforcement arrow
\draw[popred, very thick, ->, decorate, decoration={snake, amplitude=1pt, segment length=4pt}] (-5,1) -- (-5,0) node[midway, left, align=center, text width=2.5cm] {\textbf{Renforcement anthropique}\\(+ \ce{CO2}, + \ce{CH4} $\rightarrow$ + IR piégés)};
\end{tikzpicture}
\legende{Mécanisme de l'effet de serre naturel et son renforcement par les activités humaines (quartier Mokolo-Nord). Les gaz à effet de serre (GES) absorbent le rayonnement infrarouge émis par la Terre et le réémettent vers le sol, augmentant la température de surface. L'augmentation de leur concentration atmosphérique renforce ce piégeage.}
\end{popfigure}

\textbf{Attention : Piège fréquent}
Ne pas confondre « trou dans la couche d'ozone » (CFC, UV) et « effet de serre » (GES, IR). Ne pas dire « le \ce{CO2} détruit l'ozone ».

\textbf{2. Calculs détaillés des émissions (Consigne 2)}
\emph{Unités :} Tout convertir en \textbf{tonnes éq. \ce{CO2}/an} (1 t = 1 000 kg).

\begin{itemize}
    \item \textbf{A. Transport (Moto-taxis) :}
    \begin{align*}
    \text{Conso totale} &= 300 \text{ motos} \times 5 \text{ L/j} \times 300 \text{ j/an} = 450\,000 \text{ L/an} \\
    E_{\text{transport}} &= 450\,000 \text{ L} \times 2,3 \text{ kg \ce{CO2}/L} = 1\,035\,000 \text{ kg} = \mathbf{1\,035 \text{ t éq. \ce{CO2}/an}}
    \end{align*}

    \item \textbf{B. Groupes électrogènes :}
    \begin{align*}
    \text{Conso totale} &= 200 \text{ groupes} \times 10 \text{ L/j} \times 150 \text{ j/an} = 300\,000 \text{ L/an} \\
    E_{\text{groupes}} &= 300\,000 \times 2,3 = 690\,000 \text{ kg} = \mathbf{690 \text{ t éq. \ce{CO2}/an}}
    \end{align*}

    \item \textbf{C. Cuisson (Bois non renouvelé) :}
    \begin{align*}
    \text{Bois total} &= 800 \text{ ménages} \times 1,5 \text{ t/an} = 1\,200 \text{ t bois/an} \\
    \text{Part non renouvelée (50\%)} &= 600 \text{ t bois/an} \\
    \text{Émission} &= 600 \text{ t bois} \times 1,8 \text{ t \ce{CO2}/t bois} = \mathbf{1\,080 \text{ t éq. \ce{CO2}/an}}
    \end{align*}
    \emph{Note :} Si le bois est renouvelé (plantation gérée), le bilan est neutre. Ici, déforestation nette $\rightarrow$ émission nette.

    \item \textbf{D. Déchets (Décharge sauvage) :}
    \begin{align*}
    E_{\text{déchets}} &= 1\,200 \text{ t déchets} \times 0,5 \text{ t éq. \ce{CO2}/t} = \mathbf{600 \text{ t éq. \ce{CO2}/an}}
    \end{align*}

    \item \textbf{E. Déforestation (Abattage 200 arbres) :}
    \begin{align*}
    E_{\text{déforest}} &= 200 \text{ arbres} \times 25 \text{ kg \ce{CO2}/arbre} = 5\,000 \text{ kg} = \mathbf{5 \text{ t éq. \ce{CO2}/an}}
    \end{align*}
    \emph{Remarque :} Ce chiffre semble faible comparé aux combustions. Il ne compte que la perte de stock immédiat (biomasse + sol). La perte de séquestration future est comptée dans l'absorption manquante.

    \item \textbf{F. Absorption actuelle (Puits) :}
    \begin{align*}
    A_{\text{actuel}} &= 500 \text{ arbres} \times 20 \text{ kg \ce{CO2}/an} = 10\,000 \text{ kg} = \mathbf{10 \text{ t éq. \ce{CO2}/an}}
    \end{align*}
\end{itemize}

\textbf{Bilan Carbone Net 2023 :}
\begin{align*}
\text{Total Émissions} &= 1\,035 + 690 + 1\,080 + 600 + 5 = \mathbf{3\,410 \text{ t éq. \ce{CO2}/an}} \\
\text{Total Absorption} &= \mathbf{10 \text{ t éq. \ce{CO2}/an}} \\
\text{Bilan Net} &= 3\,410 - 10 = \mathbf{+3\,400 \text{ t éq. \ce{CO2}/an}} \quad (\text{Source nette massive})
\end{align*}
\emph{Interprétation :} Le quartier émet 340 fois plus qu'il n'absorbe. L'urgence est à la réduction des sources (Transport, Cuisson, Énergie, Déchets).

\textbf{3. Diagramme en barres (Consigne 3)}
\begin{popfigure}
\begin{tikzpicture}
\begin{axis}[
    ybar,
    bar width=12pt,
    width=\linewidth,
    height=9cm,
    ylabel={Émissions (t éq. \ce{CO2}/an)},
    symbolic x coords={Cuisson, Transport, Groupes, Déchets, Déforestation, Absorption},
    xtick=data,
    nodes near coords,
    nodes near coords align={vertical},
    ymin=-50, ymax=1200,
    x tick label style={rotate=30, anchor=east, font=\small},
    ymajorgrids=true,
    grid style=dashed,
    title={Diagnostic Carbone Mokolo-Nord 2023},
    title style={font=\bfseries\small},
    every node near coord/.append style={font=\tiny, rotate=90, anchor=west, yshift=2pt},
    legend style={at={(0.5,-0.2)}, anchor=north, legend columns=-1}
]
\addplot[popblue!80, fill=popblue!50] coordinates {(Cuisson,1080) (Transport,1035) (Groupes,690) (Déchets,600) (Déforestation,5)};
\addplot[popgreen!80, fill=popgreen!50] coordinates {(Absorption,-10)};
\legend{Émissions, Absorption}
\end{axis}
\end{tikzpicture}
\legende{Diagramme en barres des émissions annuelles de GES par source (t éq. \ce{CO2}/an) et absorption actuelle. La cuisson (bois non renouvelé) et le transport (moto-taxis) sont les deux postes majeurs. L'absorption (barre verte négative) est négligeable face aux émissions.}
\end{popfigure}

\textbf{4. Plan Climat Quartier (PCQ) — Proposition hiérarchisée (Consigne 4)}

\begin{center}
\small
\begin{tabularx}{\linewidth}{|l|l|X|X|X|}\hline
\textbf{Action} & \textbf{Type} & \textbf{Gain \ce{CO2} estimé / Résilience} & \textbf{Faisabilité locale} & \textbf{Justification chiffrée} \\ \hline
1. Programme « Cuisinières améliorées \& Biodigesteurs » (500 foyers ciblés) & Atténuation & \textbf{-540 t/an} (Réduction 50\% bois non renouv.) & \textbf{Haute.} Coût modéré (subventions MINEE/ONG), acceptabilité forte (économie bois, santé), délai 1-2 ans. & 400 foyers $\times$ 0,75 t bois économisé $\times$ 1,8 t \ce{CO2}/t = 540 t. Réduit déforestation \& pollution intérieure. \\ \hline
2. Solaire photovoltaïque collectif (Centrales de quartier + kits individuels) pour remplacer groupes électrogènes & Atténuation & \textbf{-550 t/an} (Remplacement 80\% conso groupes) & \textbf{Moyenne.} Investissement initial élevé (partenariat public-privé, Fonds Vert), maintenance technique nécessaire, délai 3-5 ans. & 300 000 L/an $\times$ 0,8 $\times$ 2,3 kg = 552 t. Supprime bruit \& pollution locale. \\ \hline
3. Gestion intégrée des déchets : Tri à la source + Compostage + Collecte méthanisation (biogaz) & Atténuation & \textbf{-480 t/an} (Évitement 80\% \ce{CH4} décharge) & \textbf{Moyenne.} Nécessite organisation communautaire, formation, infrastructure compostage. Délai 2 ans. & 1 200 t $\times$ 0,8 $\times$ 0,5 t = 480 t. Produit engrais (agriculture urbaine) \& biogaz (cuisson). \\ \hline
4. Reboisement urbain \& Agroforesterie : Plantation 2 000 arbres (taux survie 80\%) + Protection 500 existants & Atténuation (Puits) & \textbf{+32 t/an} (à maturité 10-15 ans) + \textbf{Stock} & \textbf{Haute.} Faible coût (pépinières locales, main d'œuvre jeunes), forte appropriation, délai court plantation, long séquestration. & 1 600 arbres $\times$ 20 kg = 32 t/an. Co-bénéfices : ombrage (adaptation), biodiversité, embellissement. \\ \hline
5. Mobilité douce \& Transport en commun : Pistes cyclables, marche, minibus électriques (remplacement progressif 50 benskins) & Atténuation & \textbf{-170 t/an} (Remplacement 50 motos $\times$ 5L $\times$ 300j $\times$ 2,3kg) & \textbf{Faible à Moyenne.} Infrastructure lourde (voirie), coût véhicules électriques, changement culturel. Délai 5-10 ans. & 50 motos $\times$ 5 $\times$ 300 $\times$ 2,3 = 172,5 t. Nécessite politique urbaine (CUD). \\ \hline
6. Végétalisation \& Gestion des eaux pluviales : Jardins de pluie, toits verts, corridors verts & Adaptation & \textbf{Résilience :} -Îlot de chaleur, -Inondations, +Sécurité alimentaire & \textbf{Haute.} Main d'œuvre locale, matériaux locaux, co-bénéfices immédiats (fraîcheur, légumes). & Répond aux canicules (Yaoundé +2°C projeté) et inondations récurrentes. Synergie Action 4. \\ \hline
7. Système d'alerte précoce \& Formation risques climatiques (Inondations, glissements) + Assurance indicielle communautaire & Adaptation & \textbf{Résilience :} Sauvegarde vies \& biens, réduction vulnérabilité pauvres & \textbf{Moyenne.} Nécessite coordination Mairie/Protection Civile/Météo Cameroun, financement assurance. & Quartier en zone pente/argile = risque glissement. Population jeune/non assurée. \\ \hline
\end{tabularx}
\end{center}

\textbf{5. Éléments pour le débat (Consigne 5)}
\begin{itemize}
    \item \textbf{Justice sociale :} Interdire les \emph{benskins} sans alternative = perte de revenus pour 300 jeunes + familles. L'action 5 doit être \emph{accompagnée} (reconversion, minibus subventionnés).
    \item \textbf{Financement :} Actions 2 \& 5 nécessitent fonds externes (Fonds Vert, Banque Mondiale, Budget d'Investissement Public). Actions 1, 3, 4, 6 peuvent démarrer par auto-financement communautaire / micro-crédit.
    \item \textbf{Efficacité vs Rapidité :} Cuisinières (Action 1) = gain immédiat, faible coût. Reboisement (Action 4) = gain différé (10 ans), mais indispensable pour neutralité 2050.
    \item \textbf{Adaptation prioritaire ?} Au Cameroun, l'adaptation (Actions 6, 7) est souvent plus urgente localement que l'atténuation (faible émetteur mondial), mais l'atténuation locale améliore la qualité de vie (air, bruit, déchets).
\end{itemize}
\tcblower
\textbf{Commentaires pédagogiques pour l'enseignant :}
\begin{itemize}
    \item \textbf{Évaluation des compétences :} Vérifier que l'élève distingue bien \emph{émission évitée} (atténuation source) et \emph{séquestration} (puits). Vérifier l'unité (t éq. \ce{CO2}/an) et les ordres de grandeur.
    \item \textbf{Erreurs fréquentes :} Oublier le facteur « non renouvelé » pour le bois (compter 100\% comme émission). Confondre kg et tonnes. Ne pas intégrer le PRG du \ce{CH4} pour les déchets (donnée fournie ici simplifiée). Proposer des actions « hors-sol » (ex: « fermer les usines » inexistantes au quartier).
    \item \textbf{Approfondissement :} Demander le calcul du « coût par tonne de \ce{CO2} évitée » pour chaque action (économie de la lutte). Introduire la notion de « \cle{co-bénéfices} » (santé, biodiversité, emploi) qui justifie l'action climat au-delà du carbone.
    \item \textbf{Contexte Cameroun :} Relier à la \cle{Contribution Déterminée au Niveau National (CDN)} du Cameroun (Accord de Paris) : -35\% émissions à l'horizon 2030. Le PCQ est la déclinaison locale de cet engagement national.
\end{itemize}
\end{exoresolu}

\courssub{Bilan de l'activité}
Cette activité a permis de mettre en évidence les points majeurs suivants :
\begin{enumerate}
    \item \textbf{Quantification du problème :} Le quartier émet \textbf{3 410 t éq. \ce{CO2}/an} pour seulement \textbf{10 t/an} absorbées. Le bilan net est massivement positif (source), typique des zones urbaines en développement où la combustion fossile (transport, électricité de secours) et la biomasse non gérée (cuisson, déforestation) dominent.
    \item \textbf{Hiérarchisation des leviers :} Les deux tiers des émissions proviennent de la \textbf{cuisson (bois non renouvelé : 32\%)} et du \textbf{transport (moto-taxis : 30\%)}. Toute politique climat efficace doit s'attaquer en priorité à ces deux postes, sans négliger l'énergie de secours (groupes : 20\%) et les déchets (18\%).
    \item \textbf{Complémentarité Atténuation / Adaptation :} Les actions d'atténuation les plus efficaces localement (cuisinières, déchets, solaire) génèrent des \textbf{co-bénéfices immédiats d'adaptation} (santé respiratoire, réduction îlots de chaleur, fertilité sols, revenus). L'adaptation (végétalisation, alertes) n'est pas un « luxe » mais une nécessité vitale face aux inondations et canicules déjà observées à Yaoundé.
    \item \textbf{Faisabilité et Justice :} Les solutions « clés en main » importées (ex: interdiction pure des motos, voitures électriques individuelles) échouent par inadéquation socio-économique. Les solutions gagnantes sont \textbf{systémiques, participatives et progressives} : cuisinières locales, compostage communautaire, reboisement par les jeunes, mini-réseaux solaires partagés.
    \item \textbf{Échelle de l'action :} Le quartier est l'échelle pertinente pour l'expérimentation et la mobilisation citoyenne, mais l'atteinte des objectifs (ex: -35\% CDN Cameroun) nécessite une \textbf{intégration dans la planification urbaine (PDU, PLU)} et un \textbf{financement pérenne} (budget communal, fiscalité carbone, fonds climat).
\end{enumerate}
\begin{aretenir}[Message clé pour le Probatoire]
Lutter contre l'effet de serre au Cameroun, c'est \textbf{réduire la pauvreté énergétique} (accès électricité propre, cuisson propre), \textbf{gérer les ressources naturelles} (forêts, sols, déchets) et \textbf{adapter le territoire} (urbanisme, alertes). Chaque tonne de \ce{CO2} évitée localement améliore le cadre de vie \emph{ici et maintenant}.
\end{aretenir}

\newpage

\coursec{Méthodes et exercices résolus}

\coursec{Lutte contre l'effet de serre}

\courssub{L'effet de serre : un phénomène naturel indispensable}

\begin{definition}[Effet de serre naturel]
Phénomène physique par lequel certains gaz de l'atmosphère, dits \cle{gaz à effet de serre (GES)}, absorbent une partie du rayonnement infrarouge émis par la surface terrestre et le réémettent dans toutes les directions, y compris vers le sol. Ce piégeage radiatif maintient la température moyenne à la surface du globe à environ $+15\ ^\circ\text{C}$, contre $-18\ ^\circ\text{C}$ en l'absence d'atmosphère.
\end{definition}

\begin{aretenir}[Les principaux gaz à effet de serre naturels]
\begin{itemize}
\item \textbf{Vapeur d'eau (\ce{H2O}) :} premier GES par sa quantité et son effet (environ $60\ \%$ de l'effet de serre naturel), sa concentration varie rapidement (cycle de l'eau).
\item \textbf{Dioxyde de carbone (\ce{CO2}) :} second contributeur naturel, cycle long (siècles à millénaires), échangé entre atmosphère, océans, biosphère et géosphère.
\item \textbf{Méthane (\ce{CH4}) :} puissant GES (pouvoir de réchauffement global $\approx 28$ fois celui du \ce{CO2} sur 100 ans), émis par les zones humides naturelles, les termites, les hydrates de méthane.
\item \textbf{Protoxyde d'azote (\ce{N2O}) et gaz fluorés :} traces, mais pouvoir de réchauffement très élevé.
\end{itemize}
\end{aretenir}

\begin{popfigure}
\begin{tikzpicture}[scale=0.9, every node/.style={font=\small}]
% Sun
\node[circle, fill=poporange, minimum size=1.5cm] (sun) at (0,5) {};
\node[above=0.2cm of sun] {\textbf{Soleil}};
\node[right=0.2cm of sun] {\tiny Rayonnement ondes courtes (UV, visible)};

% Atmosphere top
\draw[popblue, thick] (-3,3.5) -- (3,3.5) node[midway, above=0.1cm] {Atmosphère (GES : \ce{CO2}, \ce{CH4}, \ce{H2O})};

% Incoming solar
\draw[->, poporange, thick] (0,4.2) -- (0,3.7) node[midway, right] {$100\%$};
\draw[->, poporange, thick] (0.5,3.7) -- (1.5,3.7) node[midway, above] {$30\%$ réfléchi (Albédo)};
\draw[->, poporange, thick] (0,3.7) -- (0,2.8) node[midway, right] {$70\%$ absorbé};

% Ground
\draw[pattern=north east lines, pattern color=popdark] (-3,2) rectangle (3,2.5) node[midway] {Surface terrestre (Sol, Océans)};

% IR emission from ground
\draw[->, popred, thick] (0,2.5) -- (0,3.2) node[midway, left] {Rayonnement infrarouge (ondes longues)};

% GHG absorption / re-emission
\draw[popblue, decorate, decoration={snake, amplitude=1pt, segment length=4pt}] (-1,3.2) -- (-1,3.5);
\draw[popblue, decorate, decoration={snake, amplitude=1pt, segment length=4pt}] (1,3.2) -- (1,3.5);
\node[popblue, align=center] at (-2,3.35) {Absorption\\par les GES};
\node[popblue, align=center] at (2,3.35) {Réémission\\isotrope};

% Back radiation to ground
\draw[->, popred, thick] (-1,3.5) -- (-1,2.5) node[midway, left] {Contre-rayonnement};
\draw[->, popred, thick] (1,3.5) -- (1,2.5) node[midway, right] {Contre-rayonnement};

% Outgoing to space
\draw[->, popred, thick] (0,3.5) -- (0,4.2) node[midway, right] {Vers l'espace};

% Temperature
\node[draw, rounded corners, fill=popcream, anchor=south] at (0,1.3) {\textbf{Bilan :} $T_{\text{moy}} \approx +15\ ^\circ\text{C}$ (au lieu de $-18\ ^\circ\text{C}$ sans GES)};
\end{tikzpicture}
\legende{Schéma du bilan radiatif terrestre et mécanisme de l'effet de serre naturel. Le rayonnement solaire (ondes courtes) traverse l'atmosphère et chauffe le sol. Le sol émet un rayonnement infrarouge (ondes longues) que les gaz à effet de serre absorbent et réémettent en partie vers le sol, réchauffant la basse atmosphère.}
\end{popfigure}

\begin{methode}[Décrire et expliquer l'effet de serre naturel à partir d'un schéma]
\begin{enumerate}
\item \textbf{Identifier les flux d'énergie} : rayonnement solaire incident (courtes longueurs d'onde, UV/visible), fraction réfléchie vers l'espace (albédo planétaire $\approx 30\%$), fraction absorbée par la surface et l'atmosphère ($\approx 70\%$).
\item \textbf{Nommer le rayonnement terrestre} : la surface chauffée réémet un rayonnement \cle{infrarouge} (grandes longueurs d'onde) vers le haut.
\item \textbf{Repérer le rôle des GES} (\ce{CO2}, \ce{CH4}, \ce{H2O}) : ils \cle{absorbent} une partie de ce rayonnement infrarouge et le \cle{réémettent} dans toutes les directions (isotropie), notamment vers la surface (\emph{contre-rayonnement}).
\item \textbf{Conclure sur le bilan thermique} : une fraction de l'énergie thermique est piégée dans la basse atmosphère (troposphère), maintenant une température moyenne compatible avec l'eau liquide et la vie ($+15\ ^\circ\text{C}$).
\item \textbf{Rédiger la synthèse} en utilisant le vocabulaire exigé : rayonnement solaire, rayonnement infrarouge, absorption, réémission, bilan radiatif, température moyenne.
\end{enumerate}
\textbf{Réflexes Probatoire :} Ne jamais écrire que les GES « bloquent les rayons du Soleil » (ils laissent passer le visible). Préciser que l'effet de serre \emph{naturel} est \textbf{indispensable à la vie} ; seul son \emph{renforcement} anthropique est problématique.
\end{methode}

\begin{exoresolu}[Expliquer un schéma du bilan radiatif]
\textbf{Énoncé.} Sur un schéma du bilan radiatif, on lit : sur $100$ unités d'énergie solaire au sommet de l'atmosphère, $30$ sont réfléchies et $70$ absorbées. La surface émet un rayonnement infrarouge largement absorbé par les GES puis réémis vers le sol.
\begin{enumerate}
\item Calcule l'albédo planétaire.
\item Explique en 4--6 lignes le mécanisme de l'effet de serre naturel.
\item Justifie : « Sans effet de serre naturel, la vie telle que nous la connaissons serait impossible. »
\end{enumerate}
\tcblower
\textbf{Solution.}
\begin{enumerate}
\item L'\cle{albédo} $A$ est la fraction réfléchie : $A = \frac{30}{100} = 0{,}30$ soit $30\ \%$.
\item Le Soleil émet un rayonnement \emph{court} (visible/UV) qui traverse l'atmosphère et chauffe le sol. Celui-ci réémet un rayonnement \emph{infrarouge} (long). Les GES (\ce{CO2}, \ce{CH4}, \ce{H2O}) absorbent une partie de cet infrarouge et le réémettent vers le bas. Ce \emph{contre-rayonnement} réchauffe la surface et la basse atmosphère : c'est l'effet de serre.
\item Sans GES, tout l'infrarouge partirait dans l'espace : $T_{\text{moy}} \approx -18\ ^\circ\text{C}$, eau gelée partout. Grâce à l'effet de serre naturel, $T_{\text{moy}} \approx +15\ ^\circ\text{C}$, eau liquide possible. C'est donc un phénomène \emph{bénéfique et nécessaire}.
\end{enumerate}
\end{exoresolu}

\courssub{Le renforcement de l'effet de serre par les activités humaines}

\begin{definition}[Forçage radiatif]
Différence entre l'énergie entrante et sortante au sommet de l'atmosphère due à une perturbation (augmentation des GES, changement d'albédo). Un forçage \emph{positif} réchauffe le climat. Depuis 1750, le forçage total anthropique est d'environ $+2{,}7\ \text{W}\cdot\text{m}^{-2}$ (GIEC AR6), dont $\approx +2{,}0\ \text{W}\cdot\text{m}^{-2}$ pour le \ce{CO2} seul.
\end{definition}

\begin{popfigure}
\begin{tikzpicture}[scale=0.85, every node/.style={font=\small}]
% Sources
\node[draw, rounded corners, fill=poporangeL, align=center] (bois) at (-5,3) {Bois de chauffe\\\& Charbon de bois\\(Extrême-Nord, ménages)};
\node[draw, rounded corners, fill=poporangeL, align=center] (elev) at (0,3) {Élevage bovin\\\& Riziculture inondée\\(Logone, Mayo-Danay)};
\node[draw, rounded corners, fill=poporangeL, align=center] (trans) at (5,3) {Transport (taxis,\\motos, camions)\\\& Groupes électrogènes\\(Douala, Yaoundé)};

% Gases
\node[draw, circle, fill=popgreenL, minimum size=1.2cm] (co2) at (-5,1) {\ce{CO2}};
\node[draw, circle, fill=popblueL, minimum size=1.2cm] (ch4) at (0,1) {\ce{CH4}};
\node[draw, circle, fill=poppurpleL, minimum size=1.2cm] (n2o) at (5,1) {\ce{N2O}, Gaz fluorés};

% Arrows
\draw[->, thick, poporange] (bois.south) -- (co2.north) node[midway, right, align=center] {Combustion,\\ Déforestation};
\draw[->, thick, poporange] (elev.south) -- (ch4.north) node[midway, right, align=center] {Fermentation\\ entérique, Anaérobiose};
\draw[->, thick, poporange] (trans.south) -- (n2o.north) node[midway, left, align=center] {Combustion fossile,\\ Engrais, Industrie};

% Atmosphere accumulation
\node[draw, rounded corners, fill=poppinkL, align=center, minimum width=6cm] (atm) at (0,-1) {Accumulation dans l'atmosphère\\\textbf{Renforcement de l'effet de serre}\\\textbf{Hausse température moyenne ($+1,1\ ^\circ\text{C}$ depuis 1850)}};

\draw[->, thick, popgreen] (co2.south) -- (atm.north west);
\draw[->, thick, popblue] (ch4.south) -- (atm.north);
\draw[->, thick, poppurple] (n2o.south) -- (atm.north east);

% Consequences
\node[draw, rounded corners, fill=popred!20, align=center, minimum width=6cm] (cons) at (0,-3) {Conséquences :\quad Désertification Nord \quad|\quad Montée eaux littoral \quad|\quad Paludisme hauts plateaux \quad|\quad Insécurité alimentaire};
\draw[->, thick, popred] (atm.south) -- (cons.north);
\end{tikzpicture}
\legende{Principales sources anthropiques de GES au Cameroun et chaîne causale vers le renforcement de l'effet de serre.}
\end{popfigure}

\begin{methode}[Relier une activité humaine à l'émission de GES]
\begin{enumerate}
\item \textbf{Identifier l'activité} : transport, déforestation/bois de chauffe, élevage ruminant, riziculture inondée, industrie, déchets, épandage d'engrais.
\item \textbf{Associer le(s) GES émis} et le \textbf{processus} :
\begin{itemize}
\item Combustion énergies fossiles (pétrole, gaz, charbon) + combustion bois $\rightarrow$ \ce{CO2} (oxydation du carbone).
\item Fermentation entérique (ruminants), rizières inondées, décharges $\rightarrow$ \ce{CH4} (méthanogenèse anaérobie).
\item Engrais azotés, combustion biomasse, industrie chimique $\rightarrow$ \ce{N2O} (nitrification/dénitrification).
\item Réfrigération, climatisation, industrie électronique $\rightarrow$ Gaz fluorés (HFC, PFC, \ce{SF6}).
\end{itemize}
\item \textbf{Expliquer la chaîne causale} : Activité $\rightarrow$ Processus physico-chimique/biologique $\rightarrow$ Émission GES $\rightarrow$ Augmentation concentration atmosphérique $\rightarrow$ Renforcement effet de serre.
\item \textbf{Noter les effets indirects} : Déforestation = émission \ce{CO2} \textbf{ET} perte d'un puits de carbone (photosynthèse).
\end{enumerate}
\textbf{Réflexe :} Citer \emph{systématiquement} la formule chimique \textbf{et} le nom du gaz + le processus (combustion, fermentation, décomposition anaérobie).
\end{methode}

\begin{exoresolu}[Analyse d'activités émettrices en zone soudano-sahélienne]
\textbf{Énoncé.} Dans l'Extrême-Nord : (a) bois de chauffe/charbon ; (b) élevage bovin extensif ; (c) riziculture inondée (plaine du Logone). Pour chaque activité : GES émis + mécanisme.
\tcblower
\textbf{Solution.}
\begin{enumerate}
\item[(a)] \textbf{Bois/Charbon.} Combustion du carbone organique $\rightarrow$ \ce{CO2} (oxydation). \textbf{Double effet :} émission directe + réduction du couvert forestier $\rightarrow$ moins de photosynthèse $\rightarrow$ moins de \ce{CO2} absorbé.
\item[(b)] \textbf{Élevage bovin.} Ruminants $\rightarrow$ fermentation entérique dans le rumen (micro-organismes anaérobies dégradent la cellulose) $\rightarrow$ émission de \ce{CH4} (éructations).
\item[(c)] \textbf{Riziculture inondée.} Sol saturé en eau $\rightarrow$ conditions anaérobies $\rightarrow$ décomposition matière organique par archées méthanogènes $\rightarrow$ production de \ce{CH4} $\rightarrow$ diffusion vers l'atmosphère.
\end{enumerate}
\textbf{Conclusion :} Ces pratiques locales émettent \ce{CO2} et \ce{CH4}, renforçant l'effet de serre global.
\end{exoresolu}

\begin{methode}[Calculer une émission de \ce{CO2} (compétence quantitative)]
\begin{enumerate}
\item \textbf{Repérer :} Quantité consommée ($Q$) + Facteur d'émission ($FE$, en kg \ce{CO2}/unité).
\item \textbf{Appliquer :} $E = Q \times FE$.
\item \textbf{Convertir :} $1\ \text{t} = 1000\ \text{kg}$. Vérifier l'homogénéité des unités.
\item \textbf{Interpréter :} Comparer à une moyenne nationale, un objectif (ex. Accord de Paris), ou calculer une réduction potentielle.
\end{enumerate}
\textbf{Ordres de grandeur (Probatoire) :} $1\ \text{L}$ essence/gasoil $\approx 2{,}3$--$2{,}7\ \text{kg}\ \ce{CO2}$ ; $1\ \text{kg}$ bois sec $\approx 1{,}5$--$1{,}9\ \text{kg}\ \ce{CO2}$ ; Émission moyenne Cameroun $\approx 0{,}35\ \text{t}\ \ce{CO2}$/hab/an (vs $>4\ \text{t}$ Europe/USA).
\end{methode}

\begin{exoresolu}[Émissions d'un parc de taxis à Douala]
\textbf{Énoncé.} Un taxi consomme $12\ \text{L}$ gasoil/jour. $1\ \text{L}$ gasoil $\rightarrow 2{,}7\ \text{kg}\ \ce{CO2}$. Actif $300$ j/an. Ville : $5000$ taxis. Projet : remplacer $20\%$ par véhicules propres ($0\ \text{kg}\ \ce{CO2}$).
\begin{enumerate}
\item Émission annuelle d'un taxi (kg et t).
\item Émission totale du parc (t/an).
\item \ce{CO2} évité par an grâce au projet (t).
\end{enumerate}
\tcblower
\textbf{Solution.}
\begin{enumerate}
\item Conso/an = $12 \times 300 = 3600\ \text{L}$. $E = 3600 \times 2{,}7 = 9720\ \text{kg} = \mathbf{9{,}72\ t\ \ce{CO2}}$.
\item Parc total : $9{,}72 \times 5000 = \mathbf{48\,600\ t\ \ce{CO2}/an}$.
\item Évité : $20\% \times 48\,600 = \mathbf{9\,720\ t\ \ce{CO2}/an}$ (équivalent à retirer $1000$ taxis).
\end{enumerate}
\textbf{Analyse :} Ce calcul quantifie l'enjeu de la transition énergétique dans le transport urbain.
\end{exoresolu}

\courssub{Les conséquences du réchauffement climatique}

\begin{popfigure}
\begin{tikzpicture}[scale=0.8, every node/.style={font=\small, align=center}]
% Driver
\node[draw, rounded corners, fill=popred!30, align=center] (driver) at (0,4){Renforcement effet de serre\\\textbf{Hausse T° moyenne ($+1,1\ ^\circ\text{C}$)}};

% Biophysical
\node[draw, rounded corners, fill=popblueL, minimum width=3.5cm, align=center] (bio1) at (-5,2){Fonte glaces \& calottes\\\& Dilatation thermique océans};
\node[draw, rounded corners, fill=popblueL, minimum width=3.5cm, align=center] (bio2) at (0,2){Perturbation cycle de l'eau\\\& Événements extrêmes};
\node[draw, rounded corners, fill=popblueL, minimum width=3.5cm, align=center] (bio3) at (5,2){Acidification océans\\(\ce{CO2} dissous $\rightarrow$ \ce{H2CO3})};

% Impacts Env
\node[draw, rounded corners, fill=popgreenL, minimum width=3.5cm, align=center] (env1) at (-5,0){Élévation niveau des mers\\$\rightarrow$ Submersion côtière};
\node[draw, rounded corners, fill=popgreenL, minimum width=3.5cm, align=center] (env2) at (0,0){Sécheresses / Inondations\\\textit{Désertification (Nord)}};
\node[draw, rounded corners, fill=popgreenL, minimum width=3.5cm, align=center] (env3) at (5,0){Recul biodiversité\\Blanchiment coraux};

% Impacts Socio
\node[draw, rounded corners, fill=poporangeL, minimum width=3.5cm, align=center] (soc1) at (-5,-2){Menace zones côtières\\(Douala, Kribi, Limbé)};
\node[draw, rounded corners, fill=poporangeL, minimum width=3.5cm, align=center] (soc2) at (0,-2){Baisse rendements (maïs, cacao,\\mil, sorgho) $\rightarrow$ Insécurité alimentaire};
\node[draw, rounded corners, fill=poporangeL, minimum width=3.5cm, align=center] (soc3) at (5,-2){Extension paludisme (hauts plateaux)\\Migrations climatiques \\Conflits eau/terres};

% Arrows
\foreach \src/\dst in {driver/bio1, driver/bio2, driver/bio3, bio1/env1, bio2/env2, bio3/env3, env1/soc1, env2/soc2, env3/soc3}
  \draw[->, thick, popdark] (\src) -- (\dst);
\end{tikzpicture}
\legende{Chaîne causale : du forçage radiatif aux impacts socio-environnementaux au Cameroun.}
\end{popfigure}

\begin{methode}[Analyser les conséquences (environnement \& populations)]
\begin{enumerate}
\item \textbf{Classer} : Environnement (niveau mers, désertification, cycle eau, biodiversité, acidification océans) vs Populations (sécurité alimentaire, santé/paludisme, migrations, conflits ressources, infrastructures).
\item \textbf{Construire la chaîne causale complète} : Émissions GES $\rightarrow$ Forçage radiatif $\rightarrow$ $\Delta T$ $\rightarrow$ Processus biophysique $\rightarrow$ Impact localisé. \emph{Ne sauter aucune étape.}
\item \textbf{Contextualiser Cameroun} :
\begin{itemize}
\item \textbf{Nord (Extrême-Nord, Nord) :} Avancée désertique, régression Lac Tchad ($-90\%$ surface depuis 1960), raccourcissement saison pluvieuse, baisse mil/sorgho, insécurité alimentaire, migrations Sud, conflits agriculteurs/éleveurs.
\item \textbf{Littoral (Douala, Kribi, Limbé, Tiko) :} Montée eaux + subsidence + urbanisation $\rightarrow$ inondations récurrentes, érosion côtière, salinisation nappes, menace mangroves/ports.
\item \textbf{Hauts Plateaux (Ouest, Nord-Ouest) :} Remontée isothermes $\rightarrow$ extension zone de transmission paludisme vers altitudes $>1500\ \text{m}$, pression sur cultures (café, arabica).
\end{itemize}
\item \textbf{Appuyer sur données} : $+1{,}1\ ^\circ\text{C}$ global (GIEC), $+1{,}5\ ^\circ\text{C}$ seuil critique, tendance pluviométrique locale.
\end{enumerate}
\textbf{Réflexe :} Distinguer \emph{météo} (événement court) et \emph{climat} (moyenne 30 ans). Une canicule $\neq$ preuve ; la tendance décennale = preuve.
\end{methode}

\begin{exoresolu}[Analyse document : Extrême-Nord, 50 ans de données]
\textbf{Énoncé.} Hausse T° moyenne, baisse durée saison pluvieuse, surfaces mil/sorgho en baisse, Lac Tchad $-90\%$, migrations Sud.
\begin{enumerate}
\item Chaîne causale : Émissions mondiales $\rightarrow$ Baisse rendements mil/sorgho.
\item Deux conséquences sociales déductibles.
\end{enumerate}
\tcblower
\textbf{Solution.}
\begin{enumerate}
\item Émissions GES mondiales $\rightarrow$ Renforcement effet de serre $\rightarrow$ Hausse T° globale $\rightarrow$ Perturbation circulation atmosphérique sahélienne $\rightarrow$ Raccourcissement saison pluvieuse + Hausse évapotranspiration $\rightarrow$ Assèchement sols + Régression Lac Tchad $\rightarrow$ Stress hydrique cultures pluviales (mil, sorgho) $\rightarrow$ Baisse rendements.
\item \textbf{1. Insécurité alimentaire \& pauvreté :} Baisse production céréalière = moins de nourriture/revenus pour populations rurales.
\item \textbf{2. Migrations climatiques \& conflits :} Déplacements vers zones plus humides (Sud, Adamawa) $\rightarrow$ pression foncière, tensions agriculteurs/éleveurs, perte identité culturelle.
\end{enumerate}
\end{exoresolu}

\courssub{Lutter contre l'effet de serre : atténuation et adaptation}

\begin{definition}[Atténuation (Mitigation)]
Actions visant à \textbf{réduire les causes} du changement climatique : diminution des émissions de GES (efficacité énergétique, énergies renouvelables, mobilité durable) et/ou augmentation des \cle{puits de carbone} (reboisement, agroforesterie, restauration sols, capture/stockage \ce{CO2}).
\end{definition}

\begin{definition}[Adaptation]
Actions visant à \textbf{réduire la vulnérabilité et les effets} du changement climatique : variétés culturales résistantes/tolérantes, gestion intégrée ressources en eau, aménagement bassins versants (terrasses, haies vives), systèmes d'alerte précoce, protection côtière (digues, restauration mangroves), urbanisation résiliente.
\end{definition}

\begin{popfigure}
\begin{tikzpicture}[scale=0.85, every node/.style={font=\small, align=center}]
\node[draw, rounded corners, fill=popgoldL, minimum width=12cm] (title) at (0,3.5) {\textbf{Stratégies complémentaires face au changement climatique}};

% Atténuation
\node[draw, rounded corners, fill=popgreenL, minimum width=5.5cm] (att) at (-3.5,1.5) {\textbf{ATTÉNUATION (Agir sur les causes)}};
\node[draw, rounded corners, fill=popgreen!30, minimum width=5cm, align=center] (att1) at (-5.5,0){Réduire émissions\\\textbf{Énergie :} Hydro (Lom Pangar, Nachtigal), Solaire, Biogaz\\\textbf{Transport :} Bus BRT, Covoiturage, Véhicules électriques\\\textbf{Habitat :} Foyers améliorés (bois $-30$ à $-50\%$), Gaz domestique\\\textbf{Industrie :} Efficacité énergétique, Ciment bas carbone};
\node[draw, rounded corners, fill=popgreen!30, minimum width=5cm, align=center] (att2) at (-1.5,0){Augmenter puits de carbone\\\textbf{Reboisement :} Grande Muraille Verte, Reboisement communautaire\\\textbf{Agroforesterie :} Arbres + cultures (cacao, café, anacarde)\\\textbf{Gestion sols :} Agriculture de conservation, Compostage\\\textbf{Zones humides :} Restauration mangroves (Wouri, Sanaga)};

% Adaptation
\node[draw, rounded corners, fill=popblueL, minimum width=5.5cm] (adap) at (3.5,1.5) {\textbf{ADAPTATION (Agir sur les effets)}};
\node[draw, rounded corners, fill=popblue!30, minimum width=5cm, align=center] (adap1) at (1.5,0){Agriculture \& Sécurité alimentaire\\\textbf{Variétés :} Mil/sorgho cycle court, Maïs tolérant sécheresse\\\textbf{Pratiques :} Zaï, Demi-lunes, Paillage, Irrigation goutte-à-goutte\\\textbf{Élevage :} Races locales rustiques, Fourrage stocké};
\node[draw, rounded corners, fill=popblue!30, minimum width=5cm, align=center] (adap2) at (5.5,0){Territoires \& Infrastructures\\\textbf{Eau :} Barrages de retenue, Forages, Gestion intégrée bassins\\\textbf{Sols :} Terrasses, Haies vives, Cordons pierreux (anti-érosion)\\\textbf{Côte :} Restauration mangroves, Recul urbanisation, Digues\\\textbf{Santé :} Moustiquaires imprégnées, Surveillance épidémiologique paludisme};

\draw[->, thick, popgreen] (att) -- (att1);
\draw[->, thick, popgreen] (att) -- (att2);
\draw[->, thick, popblue] (adap) -- (adap1);
\draw[->, thick, popblue] (adap) -- (adap2);

% Synergy
\node[draw, rounded corners, fill=poppinkL, align=center] at (0,-2.2) {\textbf{CO-BÉNÉFICES (Synergies)} : Agroforesterie = Puits carbone (Atténuation) + Fertilité sols/Ombrage (Adaptation)\\\textbf{Mangroves} = Stockage \ce{CO2} (Atténuation) + Protection côtière/Élevage poissons (Adaptation)};
\end{tikzpicture}
\legende{Classification des mesures d'atténuation et d'adaptation avec exemples camerounais et synergies (co-bénéfices).}
\end{popfigure}

\begin{methode}[Proposer des mesures d'atténuation et d'adaptation contextualisées]
\begin{enumerate}
\item \textbf{Distinguier clairement} : Atténuation = Causes (GES) ; Adaptation = Conséquences (Vulnérabilité).
\item \textbf{Cibler la source} identifiée dans l'énoncé :
\begin{itemize}
\item Bois de chauffe $\rightarrow$ Foyers améliorés, Gaz (GPL), Biogaz, Reboisement énergie.
\item Transport $\rightarrow$ Transport en commun (BRT Douala/Yaoundé), Vélo/marche, Véhicules électriques/hybrides, Entretien moteur.
\item Agriculture $\rightarrow$ Agroforesterie, Compostage (réduit \ce{CH4}/\ce{N2O} vs engrais chimiques), Riziculture non inondée (SRI).
\item Énergie $\rightarrow$ Hydroélectricité, Solaire photovoltaïque, Biomasse durable.
\end{itemize}
\item \textbf{Vérifier faisabilité locale} : Coût, disponibilité, acceptabilité (ex. foyers améliorés peu coûteux, adoption rapide ; GPL nécessite filière distribution).
\item \textbf{Justifier scientifiquement} : « Cette mesure réduit \ce{CO2} car... » ou « Cette mesure stocke du carbone car... » ou « Cette mesure réduit la vulnérabilité car... ».
\item \textbf{Hiérarchiser} : Geste individuel $\rightarrow$ Action collective/communale $\rightarrow$ Politique nationale $\rightarrow$ Accord international.
\end{enumerate}
\textbf{Réflexe :} Un reboisement = \textbf{Atténuation} (puits de carbone). Une digue = \textbf{Adaptation} (protection). Ne pas inverser !
\end{methode}

\begin{exoresolu}[Plan d'action communal : Bafoussam]
\textbf{Énoncé.} Commune de Bafoussam : bois de chauffe massif, vieux véhicules polluants, cultures sur versants boisés (érosion). Proposer 3 mesures atténuation + 1 adaptation, justifiées.
\tcblower
\textbf{Solution.}

\textbf{Mesures d'atténuation :}
\begin{enumerate}
\item \textbf{Foyers améliorés + Promotion GPL/Biogaz.} Combustion efficace $\rightarrow$ $-30$ à $-50\%$ bois $\rightarrow$ $- \ce{CO2}$ émis + Moins de coupe $\rightarrow$ Puits de carbone préservés.
\item \textbf{Transport en commun structuré (Bus/Taxis collectifs) + Contrôle technique strict.} Mutualisation déplacements $\rightarrow$ Moins véhicules-km $\rightarrow$ Moins carburant $\rightarrow$ $- \ce{CO2}$/habitant transporté. Retrait véhicules vétustes $\rightarrow$ $- \ce{CO2}$, $- \text{NO}_x$, $- \text{particules}$.
\item \textbf{Reboisement versants + Agroforesterie (arbres fruitiers/engrais verts).} Photosynthèse $\rightarrow$ Stockage carbone biomasse/sols (Puits) + Racines stabilisent sols (co-bénéfice adaptation).
\end{enumerate}

\textbf{Mesure d'adaptation :}
\textbf{Aménagement versants (terrasses, fascines, haies vives) + Variétés culturales à cycle court/résistantes.} Ne réduit pas émissions mais protège sols érosion (pluies extrêmes) et sécurise récoltes (sécheresse) : la population s'adapte au climat nouveau.
\end{exoresolu}

\courssub{Les politiques climatiques : du local à l'international}

\begin{aretenir}[Échelonnement de la gouvernance climatique]
\begin{itemize}
\item \textbf{International :} \textbf{CCNUCC} (Rio 1992) $\rightarrow$ \textbf{Protocole de Kyoto} (1997, engagements pays développés) $\rightarrow$ \textbf{Accord de Paris} (2015, \textbf{tous} pays, objectif $<2\ ^\circ\text{C}$, viser $1{,}5\ ^\circ\text{C}$). Mécanismes : CDN (Contributions Déterminées Nationales), Bilan mondial (Global Stocktake) tous les 5 ans, Financement climat ($100$ Md\$/an pays développés $\rightarrow$ pays en développement).
\item \textbf{National (Cameroun) :} \textbf{CDN actualisée (2021)} : Engagement $-35\%$ émissions d'ici 2030 (scénario conditionnel). Piliers : \textbf{Forêt} (reboisement, gestion durable, REDD+), \textbf{Énergie} (hydroélectricité, solaire, efficacité), \textbf{Agriculture} (agroforesterie, intrants optimisés), \textbf{Déchets} (valorisation, méthanisation). Demande d'appui financier/technologique (Fonds Vert, Mécanisme pour le Développement Propre).
\item \textbf{Local :} Plans Communaux de Développement (PCD) intégrant climat, Plans d'Action Climat-Énergie, Gestion déchets municipaux (compostage, tri), Sensibilisation, Foyers améliorés.
\end{itemize}
\end{aretenir}

\begin{savaistu}[Le Cameroun et la justice climatique]
Le Cameroun incarne le principe de \textbf{responsabilité commune mais différenciée (RCMD)} : il émet $\approx 0{,}35\ \text{t}\ \ce{CO2}$/hab/an (faible émetteur) mais figure parmi les pays \textbf{les plus vulnérables} (indice ND-GAIN). L'Accord de Paris reconnaît cette asymétrie : les pays industrialisés (responsables historiques de $>70\%$ des émissions cumulées depuis 1850) doivent \textbf{réduire plus vite} et \textbf{financer l'adaptation/atténuation} du Sud. Les barrages de \textbf{Lom Pangar, Nachtigal, Memve'ele} illustrent l'atténuation nationale : énergie bas carbone remplaçant le thermique (fioul/gasoil), mais posent des défis d'adaptation sociale (réinstallation) et environnementale (méthanisation réservoirs, biodiversité).
\end{savaistu}

\begin{methode}[Présenter une politique climatique multi-échelles]
\begin{enumerate}
\item \textbf{Identifier l'échelle} demandée (individuelle, locale, nationale, internationale).
\item \textbf{Citer les instruments juridiques/opérationnels} propres à chaque échelle (voir boîte « À retenir »).
\item \textbf{Expliquer la RCMD} : Équité historique + Capacités économiques différenciées. Le Cameroun demande financement (Fonds Vert, Fonds d'Adaptation) et transfert technologique (solaire, réseaux intelligents, variétés améliorées).
\item \textbf{Montrer la complémentarité} : L'international fixe le cadre/objectifs ; le national décline en lois/stratégies (CDN) ; le local met en œuvre (terrain) ; l'individuel change les comportements. Échec à une échelle = frein aux autres.
\end{enumerate}
\end{methode}

\begin{exoresolu}[Engagement international du Cameroun (Accord de Paris)]
\textbf{Énoncé.} CDN Cameroun : engagement réduction émissions (reboisement, efficacité énergétique, EnR). Demande soutien financier/technologique.
\begin{enumerate}
\item Objectif principal Accord de Paris.
\item Justification demande soutien Cameroun.
\item Rôle hydroélectricité (Lom Pangar, Nachtigal, Memve'ele) dans la lutte.
\end{enumerate}
\tcblower
\textbf{Solution.}
\begin{enumerate}
\item Contenir hausse T° moyenne \textbf{nettement en dessous de $+2\ ^\circ\text{C}$} vs ère préindustrielle, poursuivre efforts pour limiter à $+1{,}5\ ^\circ\text{C}$.
\item Cameroun = \textbf{Faible émetteur} ($<0{,}4\ \text{t}\ \ce{CO2}$/hab/an) mais \textbf{Forte vulnérabilité} (désertification, côte, santé, agriculture). Principe RCMD : pays développés (émetteurs historiques) ont obligation morale/juridique (Art. 9 Accord Paris) d'apporter appui financier ($100$ Md\$/an objectif collectif), technologique et capacitaire aux pays en développement pour action climat.
\item Hydroélectricité = Énergie \textbf{renouvelable}, \textbf{bas carbone} (pas combustion en fonctionnement). Remplace groupes électrogènes/centrales thermiques (fioul/gasoil) $\rightarrow$ Évite émissions \ce{CO2} massives. Alimente ménages/industries $\rightarrow$ Réduit pression bois de chauffe (déforestation). Levier majeur \textbf{atténuation} dans CDN Cameroun.
\end{enumerate}
\textbf{Conclusion :} Solidarité internationale indispensable : objectifs communs, efforts différenciés selon responsabilités et capacités.
\end{exoresolu}

\begin{attention}[Les 5 erreurs fatales au Probatoire]
\begin{enumerate}
\item \textbf{Confondre les rayonnements :} « Les GES bloquent les rayons du Soleil » \textbf{FAUX}. L'atmosphère est transparente au solaire (visible) ; les GES absorbent le \emph{rayonnement infrarouge émis par la Terre}. \textit{Préciser toujours : Solaire (courtes $\lambda$) vs Terrestre (longues $\lambda$).}
\item \textbf{Démoniser l'effet de serre :} L'effet de serre \emph{naturel} est \textbf{indispensable} ($+15\ ^\circ\text{C}$ vs $-18\ ^\circ\text{C}$). Le problème = \emph{renforcement anthropique}. \textit{Le correcteur attend cette nuance fondamentale.}
\item \textbf{Confondre Atténuation / Adaptation :} Atténuation = Causes (GES, puits) ; Adaptation = Conséquences (vulnérabilité). Classer « Reboisement » en Adaptation ou « Digue » en Atténuation = \textbf{0 point}. \textit{Mémoriser : Reboisement = Puits = Atténuation.}
\item \textbf{Oublier unités/conversions :} Émission = nombre + unité (kg ou t \ce{CO2}). $1\ \text{t} = 1000\ \text{kg}$. $L \times \text{kg/L} = \text{kg}$ (pas t direct). \textit{Vérifier toujours l'homogénéité.}
\item \textbf{Confondre Trou ozone / Effet de serre :} \textbf{Deux problèmes distincts.} Couche ozone (stratosphère) = filtre UV (trou = CFC). Effet de serre (troposphère) = piégeage chaleur (GES = \ce{CO2}, \ce{CH4}...). \textit{Jamais écrire : « Le trou dans la couche d'ozone réchauffe la planète ».}
\end{enumerate}
\end{attention}

\begin{exercice}[Entraînement Probatoire - QCM \& Rédaction]
\begin{enumerate}
\item \textbf{QCM.} L'albédo planétaire moyen est d'environ : A) $0{,}1$ B) $0{,}3$ C) $0{,}7$ D) $0{,}9$.
\item \textbf{QCM.} Le principal GES émis par la riziculture inondée est : A) \ce{CO2} B) \ce{CH4} C) \ce{N2O} D) \ce{H2O}.
\item \textbf{Calcul.} Une famille utilise $2\ \text{kg}$ de bois/jour pour cuisiner ($FE = 1{,}8\ \text{kg}\ \ce{CO2}/\text{kg}$ bois). Calculer l'émission annuelle en tonnes. Proposer une mesure d'atténuation et une d'adaptation pour cette famille.
\item \textbf{Rédaction (10 lignes).} En t'appuyant sur l'exemple du Lac Tchad, explique comment le renforcement de l'effet de serre peut entraîner des migrations climatiques et des conflits d'usage de l'eau.
\item \textbf{Argumentaire.} « Le Cameroun pollue peu, il n'a pas à réduire ses émissions. » Réponds en mobilisant les notions de CDN, RCMD et co-bénéfices développement/climat.
\end{enumerate}
\end{exercice}

\begin{corrige}[Exercice Entraînement]
\begin{enumerate}
\item \textbf{B} ($0{,}3$ soit $30\%$).
\item \textbf{B} (\ce{CH4}, fermentation anaérobie dans sols inondés).
\item Émission journalière = $2 \times 1{,}8 = 3{,}6\ \text{kg}\ \ce{CO2}$. Annuelle = $3{,}6 \times 365 = 1314\ \text{kg} = \mathbf{1{,}314\ t\ \ce{CO2}}$. \textbf{Atténuation :} Foyer amélioré (réduit conso bois $\rightarrow$ $- \ce{CO2}$, $-$ déforestation). \textbf{Adaptation :} Diversification alimentaire (cultures moins dépendantes bois pour cuisson ?) ou Aménagement cuisine (ventilation) pour réduire pollution air intérieur (santé). \textit{Note : L'adaptation ici porte sur la vulnérabilité sanitaire/énergétique du ménage.}
\item \textbf{Chaîne :} Émissions GES mondiales $\rightarrow$ $\uparrow T°$ $\rightarrow$ $\uparrow$ Évaporation + $\downarrow$ Pluies Sahel $\rightarrow$ Régression Lac Tchad ($-90\%$) $\rightarrow$ Perte ressources pêche/élevage/agriculture (recession) $\rightarrow$ Perte revenus/sécurité alimentaire $\rightarrow$ Migrations populations riveraines vers zones humides Sud $\rightarrow$ Pression foncière/eau sur zones d'accueil $\rightarrow$ Conflits agriculteurs/éleveurs/autochtones/allochtones.
\item \textbf{Argumentaire structuré :}
\begin{enumerate}
\item \textbf{Équité (RCMD) :} Vrai, émissions faibles ($0{,}35\ \text{t}$/hab), mais responsabilité \emph{commune} : tout pays doit agir selon ses capacités (Art. 4 Accord Paris).
\item \textbf{Intérêt national (Co-bénéfices) :} Réduire émissions = Développement durable. Ex: Hydroélectricité/Solaire = Énergie accès + Moins fioul importé + Moins déforestation bois. Agroforesterie = Rendements + Sols + Carbone. Foyers améliorés = Santé (fumée) + Économies + Forêt.
\item \textbf{Accès financement :} CDN conditionnelle ($-35\%$) \textbf{exige} appui international (Fonds Vert, technologie). Ne pas s'engager = Perdre financements climat.
\item \textbf{Vulnérabilité :} Cameroun paye déjà le prix (Lac Tchad, Côte, Paludisme). Atténuation globale (dont sa part) limite réchauffement futur $\rightarrow$ Protège son propre territoire.
\end{enumerate}
\end{enumerate}
\end{corrige}

\newpage

\coursec{Exercices}

\courssub{Je vérifie mes connaissances}

\begin{exercice}[QCM : Gaz à effet de serre et terminologie]
Parmi les affirmations suivantes, une seule est correcte. La recopier en justifiant brièvement le choix.
\begin{enumerate}
    \item Le principal gaz à effet de serre d'origine anthropique est le méthane (\ce{CH4}).
    \item L'effet de serre naturel est un phénomène nuisible qu'il faut éliminer totalement.
    \item L'atténuation consiste à réduire les sources d'émission de GES, tandis que l'adaptation vise à limiter la vulnérabilité aux impacts du changement climatique.
    \item Le trou de la couche d'ozone est la cause principale du réchauffement climatique actuel.
\end{enumerate}
\end{exercice}

\begin{exercice}[Vrai ou Faux justifié]
Dire si chacune des affirmations suivantes est vraie ou fausse, en justifiant par une phrase précise.
\begin{enumerate}
    \item La vapeur d'eau (\ce{H2O}) est le gaz à effet de serre qui contribue le plus à l'effet de serre naturel.
    \item Le pouvoir de réchauffement global (PRG) du méthane sur 100 ans est inférieur à celui du dioxyde de carbone.
    \item La déforestation dans le bassin du Congo contribue au renforcement de l'effet de serre en réduisant le puits de carbone forestier.
    \item L'augmentation de la température moyenne globale est uniforme sur toute la surface du globe (mêmes valeurs aux pôles et à l'équateur).
\end{enumerate}
\end{exercice}

\begin{exercice}[Définitions et notions clés]
Donner une définition rigoureuse pour chacun des termes suivants :
\begin{enumerate}
    \item Forçage radiatif.
    \item Empreinte carbone.
    \item Effet de serre anthropique (ou renforcé).
    \item Résilience climatique (dans le contexte de l'adaptation).
\end{enumerate}
\end{exercice}

\begin{exercice}[Calcul direct : Émissions d'un véhicule]
Un taxi-brousse circulant à Yaoundé consomme en moyenne \SI{12}{L} de carburant (essence) par jour de travail. On admet que la combustion d'un litre d'essence émet \SI{2,3}{kg} de \ce{CO2}.
\begin{enumerate}
    \item Calculer la masse de \ce{CO2} émise par ce véhicule en une journée.
    \item Calculer la masse de \ce{CO2} émise en une année (300 jours de travail).
    \item Exprimer ce résultat en tonnes équivalent \ce{CO2} (\SI{1}{t} = \SI{1000}{kg}).
\end{enumerate}
\end{exercice}

\courssub{Je m'entraîne}

\begin{exercice}[Analyse du bilan radiatif terrestre]
Le schéma ci-dessous représente le bilan radiatif simplifié de la Terre (valeurs moyennes exprimées en unités, le rayonnement solaire incident au sommet de l'atmosphère valant 100 unités). Compléter les légendes des flèches (A, B, C, D, E) en choisissant parmi : \emph{Rayonnement solaire incident}, \emph{Rayonnement solaire réfléchi (albédo)}, \emph{Rayonnement solaire absorbé par la surface}, \emph{Rayonnement infrarouge émis par la surface}, \emph{Rayonnement infrarouge réémis vers la surface par l'atmosphère (effet de serre)}.

\begin{popfigure}
\begin{tikzpicture}[scale=0.95, every node/.style={font=\small}]
    % Soleil
    \node[draw, circle, fill=poporange!80, minimum size=1.1cm] (sun) at (-4.5,4.5) {Soleil};
    % Atmosphère (couche)
    \draw[popblue, thick] (-4,3.2) -- (5,3.2);
    \draw[popblue, thick] (-4,2.6) -- (5,2.6);
    \node[popblue] at (4.2,2.9) {Atmosphère};
    % Surface
    \draw[popgreen, very thick] (-4,0.5) -- (5,0.5);
    \node[popgreen] at (4.2,0.1) {Surface terrestre};
    % A : rayonnement solaire incident (100)
    \draw[->, very thick, poporange] (sun.east) -- (-1.5,3.6) node[midway, above] {\textbf{A} (100)};
    % B : réfléchi vers l'espace (30)
    \draw[->, thick, poporange] (-1.5,3.4) -- (-3.5,4.6) node[midway, above left] {\textbf{B} (30)};
    % C : absorbé par la surface (50)
    \draw[->, very thick, poporange] (-0.5,3.4) -- (-0.5,0.6) node[midway, right] {\textbf{C} (50)};
    % D : IR émis par la surface (vers le haut)
    \draw[->, very thick, poppink] (1.2,0.6) -- (1.2,3.4) node[midway, right] {\textbf{D} (117)};
    % E : IR réémis vers la surface par l'atmosphère
    \draw[->, very thick, poppink] (2.6,2.9) -- (2.6,0.6) node[midway, right] {\textbf{E} (96)};
    % IR sortant vers l'espace
    \draw[->, thick, poppink] (3.8,3.4) -- (4.8,4.6) node[midway, above right] {Vers l'espace};
\end{tikzpicture}
\legende{Bilan radiatif terrestre simplifié (valeurs moyennes, rayonnement solaire incident = 100 unités). Flèches oranges : rayonnement solaire (courtes longueurs d'onde) ; flèches roses : rayonnement infrarouge thermique (grandes longueurs d'onde).}
\end{popfigure}

\begin{enumerate}
    \item Légender les flèches A, B, C, D, E.
    \item Expliquer pourquoi la valeur de la flèche D (117 unités) est supérieure à 100 (rayonnement solaire incident).
    \item Identifier sur le schéma le phénomène d'\emph{effet de serre naturel}.
\end{enumerate}
\end{exercice}

\begin{exercice}[Interprétation de données : Concentration en \ce{CO2} et température]
Le tableau ci-dessous donne l'évolution de la concentration atmosphérique en \ce{CO2} (en ppm) et de l'anomalie de température moyenne globale (en $^\circ$C par rapport à la moyenne 1951-1980) depuis 1900.

\begin{center}
\begin{tabularx}{\linewidth}{|c|X|X|}\hline
\rowcolor{popblueL}\textbf{Année} & \textbf{Concentration \ce{CO2} (ppm)} & \textbf{Anomalie T ($^\circ$C)} \\ \hline
1900 & 295 & $-0{,}1$ \\ \hline
1950 & 310 & $-0{,}05$ \\ \hline
1980 & 339 & $+0{,}15$ \\ \hline
2000 & 370 & $+0{,}40$ \\ \hline
2020 & 414 & $+1{,}00$ \\ \hline
\end{tabularx}
\end{center}

\begin{enumerate}
    \item Représenter graphiquement (sur papier millimétré ou à l'aide d'un tableur) les deux courbes d'évolution sur un même graphique (deux axes des ordonnées, un axe des abscisses commun : le temps). Décrire la construction attendue.
    \item Décrire l'allure des deux courbes sur la période 1900-2020.
    \item Mettre en évidence la relation entre les deux grandeurs.
    \item En déduire l'influence des activités humaines (industrialisation, énergies fossiles) sur cette évolution.
\end{enumerate}
\end{exercice}

\begin{exercice}[Étude de cas : Émissions d'un village de l'Extrême-Nord]
Dans un village de l'Extrême-Nord du Cameroun, on observe cinq activités principales sources de gaz à effet de serre (GES). Compléter le tableau suivant en associant à chaque activité le principal GES émis (\ce{CO2}, \ce{CH4}, \ce{N2O}) et en proposant une mesure concrète de réduction adaptée au contexte local.

\begin{center}
\begin{tabularx}{\linewidth}{|l|c|X|}\hline
\rowcolor{popblueL}\textbf{Activité} & \textbf{Principal GES émis} & \textbf{Mesure de réduction (contexte local)} \\ \hline
Brûlis agricoles (nettoyage des champs) & & \\ \hline
Moto-taxis (moteurs essence, entretien aléatoire) & & \\ \hline
Élevage extensif de zébus (nombreux troupeaux) & & \\ \hline
Décharge à ciel ouvert (décomposition des déchets organiques) & & \\ \hline
Groupe électrogène diesel (électricité du village) & & \\ \hline
\end{tabularx}
\end{center}
\end{exercice}

\begin{exercice}[Classer : Atténuation ou Adaptation ?]
Classer les mesures suivantes dans un tableau à deux colonnes : « Atténuation » (réduction des causes) ou « Adaptation » (gestion des conséquences). Justifier brièvement chaque classement.
\begin{itemize}
    \item Reboisement de 10 000 ha sur les flancs du Mont Cameroun.
    \item Construction de digues de protection à Limbé contre la montée des eaux.
    \item Distribution de semences de maïs à cycle court et résistant à la sécheresse dans le Nord.
    \item Installation de panneaux solaires sur les toits des lycées de Douala.
    \item Mise en place d'un système d'alerte précoce aux inondations à Maroua.
    \item Promotion des foyers améliorés (réduction du bois de chauffe) dans les zones rurales.
\end{itemize}
\end{exercice}

\courssub{Je me prépare au Probatoire}

\begin{exercice}[Sujet type Probatoire : Mécanisme de l'effet de serre et bilan radiatif]
\textbf{Document 1 :} Schéma simplifié du bilan radiatif de la Terre (flux radiatifs numérotés 1 à 5, valeurs en \SI{}{W.m^{-2}}).
\textbf{Document 2 :} Spectre d'absorption de l'atmosphère terrestre : quasi-transparence dans le domaine visible ; forte absorption dans l'infrarouge thermique par \ce{H2O}, \ce{CO2} et \ce{CH4}.

\begin{enumerate}
    \item À l'aide du Document 1, légender les cinq flux radiatifs principaux en précisant leur nature (rayonnement solaire de courtes longueurs d'onde / rayonnement infrarouge de grandes longueurs d'onde) et leur origine/destination.
    \item En vous appuyant sur le Document 2, expliquer en 5 lignes maximum le mécanisme physique de l'effet de serre naturel. Utiliser les termes : \emph{rayonnement solaire}, \emph{absorption par la surface}, \emph{réémission infrarouge}, \emph{absorption par les GES}, \emph{réémission vers le sol}.
    \item Distinguer clairement « effet de serre naturel » et « effet de serre renforcé (anthropique) ». Citer deux activités humaines majeures responsables du renforcement.
    \item Calculer la température d'équilibre théorique de la Terre sans effet de serre (modèle du corps noir) sachant : constante solaire $S_0 = \SI{1361}{W.m^{-2}}$, albédo moyen $a = 0{,}3$, constante de Stefan-Boltzmann $\sigma = \SI{5,67e-8}{W.m^{-2}.K^{-4}}$. Formule : $T^4 = \dfrac{S_0(1-a)}{4\sigma}$. Comparer à la température moyenne observée ($+15\,^\circ$C) et conclure sur le rôle de l'atmosphère.
\end{enumerate}
\end{exercice}

\begin{exercice}[Sujet type Probatoire : Dynamique du carbone et responsabilité humaine]
\textbf{Document :} Courbes conjointes de la concentration atmosphérique en \ce{CO2} (ppm) et de l'anomalie de température globale ($^\circ$C) de 1850 à 2020 (courbes de type « crosse de hockey » : quasi stables avant 1900, puis croissance rapide). Données chiffrées : 1850 : 285 ppm / $-0{,}4\,^\circ$C ; 2020 : 415 ppm / $+1{,}1\,^\circ$C. Émissions anthropiques cumulées depuis 1850 : environ \SI{2400}{Gt} de \ce{CO2}.

\begin{enumerate}
    \item Décrire l'évolution conjointe des deux courbes sur la période 1850-2020. Préciser le point de rupture (inflexion) majeur.
    \item En utilisant la notion de corrélation, analyser le lien entre concentration en \ce{CO2} et température. Cette corrélation suffit-elle à prouver la causalité ? Justifier en citant le mécanisme physique.
    \item L'analyse isotopique du carbone (rapports \ce{^{13}C}/\ce{^{12}C} et \ce{^{14}C}) permet d'identifier l'origine du \ce{CO2} atmosphérique. Expliquer pourquoi la baisse du rapport \ce{^{13}C}/\ce{^{12}C} et la diminution du \ce{^{14}C} atmosphérique (effet Suess) confirment l'origine fossile (charbon, pétrole, gaz) de l'excès de \ce{CO2}.
    \item Le Cameroun, bien que faible émetteur (environ 0,1\,\% des émissions mondiales), s'est engagé dans sa Contribution Déterminée au Niveau National (CDN). Citer deux engagements d'atténuation et un engagement d'adaptation présents dans la CDN du Cameroun (actualisée en 2021).
\end{enumerate}
\end{exercice}

\begin{exercice}[Sujet type Probatoire : Synthèse argumentée — le Cameroun face au climat]
\textbf{Consigne :} Rédiger un texte argumenté structuré (introduction, développement en deux parties distinctes, conclusion) d'environ 15 à 20 lignes sur le thème : \emph{« Le Cameroun face au changement climatique : constats, vulnérabilités et réponses à multiples échelles »}.

Votre texte doit impérativement intégrer les éléments suivants :
\begin{itemize}
    \item \textbf{Deux conséquences concrètes observées ou projetées au Cameroun} (ex. : montée des eaux et érosion côtière à Kribi/Limbé, sécheresses récurrentes et désertification dans le Grand Nord, recrudescence du paludisme en altitude, baisse des rendements du cacao et du café).
    \item \textbf{Trois mesures de lutte à des échelles différentes} :
    \begin{enumerate}
        \item une mesure à l'échelle \emph{individuelle/locale} (ex. : agroforesterie, foyer amélioré, tri des déchets, mobilité douce) ;
        \item une mesure à l'échelle \emph{nationale} (ex. : Plan National d'Adaptation, promotion des énergies renouvelables, gestion durable des forêts du bassin du Congo, reboisement « Opération Sahel Vert ») ;
        \item une mesure à l'échelle \emph{internationale} (ex. : Accord de Paris, mécanisme REDD+, Fonds Vert pour le Climat, transferts de technologies propres).
    \end{enumerate}
\end{itemize}
La rédaction sera évaluée sur la rigueur scientifique du vocabulaire (GES, forçage radiatif, résilience, atténuation, adaptation), la cohérence logique et l'ancrage dans le contexte camerounais.
\end{exercice}



\newpage

\coursec{Corrigés des exercices}

\begin{corrige}[Exercice 1 — QCM : Gaz à effet de serre et terminologie]
La bonne réponse est la \textbf{3}.
\begin{itemize}
    \item \textbf{1. Fausse} : le principal GES d'origine anthropique, tant en quantité émise qu'en contribution au forçage radiatif, est le dioxyde de carbone \ce{CO2} (environ les trois quarts des émissions anthropiques). Le méthane \ce{CH4} arrive en deuxième position.
    \item \textbf{2. Fausse} : l'effet de serre naturel est \emph{indispensable} à la vie : sans lui, la température moyenne du globe serait d'environ $-18\,^\circ$C au lieu de $+15\,^\circ$C. C'est son \emph{renforcement} par les activités humaines qui pose problème, pas le phénomène lui-même.
    \item \textbf{3. Vraie} : c'est la définition exacte des deux stratégies complémentaires de lutte contre le changement climatique : l'\cle{atténuation} agit sur les causes (réduction des émissions, puits de carbone), l'\cle{adaptation} agit sur les conséquences (réduction de la vulnérabilité).
    \item \textbf{4. Fausse} : le trou de la couche d'ozone (dû aux CFC, qui laissent passer les UV) et le réchauffement climatique (dû aux GES, qui piègent l'infrarouge) sont deux problèmes \emph{distincts}, de causes et de mécanismes différents. Confusion classique à éviter absolument au Probatoire.
\end{itemize}
\end{corrige}

\begin{corrige}[Exercice 2 — Vrai ou Faux justifié]
\begin{enumerate}
    \item \textbf{Vrai} : la vapeur d'eau est le principal contributeur de l'effet de serre \emph{naturel} (environ 60\,\% de l'effet total), devant le \ce{CO2}. Mais sa concentration dépend de la température de l'air : elle agit comme une \emph{amplification} (rétroaction) et non comme la cause initiale du réchauffement actuel.
    \item \textbf{Faux} : le PRG du méthane sur 100 ans vaut environ 28 à 30 fois celui du \ce{CO2} (pris comme référence, PRG = 1). Le méthane est donc un gaz beaucoup plus \emph{puissant}, mais moins abondant et moins persistant (durée de vie $\approx$ 12 ans contre des siècles pour le \ce{CO2}).
    \item \textbf{Vrai} : la forêt du bassin du Congo (deuxième massif forestier tropical du monde, largement présent au Cameroun) stocke du carbone dans sa biomasse et ses sols. Sa destruction libère du \ce{CO2} (brûlis, décomposition) \emph{et} supprime un puits qui absorbait le \ce{CO2} atmosphérique par photosynthèse : double effet aggravant.
    \item \textbf{Faux} : le réchauffement est \emph{hétérogène} : il est environ deux à trois fois plus rapide dans les régions polaires (amplification arctique) qu'à l'équateur, et plus marqué sur les continents que sur les océans, dont l'inertie thermique est grande.
\end{enumerate}
\end{corrige}

\begin{corrige}[Exercice 3 — Définitions et notions clés]
\begin{enumerate}
    \item \textbf{Forçage radiatif} : variation du bilan radiatif de la Terre (différence entre l'énergie reçue et l'énergie renvoyée vers l'espace), exprimée en $ \text{W}\cdot\text{m}^{-2} $, provoquée par un facteur externe (augmentation des GES, variation solaire, aérosols). Un forçage positif entraîne un réchauffement.
    \item \textbf{Empreinte carbone} : quantité totale de gaz à effet de serre émise directement ou indirectement par une activité, un produit, un individu ou un pays, exprimée en tonnes équivalent \ce{CO2} (t\ce{CO2}e).
    \item \textbf{Effet de serre anthropique (renforcé)} : accroissement de l'effet de serre naturel dû à l'augmentation des concentrations de GES (\ce{CO2}, \ce{CH4}, \ce{N2O}...) liée aux activités humaines, qui accroît le forçage radiatif et provoque le réchauffement climatique.
    \item \textbf{Résilience climatique} : capacité d'un système (population, écosystème, infrastructure) à résister aux chocs climatiques, à s'y adapter et à récupérer après une perturbation (sécheresse, inondation), tout en maintenant ses fonctions essentielles.
\end{enumerate}
\end{corrige}

\begin{corrige}[Exercice 4 — Calcul direct : Émissions d'un véhicule]
\begin{enumerate}
    \item Masse journalière (consommation 12 L/j, facteur d'émission 2,3 kg \ce{CO2}/L) :
    \[ m_j = 12\,\text{L/j} \times 2{,}3\,\text{kg/L} = \SI{27,6}{kg} \text{ de } \ce{CO2} \text{ par jour.} \]
    \item Masse annuelle (300 jours d'activité) :
    \[ m_a = 27{,}6 \times 300 = \SI{8280}{kg} \text{ de } \ce{CO2} \text{ par an.} \]
    \item Conversion en tonnes :
    \[ m_a = \frac{8280}{1000} = \SI{8,28}{t} \text{ équivalent } \ce{CO2} \text{ par an.} \]
\end{enumerate}
\emph{Commentaire} : plus de 8 tonnes de \ce{CO2} par an pour un seul véhicule : ce résultat illustre l'intérêt des politiques de transport collectif, de l'entretien des moteurs et du renouvellement du parc automobile dans les villes camerounaises (Yaoundé, Douala).
\end{corrige}

\begin{corrige}[Exercice 5 — Analyse du bilan radiatif terrestre]
\begin{enumerate}
    \item \textbf{A} : rayonnement solaire incident ; \textbf{B} : rayonnement solaire réfléchi vers l'espace (albédo, environ 30\,\%) ; \textbf{C} : rayonnement solaire absorbé par la surface (environ 50 unités, le reste étant absorbé par l'atmosphère) ; \textbf{D} : rayonnement infrarouge émis par la surface chauffée ; \textbf{E} : rayonnement infrarouge réémis vers la surface par l'atmosphère (contre-rayonnement).
    \item La surface reçoit de l'énergie de \emph{plusieurs} sources : le rayonnement solaire absorbé (C, 50 unités) \emph{et} le rayonnement infrarouge renvoyé par l'atmosphère (E, 96 unités), auxquels s'ajoutent les flux de convection et d'évaporation. Pour rester en équilibre thermique, elle doit évacuer autant d'énergie qu'elle en reçoit : son émission infrarouge D (117 unités) dépasse donc les 100 unités du flux solaire incident. Il n'y a \emph{aucune violation} du principe de conservation de l'énergie : à l'échelle du système global Terre-atmosphère, l'énergie renvoyée vers l'espace (B + infrarouge sortant) égale exactement l'énergie solaire reçue.
    \item L'effet de serre naturel correspond à la flèche \textbf{E} : les gaz à effet de serre (\ce{H2O}, \ce{CO2}, \ce{CH4}) absorbent une grande partie du rayonnement infrarouge émis par la surface (D) et le réémettent dans toutes les directions, notamment vers le sol, ce qui maintient la température moyenne à $+15\,^\circ$C au lieu de $-18\,^\circ$C.
\end{enumerate}
\textbf{Point de vigilance} : ne jamais écrire que « la Terre émet plus d'énergie qu'elle n'en reçoit » sans préciser le niveau (surface seule \emph{vs} système Terre-atmosphère) : c'est l'erreur la plus fréquente sur ce type de schéma.
\end{corrige}

\begin{corrige}[Exercice 6 — Interprétation de données : Concentration en \ce{CO2} et température]
\begin{enumerate}
    \item Construction attendue : axe des abscisses = années (1900 à 2020) ; axe des ordonnées de gauche = concentration en \ce{CO2} (gradué par exemple de 280 à 420 ppm) ; axe des ordonnées de droite = anomalie de température (de $-0{,}2$ à $+1{,}2\,^\circ$C). On place les cinq points de chaque série, on les relie, et on légende chaque courbe.
    \item Les deux courbes sont \emph{croissantes} sur toute la période, avec une croissance lente avant 1950 puis de plus en plus rapide (allure exponentielle) : la concentration passe de 295 à 414 ppm (+40\,\%) et l'anomalie de $-0{,}1$ à $+1{,}0\,^\circ$C.
    \item Les deux grandeurs évoluent de façon concomitante : il existe une \emph{corrélation positive forte} entre la concentration en \ce{CO2} et la température moyenne globale.
    \item L'accélération après 1950 coïncide avec l'explosion de la consommation mondiale de charbon, de pétrole et de gaz (industrie, transports, production d'électricité). Or le mécanisme physique est connu : le \ce{CO2} absorbe le rayonnement infrarouge émis par la Terre et le réémet vers le sol. La corrélation, appuyée par ce mécanisme et par les preuves isotopiques, permet d'attribuer le réchauffement récent aux émissions anthropiques de GES.
\end{enumerate}
\textbf{Point de vigilance} : exiger une phrase du type « corrélation + mécanisme physique $\Rightarrow$ causalité » ; une simple constatation de parallélisme des courbes ne suffit pas à conclure.
\end{corrige}

\begin{corrige}[Exercice 7 — Étude de cas : Émissions d'un village de l'Extrême-Nord]
\begin{center}
\begin{tabularx}{\linewidth}{|l|c|X|}\hline
\rowcolor{popblueL}\textbf{Activité} & \textbf{GES} & \textbf{Mesure de réduction} \\ \hline
Brûlis agricoles & \ce{CO2} & Remplacer le brûlis par l'enfouissement des résidus ou le compostage ; paillage (agriculture de conservation). \\ \hline
Moto-taxis & \ce{CO2} & Entretien régulier des moteurs, covoiturage, promotion des déplacements à vélo ou à pied sur courtes distances. \\ \hline
Élevage de zébus & \ce{CH4} & Améliorer l'alimentation du bétail (fourrages de qualité), récupérer le fumier pour le biogaz ou le compost. \\ \hline
Décharge à ciel ouvert & \ce{CH4} & Trier et composter les déchets organiques ; capter le biogaz ; enfouissement contrôlé. \\ \hline
Groupe électrogène diesel & \ce{CO2} & Installer un mini-réseau solaire photovoltaïque avec batteries (ensoleillement exceptionnel de l'Extrême-Nord). \\ \hline
\end{tabularx}
\end{center}
\emph{Rappel} : la fermentation entérique des ruminants et la décomposition anaérobie des déchets produisent du méthane \ce{CH4} ; la combustion des carburants fossiles et de la biomasse produit essentiellement du \ce{CO2}. Le \ce{N2O} provient surtout des engrais azotés et des sols cultivés.
\end{corrige}

\begin{corrige}[Exercice 8 — Classer : Atténuation ou Adaptation ?]
\begin{center}
\begin{tabularx}{\linewidth}{|X|X|}\hline
\rowcolor{popblueL}\textbf{Atténuation (agir sur les causes)} & \textbf{Adaptation (agir sur les conséquences)} \\ \hline
Reboisement du Mont Cameroun : la forêt absorbe le \ce{CO2} par photosynthèse (puits de carbone). & Digues à Limbé : protection contre la montée du niveau marin, conséquence du réchauffement. \\ \hline
Panneaux solaires à Douala : énergie renouvelable qui remplace des énergies fossiles émettrices de \ce{CO2}. & Semences résistantes à la sécheresse : maintien des récoltes malgré l'aridification du Nord. \\ \hline
Foyers améliorés : moins de bois brûlé, donc moins de \ce{CO2} émis et moins de déforestation. & Alerte précoce aux inondations à Maroua : réduction de la vulnérabilité des populations face aux crues. \\ \hline
\end{tabularx}
\end{center}
\textbf{Critère de classement à retenir} : la mesure réduit-elle les émissions ou augmente-elle les puits de carbone ? $\rightarrow$ atténuation. La mesure réduit-elle les dégâts d'un impact climatique ? $\rightarrow$ adaptation.
\end{corrige}

\begin{corrige}[Exercice 9 — Sujet type Probatoire : Mécanisme de l'effet de serre et bilan radiatif]
\begin{enumerate}
    \item Flux 1 : rayonnement solaire incident (courtes longueurs d'onde), du Soleil vers la Terre. Flux 2 : rayonnement solaire réfléchi vers l'espace par les nuages, l'atmosphère et les surfaces claires (albédo). Flux 3 : rayonnement solaire absorbé par la surface terrestre. Flux 4 : rayonnement infrarouge (grandes longueurs d'onde) émis par la surface chauffée vers l'atmosphère. Flux 5 : rayonnement infrarouge réémis par l'atmosphère (GES) vers la surface — c'est l'effet de serre.
    \item Le rayonnement solaire, essentiellement visible, traverse l'atmosphère quasi transparente et subit une absorption par la surface, qui s'échauffe. La surface chauffée effectue une réémission infrarouge vers l'espace. Mais l'absorption par les GES (\ce{H2O}, \ce{CO2}, \ce{CH4}) piège une grande partie de cet infrarouge ; les GES procèdent alors à une réémission vers le sol, qui maintient la surface plus chaude qu'elle ne le serait sans atmosphère.
    \item L'\emph{effet de serre naturel} est le piégeage infrarouge par les GES présents naturellement ; il est indispensable à la vie ($+15\,^\circ$C en moyenne). L'\emph{effet de serre renforcé} est son amplification par l'augmentation des concentrations de GES due aux activités humaines : combustion des énergies fossiles (charbon, pétrole, gaz) et déforestation (qui réduit les puits de carbone).
    \item Calcul de la température d'équilibre sans effet de serre :
    \begin{align*}
    T^4 &= \frac{S_0(1-a)}{4\sigma} = \frac{1361 \times (1-0{,}3)}{4 \times 5{,}67 \times 10^{-8}} = \frac{952{,}7}{2{,}268 \times 10^{-7}} \approx 4{,}20 \times 10^{9}\ \text{K}^4 \\
    T &\approx \sqrt[4]{4{,}20 \times 10^{9}} \approx \SI{255}{K} \approx -18\,^\circ\text{C}
    \end{align*}
    La température d'équilibre sans effet de serre ($-18\,^\circ$C) est inférieure d'environ $33\,^\circ$C à la température moyenne observée ($+15\,^\circ$C). Conclusion : l'atmosphère, grâce à ses gaz à effet de serre, réchauffe naturellement la surface d'environ $33\,^\circ$C et rend la planète habitable.
\end{enumerate}
\textbf{Barème indicatif (20 pts)} : légende des flux (5 pts) ; mécanisme en 5 lignes avec les 5 termes imposés (5 pts) ; distinction naturel/renforcé + 2 activités (4 pts) ; calcul littéral et numérique correct avec conversion K/$^\circ$C (4 pts) ; conclusion (2 pts).

\textbf{Points de vigilance} : citer la formule avant l'application numérique ; ne pas oublier le facteur 4 (la Terre intercepte le soleil sur un disque mais rayonne sur toute sa sphère) ; convertir les kelvins en degrés Celsius ($T(^\circ\text{C}) = T(\text{K}) - 273$) ; rédiger la question 2 avec exactement les cinq termes imposés, soulignés si possible.
\end{corrige}

\begin{corrige}[Exercice 10 — Sujet type Probatoire : Dynamique du carbone et responsabilité humaine]
\begin{enumerate}
    \item De 1850 à environ 1900-1950, les deux courbes sont quasi stables (285 ppm ; anomalies voisines de $-0{,}4\,^\circ$C). Après 1950, elles s'élèvent fortement et simultanément jusqu'à 415 ppm et $+1{,}1\,^\circ$C en 2020. Le point de rupture majeur se situe au milieu du XX\textsuperscript{e} siècle, et coïncide avec l'industrialisation massive et l'explosion de la consommation d'énergies fossiles.
    \item Il existe une forte corrélation positive entre concentration en \ce{CO2} et température. Mais une corrélation seule ne prouve \emph{pas} une causalité. Ici, la causalité repose sur un mécanisme physique établi : le \ce{CO2} absorbe le rayonnement infrarouge émis par la surface et le réémet vers le sol (effet de serre) ; toute augmentation de sa concentration accroît le forçage radiatif, donc la température. Corrélation + mécanisme + preuves isotopiques = attribution solide.
    \item Les végétaux, et donc les combustibles fossiles qui en dérivent, sont appauvris en \ce{^{13}C} (la photosynthèse discrimine les isotopes lourds) : l'apport massif de \ce{CO2} fossile fait donc \emph{baisser} le rapport \ce{^{13}C}/\ce{^{12}C} atmosphérique. De plus, le \ce{^{14}C} est radioactif (période $\approx$ 5 730 ans) : il a totalement disparu des combustibles fossiles vieux de millions d'années. Leur combustion dilue donc le \ce{^{14}C} atmosphérique (effet Suess). Ces deux signatures isotopiques démontrent que l'excès de \ce{CO2} provient bien des énergies fossiles, et non des volcans ou du dégazage océanique.
    \item Atténuation (CDN 2021 du Cameroun) : réduction des émissions de l'ordre de 35\,\% à l'horizon 2030 par rapport au scénario tendanciel (conditionnée à l'appui international) ; développement des énergies renouvelables (hydroélectricité, solaire) et gestion durable des forêts / reboisement (lutte contre la déforestation du bassin du Congo). Adaptation : renforcement de la résilience des secteurs agricole, de la santé et des zones côtières (Plans Nationaux d'Adaptation, agriculture résiliente au climat).
\end{enumerate}
\textbf{Barème indicatif (20 pts)} : description des courbes et rupture (4 pts) ; analyse corrélation/causalité avec mécanisme (5 pts) ; explication des deux signatures isotopiques (6 pts) ; engagements de la CDN (5 pts).

\textbf{Points de vigilance} : distinguer clairement \emph{corrélation} (constat statistique) et \emph{causalité} (mécanisme) ; pour l'effet Suess, bien préciser que le \ce{^{14}C} a disparu des fossiles \emph{par décroissance radioactive} ; citer des engagements chiffrés ou nommés (35\,\%, REDD+, énergies renouvelables) plutôt que des formules vagues.
\end{corrige}

\begin{corrige}[Exercice 11 — Sujet type Probatoire : Synthèse argumentée — le Cameroun face au climat]
\textbf{Introduction attendue} : définition du problème (renforcement de l'effet de serre par les émissions anthropiques de GES, forçage radiatif positif) et annonce d'un plan en deux parties (constats/vulnérabilités, puis réponses à plusieurs échelles).

\textbf{Première partie — constats et vulnérabilités} : le candidat cite au moins deux impacts contextualisés, par exemple :
\begin{itemize}
    \item l'érosion côtière et la montée du niveau marin menaçant Limbé et Kribi (littoral atlantique) ;
    \item l'avancée de la désertification et les sécheresses récurrentes dans l'Extrême-Nord (recul du lac Tchad), qui fragilisent l'agriculture et l'élevage ;
    \item l'extension du paludisme vers les hautes terres de l'Ouest, autrefois épargnées, du fait de températures plus clémentes pour le moustique vecteur ;
    \item la baisse projetée des rendements du cacao et du café, piliers de l'économie rurale.
\end{itemize}

\textbf{Deuxième partie — les réponses à trois échelles} :
\begin{itemize}
    \item \emph{Locale/individuelle} : adoption des foyers améliorés (moins de bois de chauffe, donc moins de \ce{CO2} et moins de déforestation), agroforesterie cacaoyère, compostage des déchets organiques (réduction du \ce{CH4}), mobilité douce.
    \item \emph{Nationale} : mise en œuvre du Plan National d'Adaptation aux Changements Climatiques, développement hydroélectrique et solaire, gestion durable des forêts camerounaises du bassin du Congo (puits de carbone majeur), campagnes de reboisement dans le Grand Nord (« Opération Sahel Vert »).
    \item \emph{Internationale} : engagement du Cameroun dans l'Accord de Paris (CDN actualisée 2021), participation au mécanisme REDD+ (rémunération de la réduction des émissions liées à la déforestation), accès au Fonds Vert pour le Climat pour financer l'adaptation, transferts de technologies propres.
\end{itemize}

\textbf{Conclusion attendue} : distinction claire entre atténuation (réduire les émissions) et adaptation (accroître la résilience) ; le Cameroun, faible émetteur (environ 0,1\,\% des émissions mondiales) mais très vulnérable, doit combiner les deux stratégies avec l'appui de la coopération internationale.

\textbf{Barème indicatif (20 pts)} : rigueur du vocabulaire scientifique — GES, forçage radiatif, résilience, atténuation, adaptation (4 pts) ; exactitude des faits et des exemples (4 pts) ; structure logique en trois parties (4 pts) ; ancrage dans le contexte camerounais (4 pts) ; qualité de la langue (4 pts).

\textbf{Points de vigilance} : respecter impérativement la structure imposée (introduction, deux parties, conclusion) ; employer les termes scientifiques à bon escient (ne pas écrire « pollution de l'ozone » pour « effet de serre ») ; attribuer chaque mesure à la bonne échelle ; éviter les généralités non étayées (« il faut protéger la nature ») au profit de mesures concrètes et nommées.
\end{corrige}

\newpage

\coursec{Fiche bilan}

\begin{popfigure}
\begin{tikzpicture}[
  mindmap,
  grow cyclic,
  every node/.style={concept, execute at begin node=\hskip0pt},
  concept color=popblue,
  level 1/.style={level distance=4.5cm, sibling angle=360/6},
  level 2/.style={level distance=3cm, sibling angle=45},
  text=white,
  font=\small\bfseries
]
\node[concept, scale=1.2, align=center]{Lutte contre\\ l'effet de serre}
  child[concept color=popgreen] { node[concept, align=center]{1. Effet de serre\\ naturel}
    child { node[concept, scale=0.7] {Rayonnement solaire} }
    child { node[concept, scale=0.7] {Absorption IR par GES} }
    child { node[concept, scale=0.7] {Réémission vers le sol} }
    child { node[concept, scale=0.7] {Température +15 °C} }
  }
  child[concept color=poporange] { node[concept, align=center]{2. Renforcement\\ anthropique}
    child { node[concept, scale=0.7] {Combustion fossiles} }
    child { node[concept, scale=0.7] {Déforestation} }
    child { node[concept, scale=0.7] {Élevage (CH₄)} }
    child { node[concept, scale=0.7] {Industrie (N₂O, CFC)} }
  }
  child[concept color=poppurple] { node[concept, align=center]{3. Conséquences\\ climatiques}
    child { node[concept, scale=0.7] {Hausse T° moyenne} }
    child { node[concept, scale=0.7] {Fonte glaces / Mer} }
    child { node[concept, scale=0.7] {Événements extrêmes} }
    child { node[concept, scale=0.7] {Biodiversité menacée} }
  }
  child[concept color=poppink] { node[concept, align=center]{4. Atténuation\\ (Réduction)}
    child { node[concept, scale=0.7] {Énergies renouvelables} }
    child { node[concept, scale=0.7] {Efficacité énergétique} }
    child { node[concept, scale=0.7] {Reboisement / Puits} }
    child { node[concept, scale=0.7] {Mobilité douce} }
  }
  child[concept color=popgold] { node[concept, align=center]{5. Adaptation\\ (Résilience)}
    child { node[concept, scale=0.7] {Agriculture résiliente} }
    child { node[concept, scale=0.7] {Gestion eau / Risques} }
    child { node[concept, scale=0.7] {Urbanisme durable} }
    child { node[concept, scale=0.7] {Systèmes d'alerte} }
  }
  child[concept color=popdark] { node[concept, align=center]{6. Politiques\\ \& Cameroun}
    child { node[concept, scale=0.7] {Accord de Paris} }
    child { node[concept, scale=0.7] {CDN Cameroun} }
    child { node[concept, scale=0.7] {Plan Climat} }
    child { node[concept, scale=0.7] {Éducation / Citoyenneté} }
  };
\end{tikzpicture}
\legende{Carte mentale récapitulative de la leçon ND11 : mécanisme naturel, causes humaines, impacts, solutions d'atténuation et d'adaptation, cadres politiques (international et national).}
\end{popfigure}

\begin{bilan}
\courssub{Définitions clés}
\begin{itemize}
  \item \textbf{Effet de serre naturel} : Phénomène par lequel certains gaz de l'atmosphère (GES) absorbent une partie du rayonnement infrarouge émis par la Terre et le réémettent vers le sol, maintenant une température moyenne de \SI{15}{\celsius} (au lieu de \SI{-18}{\celsius} sans atmosphère).
  \item \textbf{Gaz à effet de serre (GES)} : Gaz absorbant le rayonnement infrarouge. Principaux : \ce{H2O} (vapeur d'eau), \ce{CO2} (dioxyde de carbone), \ce{CH4} (méthane), \ce{N2O} (protoxyde d'azote), CFC/HFC.
  \item \textbf{Forçage radiatif} : Différence entre l'énergie entrante et sortante au sommet de l'atmosphère (W·m⁻²). Positif = réchauffement.
  \item \textbf{Atténuation} : Actions réduisant les émissions de GES ou augmentant les puits de carbone (ex. : transition énergétique, reboisement).
  \item \textbf{Adaptation} : Ajustements face aux effets actuels ou futurs du changement climatique (ex. : variétés résistantes à la sécheresse, digues).
  \item \textbf{Puits de carbone} : Réservoir absorbant plus de carbone qu'il n'en émet (forêts, océans, sols).
\end{itemize}

\courssub{Mécanismes à maîtriser (schéma explicatif)}
\begin{enumerate}
  \item \textbf{Bilan radiatif} : \SI{340}{W.m^{-2}} entrants $\rightarrow$ \SI{100}{W.m^{-2}} réfléchis (albédo) $\rightarrow$ \SI{240}{W.m^{-2}} absorbés par le système Terre-atmosphère.
  \item \textbf{Rayonnement terrestre} : Sol émet \SI{390}{W.m^{-2}} (IR). Atmosphère (GES) en absorbe \SI{350}{W.m^{-2}} et réémet \SI{324}{W.m^{-2}} vers le sol (contre-rayonnement).
  \item \textbf{Renforcement} : Hausse [GES] $\rightarrow$ absorption IR accrue $\rightarrow$ contre-rayonnement $\uparrow$ $\rightarrow$ T° sol $\uparrow$ $\rightarrow$ nouvel équilibre plus chaud.
\end{enumerate}

\courssub{Tableau récapitulatif : GES majeurs et origines anthropiques}
\begin{center}
\begin{tabularx}{\linewidth}{|l|X|X|X|}\hline
\rowcolor{popblueL}
\textbf{Gaz} & \textbf{Pouvoir de réchauffement global (PRG\footnote{Sur 100 ans, ref \ce{CO2}=1.})} & \textbf{Principales sources anthropiques (Cameroun \& Monde)} & \textbf{Durée de vie atm.} \\ \hline
\ce{CO2} & 1 & Combustion fossiles (transports, industrie, centrales), déforestation, cimenteries (Figuil) & \SIrange{100}{1000}{ans} \\ \hline
\ce{CH4} & 28--30 & Élevage (fermentation entérique, déjections), riziculture, fuites gaz/charbon, décharges (Nkolfoulou, Yaoundé) & \SI{12}{ans} \\ \hline
\ce{N2O} & 265 & Engrais azotés (agriculture intensive), combustion biomasse, industrie chimique & \SI{121}{ans} \\ \hline
CFC / HFC & \numrange{1000}{10000} & Réfrigération, climatisation, aérosols (interdits/protocole Montréal/Kigali) & \SIrange{10}{1000}{ans} \\ \hline
\end{tabularx}
\end{center}

\courssub{Méthodes express (Savoir-faire probatoire)}
\begin{methode}[Expliquer l'effet de serre à partir d'un schéma]
  \begin{enumerate}
    \item Identifier les flux : rayonnement solaire entrant (courtes longueurs d'onde), rayonnement infrarouge sortant (longues longueurs d'onde).
    \item Montrer l'absorption par les GES (liaisons moléculaires vibrantes) et la réémission omnidirectionnelle.
    \item Conclure : le piégeage de la chaleur maintient T° compatible avec la vie ; l'augmentation des GES épaissit la « couverture ».
  \end{enumerate}
\end{methode}

\begin{methode}[Relier activité humaine $\rightarrow$ émission GES]
  \begin{enumerate}
    \item Identifier le secteur : Transport, Énergie, Agriculture/Forêt, Industrie, Déchets, Résidentiel.
    \item Préciser le processus : Combustion ($\ce{C} + \ce{O2} \rightarrow \ce{CO2}$), Fermentation anaérobie ($\ce{C6H12O6} \rightarrow \ce{CH4} + \ce{CO2}$), Nitrification/Dénitrification ($\ce{NH4+} \rightarrow \ce{N2O}$).
    \item Quantifier si possible (ex. : 1 L essence $\approx$ \SI{2,3}{kg} \ce{CO2}).
  \end{enumerate}
\end{methode}

\begin{methode}[Proposer actions atténuation/adaptation (Contexte Cameroun)]
  \begin{enumerate}
    \item \textbf{Atténuation} : Hydroélectricité (barrages Nachtigal, Lom Pangar, Memve'ele) $\rightarrow$ moins thermique ; Agroforesterie/cacao-ombragé $\rightarrow$ puits carbone ; Transport en commun (bus BRT Douala/Yaoundé) ; Cuiseurs améliorés (réduit bois/charbon).
    \item \textbf{Adaptation} : Variétés courtes cycles/résistantes (maïs, mil, sorgho) ; Agroécologie (zaï, haies vives) ; Protection mangroves (Wouri, Rio del Rey) contre érosion côtière ; Systèmes d'alerte crues (Bénoué, Logone).
  \end{enumerate}
\end{methode}
\end{bilan}

\begin{aretenir}[Checklist avant le Probatoire]
  $\square$ Je sais définir l'effet de serre naturel et donner la température moyenne terrestre grâce à lui (\SI{15}{\celsius}).\\
  $\square$ Je cite au moins 4 GES et leurs sources anthropiques principales (dont contexte Cameroun).\\
  $\square$ J'explique le mécanisme sur un schéma : absorption IR par liaisons moléculaires $\rightarrow$ réémission vers le sol.\\
  $\square$ Je distingue \textbf{forçage radiatif} (cause) et \textbf{réchauffement} (conséquence).\\
  $\square$ Je calcule/estime un bilan carbone simple (ex. trajet, consommation électricité).\\
  $\square$ Je différencie clairement \textbf{atténuation} (agir sur les causes) et \textbf{adaptation} (agir sur les effets).\\
  $\square$ Je donne 3 exemples d'atténuation et 3 d'adaptation pertinents pour le Cameroun.\\
  $\square$ Je connais les grands accords : UNFCCC (Rio 1992), Kyoto (1997), Paris (2015, objectif \SI{1,5}{\textendash}\SI{2}{\celsius}).\\
  $\square$ Je sais ce qu'est une CDN (Contribution Déterminée au niveau National) et l'objectif Cameroun (\SI{-35}{\%} émissions 2030 vs scénario tendanciel).\\
  $\square$ J'évite les confusions : trou dans la couche d'ozone $\neq$ effet de serre ; vapeur d'eau = GES principal mais rétroaction, pas forçage initial direct.\\
  $\square$ Je sais argumenter : « Pourquoi lutter ? » (coût inaction $>$ coût action ; justice climatique ; santé ; biodiversité).\\
  $\square$ Je maîtrise le vocabulaire précis : anthropique, résilience, séquestration, transition énergétique, développement durable.
\end{aretenir}

\findecours

\end{document}
